% Lesson 19 — Vocabulary: Words and expressions related to requesting for services — Anglais 1ère A — Cours signé OnBuch+
% Programme officiel MINESEC. Compiler avec : tectonic ENA19-vocabulary-words-and-expressions-related-to-requesting.tex
\newcommand{\DOCMATIERE}{Anglais}
\newcommand{\DOCNIVEAU}{1ère A}
\newcommand{\DOCMODULE}{Chapter 2 : Economic Life and Occupations}
\newcommand{\DOCLECON}{ENA19}
\newcommand{\DOCTITRE}{Lesson 19 — Vocabulary: Words and expressions related to requesting for services}
% OnBuch+ — préambule « cours signés » ENRICHI (compilé avec tectonic / XeLaTeX)
% Système visuel « Pop » d'OnBuch+ : fond crème (jamais blanc), bordures encre
% épaisses, ombres « dures » décalées, titres Archivo Black en capitales.
% Version MATHÉMATIQUES : graphes (pgfplots/tikz), boîtes pédagogiques
% (définition, propriété, méthode, attention, le savais-tu, exercice résolu,
% exercice, TP, bilan), page de garde et signature OnBuch+ ancrée en bas.
%
% Macros attendues dans main.tex AVANT \input{preamble} :
%   \DOCMATIERE  \DOCNIVEAU  \DOCMODULE  \DOCLECON (numéro)  \DOCTITRE
\documentclass[11pt]{article}

\usepackage{fontspec}
\usepackage[english,french]{babel} % français = langue principale ; l'anglais via \eng{..} / textebox
\usepackage[a4paper,margin=2cm,top=3.4cm,bottom=2.8cm,headheight=1.4cm,footskip=1.3cm]{geometry}
\usepackage{amsmath,amssymb,amsfonts}
\usepackage{mathtools} % \xmapsto, \coloneqq... souvent utilisés par les modèles
\usepackage{cancel} % simplifications barrées (bilans de forces)
% \abs{z} (module, valeur absolue) et \norm{u} (norme), utilisés par les modèles
\providecommand{\abs}{}\let\abs\relax\DeclarePairedDelimiter{\abs}{\lvert}{\rvert}
\providecommand{\norm}{}\let\norm\relax\DeclarePairedDelimiter{\norm}{\lVert}{\rVert}
\usepackage[table]{xcolor} % [table] : fournit aussi \rowcolors (alternance de lignes)
\usepackage{pagecolor}
\usepackage{tikz}
\usetikzlibrary{arrows.meta,positioning,calc,shapes.geometric,shapes.misc,shapes.arrows,decorations.pathmorphing,decorations.pathreplacing,decorations.markings,patterns,shadows,mindmap,trees,fit,backgrounds,matrix,intersections,angles,quotes}
\usepackage{pgfplots}
\usepackage[version=4]{mhchem} % équations nucléaires \ce{^{235}_{92}U}
\usepackage[european,siunitx]{circuitikz} % schémas électriques
\pgfplotsset{compat=1.18}
\usepgfplotslibrary{fillbetween,groupplots} % aires entre courbes (calcul intégral)
\usepackage{enumitem}
\usepackage{fancyhdr}
% [most] charge toutes les bibliothèques tcolorbox (listings, minted, raster,
% documentation...) et gonfle inutilement la mémoire TeX sur un document long
% (~300 boîtes) : on ne charge que ce qui est réellement utilisé.
\usepackage[skins,breakable]{tcolorbox}
\usepackage{graphicx}
\usepackage{colortbl}
\usepackage{booktabs}
\usepackage{tabularx}
\usepackage{diagbox} % case d'en-tête diagonale des tableaux à double entrée
% Colonnes extensibles centrée / gauche / droite (C, L, R), souvent utilisées par les modèles
\newcolumntype{C}{>{\centering\arraybackslash}X}
\newcolumntype{L}{>{\raggedright\arraybackslash}X}
\newcolumntype{R}{>{\raggedleft\arraybackslash}X}
\usepackage{array}
\usepackage{multicol}
\usepackage{multirow}
\usepackage{siunitx}
% parse-numbers=false : accepte aussi \SI{2,5 \times 10^{-3}}{...}
\sisetup{locale=FR,output-decimal-marker={,},group-minimum-digits=4,parse-numbers=false}
\DeclareSIUnit\ohm{\ensuremath{\symup{\Omega}}} % U+2126 absent de la police
\DeclareSIUnit\division{div} % réglages d'oscilloscope (ms/div, V/div)
\DeclareSIUnit\tr{tr} % tours (vitesse de rotation, ex. \SI{1500}{\tr\per\minute})
\DeclareSIUnit\molar{M} % concentration molaire (ex. \SI{3}{\molar}, \SI{1}{\micro\molar})
\DeclareSIUnit\an{an} % année (le modèle confond parfois avec \year, un compteur TeX) — ex. \SI{5}{\kilo\meter\per\an}
\DeclareSIUnit\FCFA{FCFA} % franc CFA (ex. \SI{500}{\FCFA\per\kilo\gram})
\DeclareSIUnit\degreeBrix{\textdegree Brix} % teneur en sucre d'un sirop/jus (réfractométrie)
\DeclareSIUnit\month{mois} % le modèle utilise parfois \month, un compteur TeX, comme unité de durée
\usepackage{qrcode}
% \degree : commande fréquemment utilisée par les modèles pour un angle ou une
% température en dehors de \SI{}{} (non fournie par les paquets chargés).
\providecommand{\degree}{^{\circ}}
% \qed (fin de démonstration), utilisé par les modèles sans amsthm
\providecommand{\qed}{\hfill$\square$}
% \barre : double barre des valeurs interdites dans un tableau de variations (tkz-tab)
\providecommand{\barre}{\ensuremath{\|}}
% environnement proof (amsthm non chargé) utilisé par les modèles
\ifdefined\proof\else\newenvironment{proof}[1][Démonstration]{\par\noindent\textit{#1.}\ }{\qed\par}\fi
\usepackage{lastpage}
\usepackage{hyperref}
\hypersetup{hidelinks}

% ============================================================
% Palette « Pop » OnBuch+ — mêmes hex que la classe P (plus_ui.dart)
% ============================================================
\definecolor{poporange}{HTML}{F59321}
\definecolor{popdark}{HTML}{E07A0C}
\definecolor{popcream}{HTML}{FFF6E8}
\definecolor{popink}{HTML}{1C1714}
\definecolor{poppurple}{HTML}{7A5AE0}
\definecolor{popgreen}{HTML}{2E9E6A}
\definecolor{poppink}{HTML}{E0457B}
\definecolor{popblue}{HTML}{2F7DE1}
\definecolor{popgold}{HTML}{C98A12}
\definecolor{popmuted}{HTML}{8A7F76}
\definecolor{popline}{HTML}{EADFD2}
\definecolor{popgray}{HTML}{8A7F76}
\colorlet{popgrey}{popgray}
% Alias tolérants (le modèle nomme parfois les couleurs différemment)
\colorlet{popwhite}{white}
\colorlet{popblack}{popink}
\colorlet{popred}{poppink}
\colorlet{popyellow}{popgold}
% Teintes claires (fonds de tableaux/encadrés) — le modèle nomme parfois
% les couleurs avec un suffixe L (ex. popblueL) sans les avoir définies.
\colorlet{poporangeL}{poporange!20}
\colorlet{popdarkL}{popdark!20}
\colorlet{poppurpleL}{poppurple!20}
\colorlet{popgreenL}{popgreen!20}
\colorlet{poppinkL}{poppink!20}
\colorlet{popblueL}{popblue!20}
\colorlet{popgoldL}{popgold!20}
\colorlet{popredL}{popred!20}
\colorlet{popgrayL}{popgray!20}
% Teintes claires pour fonds de tableaux / zones
\colorlet{popblueL}{popblue!10!white}
\colorlet{poporangeL}{poporange!14!white}
\colorlet{popgreenL}{popgreen!12!white}
\colorlet{poppurpleL}{poppurple!12!white}
\colorlet{poppinkL}{poppink!12!white}
\colorlet{popgoldL}{popgold!14!white}

\pagecolor{popcream}

% ============================================================
% Typographie
% ============================================================
\defaultfontfeatures{Ligatures=TeX}
\setmainfont{PlusJakartaSans-Medium.ttf}[Path=../../../fonts/,BoldFont=PlusJakartaSans-Bold.ttf]
% Symboles, API et emojis absents de la police principale : repli sur DejaVu Sans (caractères actifs)
\newfontface\symfont{DejaVuSans.ttf}[Path=../../../fonts/]
\makeatletter
\newcommand\symactive[1]{\catcode#1=13 \begingroup\lccode`\~=#1 \lowercase{\endgroup\def~}{{\symfont\char#1}}}
\newcommand\symblank[1]{\catcode#1=13 \begingroup\lccode`\~=#1 \lowercase{\endgroup\def~}{}}
\newcommand\symhyph[1]{\catcode#1=13 \begingroup\lccode`\~=#1 \lowercase{\endgroup\def~}{\mbox{-}}}
\newcount\symc
\newcommand\symrange[2]{\symc=#1 \@whilenum\symc<#2 \do{\expandafter\symactive\expandafter{\the\symc}\advance\symc by 1 }}
\newcommand\symblankrange[2]{\symc=#1 \@whilenum\symc<#2 \do{\expandafter\symblank\expandafter{\the\symc}\advance\symc by 1 }}
\makeatother
\symrange{"0250}{"02FF}\symactive{"2713}\symactive{"2714}\symactive{"2717}\symactive{"2718}\symactive{"2610}\symactive{"2611}\symactive{"2612}
\symhyph{"2011}\symhyph{"2E1A}\symblankrange{"1F300}{"1FAFF}\symblank{"FE0F}
% Maths en sans-serif (Fira Math) avec chiffres et lettres droites de Plus
% Jakarta, pour que les formules se fondent dans le texte.
\usepackage[math-style=ISO,bold-style=ISO]{unicode-math}
\setmathfont{FiraMath-Regular.otf}[Path=../../../fonts/]
\setmathfont{PlusJakartaSans-Medium.ttf}[Path=../../../fonts/,range={up/num,up/{Latin,latin},\mathup}]
\usepackage{icomma} % virgule décimale sans espace en mode math
% Caractères mathématiques Unicode absents des polices de texte (Plus Jakarta,
% Archivo Black) : rendus par leur équivalent mathématique, en texte comme en
% formule — sinon ils s'affichent en carré vide (ex. « Arithmétique dans ℤ »).
\usepackage{newunicodechar}
\newunicodechar{ℝ}{\ensuremath{\mathbb{R}}}
\newunicodechar{ℤ}{\ensuremath{\mathbb{Z}}}
\newunicodechar{ℂ}{\ensuremath{\mathbb{C}}}
\newunicodechar{ℕ}{\ensuremath{\mathbb{N}}}
\newunicodechar{ℚ}{\ensuremath{\mathbb{Q}}}
\newunicodechar{𝔻}{\ensuremath{\mathbb{D}}}
\newunicodechar{α}{\ensuremath{\alpha}}
\newunicodechar{β}{\ensuremath{\beta}}
\newunicodechar{γ}{\ensuremath{\gamma}}
\newunicodechar{δ}{\ensuremath{\delta}}
\newunicodechar{ε}{\ensuremath{\varepsilon}}
\newunicodechar{θ}{\ensuremath{\theta}}
\newunicodechar{λ}{\ensuremath{\lambda}}
\newunicodechar{μ}{\ensuremath{\mu}}
\newunicodechar{σ}{\ensuremath{\sigma}}
\newunicodechar{τ}{\ensuremath{\tau}}
\newunicodechar{φ}{\ensuremath{\varphi}}
\newunicodechar{ψ}{\ensuremath{\psi}}
\newunicodechar{ω}{\ensuremath{\omega}}
\newunicodechar{Σ}{\ensuremath{\Sigma}}
% majuscules produites par \MakeUppercase dans les titres (α-aminés -> Α)
\newunicodechar{Α}{\ensuremath{\alpha}}
\newunicodechar{Β}{\ensuremath{\beta}}
\newunicodechar{Γ}{\ensuremath{\Gamma}}
\newunicodechar{Δ}{\ensuremath{\Delta}}
\newunicodechar{Ε}{\ensuremath{\varepsilon}}
\newunicodechar{Θ}{\ensuremath{\Theta}}
\newunicodechar{Λ}{\ensuremath{\Lambda}}
\newunicodechar{Μ}{\ensuremath{\mu}}
\newunicodechar{Τ}{\ensuremath{\tau}}
\newunicodechar{Φ}{\ensuremath{\Phi}}
\newunicodechar{Ψ}{\ensuremath{\Psi}}
\newunicodechar{Ω}{\ensuremath{\Omega}}
\newunicodechar{↦}{\ensuremath{\mapsto}}
\newunicodechar{∈}{\ensuremath{\in}}
\newunicodechar{∉}{\ensuremath{\notin}}
\newunicodechar{∧}{\ensuremath{\wedge}}
\newunicodechar{⇔}{\ensuremath{\Leftrightarrow}}
\newunicodechar{⟺}{\ensuremath{\Longleftrightarrow}}
\newunicodechar{⇒}{\ensuremath{\Rightarrow}}
\newunicodechar{⟹}{\ensuremath{\Longrightarrow}}
\newunicodechar{∘}{\ensuremath{\circ}}
\newunicodechar{∪}{\ensuremath{\cup}}
\newunicodechar{∩}{\ensuremath{\cap}}
\newunicodechar{≡}{\ensuremath{\equiv}}
\newunicodechar{⊂}{\ensuremath{\subset}}
\newunicodechar{⊥}{\ensuremath{\perp}}
\newunicodechar{∃}{\ensuremath{\exists}}
\newunicodechar{∀}{\ensuremath{\forall}}
\newunicodechar{∛}{\ensuremath{\sqrt[3]{\,}}}
\newunicodechar{∜}{\ensuremath{\sqrt[4]{\,}}}
\newunicodechar{⃗}{\ensuremath{{}^{\to}}}
\newunicodechar{ⁿ}{\ifmmode^{n}\else\textsuperscript{\ensuremath{n}}\fi}
\newunicodechar{ˣ}{\ifmmode^{x}\else\textsuperscript{\ensuremath{x}}\fi}
\newunicodechar{ᵇ}{\ifmmode^{b}\else\textsuperscript{\ensuremath{b}}\fi}
\newunicodechar{ʳ}{\ifmmode^{r}\else\textsuperscript{\ensuremath{r}}\fi}
\newunicodechar{ᵘ}{\ifmmode^{u}\else\textsuperscript{\ensuremath{u}}\fi}
\newunicodechar{ʸ}{\ifmmode^{y}\else\textsuperscript{\ensuremath{y}}\fi}
\newunicodechar{ᵃ}{\ifmmode^{a}\else\textsuperscript{\ensuremath{a}}\fi}
\newunicodechar{ᵉ}{\ifmmode^{e}\else\textsuperscript{\ensuremath{e}}\fi}
\newunicodechar{ᵏ}{\ifmmode^{k}\else\textsuperscript{\ensuremath{k}}\fi}
\newunicodechar{ᵖ}{\ifmmode^{p}\else\textsuperscript{\ensuremath{p}}\fi}
\newunicodechar{ᴺ}{\ifmmode^{N}\else\textsuperscript{\ensuremath{N}}\fi}
\newunicodechar{ᶜ}{\ifmmode^{c}\else\textsuperscript{\ensuremath{c}}\fi}
\newunicodechar{ʰ}{\ifmmode^{h}\else\textsuperscript{\ensuremath{h}}\fi}
\newunicodechar{ᵠ}{\ifmmode^{q}\else\textsuperscript{\ensuremath{q}}\fi}
\newunicodechar{ₙ}{\ifmmode_{n}\else\textsubscript{\ensuremath{n}}\fi}
\newunicodechar{ₐ}{\ifmmode_{a}\else\textsubscript{\ensuremath{a}}\fi}
\newunicodechar{ᵢ}{\ifmmode_{i}\else\textsubscript{\ensuremath{i}}\fi}
\newunicodechar{ᵣ}{\ifmmode_{r}\else\textsubscript{\ensuremath{r}}\fi}
\newunicodechar{ₖ}{\ifmmode_{k}\else\textsubscript{\ensuremath{k}}\fi}
\newunicodechar{ᵦ}{\ifmmode_{\beta}\else\textsubscript{\ensuremath{\beta}}\fi}
\newunicodechar{ₓ}{\ifmmode_{x}\else\textsubscript{\ensuremath{x}}\fi}
\newfontfamily\popdisplay{ArchivoBlack-Regular.ttf}[Path=../../../fonts/]
\newfontfamily\poplogo{SpaceGrotesk-Bold.ttf}[Path=../../../fonts/]
\newfontfamily\popsemi{PlusJakartaSans-SemiBold.ttf}[Path=../../../fonts/]
\newfontfamily\popxbold{PlusJakartaSans-ExtraBold.ttf}[Path=../../../fonts/]
\newfontfamily\popxboldls{PlusJakartaSans-ExtraBold.ttf}[Path=../../../fonts/,LetterSpace=3.0]

\setlist[enumerate]{leftmargin=1.7em,itemsep=5pt,topsep=4pt}
\setlist[itemize]{leftmargin=1.7em,itemsep=4pt,topsep=4pt,label={\color{poporange}\textbullet}}
\setlength{\parindent}{0pt}
\setlength{\parskip}{5pt}
\renewcommand{\arraystretch}{1.25}

% pgfplots : style Pop par défaut
\pgfplotsset{
  every axis/.append style={axis line style={popink,line width=1pt},tick style={popink},
    grid style={popline,line width=0.5pt},label style={font=\small},tick label style={font=\footnotesize}},
  popcurve/.style={poporange,line width=1.8pt,no markers,smooth},
}
% Même style disponible dans un \draw TikZ simple (hors pgfplots)
\tikzset{popcurve/.style={poporange,line width=1.8pt}}
\tikzset{popgrid/.style={popline,very thin}}
\colorlet{popcurve}{poporange} % le nom du style est parfois utilisé comme couleur

% ============================================================
% Pill / squiggle
% ============================================================
\newcommand{\poppill}[3]{%
  \tikz[baseline=-0.55ex]\node[draw=popink,line width=1.2pt,rounded corners=2.6pt,%
    fill=#2,inner xsep=6.5pt,inner ysep=3pt]{%
    \popxboldls\fontsize{7}{7}\selectfont\color{#3}\MakeUppercase{#1}};%
}
\newcommand{\popsquiggle}[1]{%
  \tikz[baseline=-0.4ex]{
    \draw[#1,line width=2.6pt,line cap=round]
      plot [smooth] coordinates {(0,0.09) (0.34,0.01) (0.68,0.21) (1.02,0.01) (1.36,0.10)};
  }%
}
% Mot-clé surligné (marqueur orange sous le texte)
\newcommand{\cle}[1]{{\popxbold\color{popdark}#1}}

% ============================================================
% En-tête / pied
% ============================================================
\pagestyle{fancy}
\fancyhf{}
\renewcommand{\headrulewidth}{0pt}
\renewcommand{\footrulewidth}{0pt}
\lhead{\raisebox{-0.15ex}{\poplogo\fontsize{13}{13}\selectfont\color{popink}OnBuch\color{poporange}+}}
\rhead{\poppill{\DOCMATIERE\ \textbullet\ \DOCNIVEAU\ \textbullet\ Leçon \DOCLECON}{white}{popink}}
\lfoot{\popxbold\fontsize{7.6}{7.6}\selectfont\color{popmuted}\MakeUppercase{Cours signé OnBuch+}\ \textbullet\ {\popsemi\DOCTITRE}}
\rfoot{\tikz[baseline=-0.6ex]\node[rounded corners=8.5pt,fill=popink,minimum height=17pt,minimum width=17pt,inner xsep=5pt,inner sep=0]{\popxbold\fontsize{8}{8}\selectfont\color{popcream}\thepage\,/\,\pageref*{LastPage}};}

% ============================================================
% Titres
% ============================================================
\newcommand{\chaptitre}[2]{%
  \begin{tcolorbox}[enhanced,colback=poporange,colframe=popink,boxrule=2.6pt,arc=11pt,
      drop shadow=popink,left=15pt,right=15pt,top=12pt,bottom=11pt,before skip=3pt,after skip=18pt]
    \poppill{#2}{popink}{popcream}\par\vspace{9pt}
    {\raggedright\hyphenpenalty=10000\exhyphenpenalty=10000\popdisplay\fontsize{22}{24}\selectfont\color{popcream}\MakeUppercase{#1}\par}
  \end{tcolorbox}
}
% Section numérotée : pastille orange avec le numéro + titre Archivo
\newcounter{popsec}
\newcommand{\coursec}[1]{%
  \stepcounter{popsec}\setcounter{popsub}{0}%
  \par\vspace{16pt}\noindent
  \begin{minipage}{\linewidth}
  \tikz[baseline=-0.7ex]\node[draw=popink,line width=1.6pt,fill=poporange,rounded corners=4pt,minimum size=22pt,inner sep=0]{\popdisplay\fontsize{12}{12}\selectfont\color{popcream}\thepopsec};%
  \hspace{8pt}{\popdisplay\fontsize{15}{17}\selectfont\color{popink}\MakeUppercase{#1}}\par
  \vspace{4pt}\noindent{\color{poporange}\rule{36pt}{3.6pt}}
  \end{minipage}\par\nopagebreak\vspace{8pt}
}
\newcounter{popsub}
\newcommand{\courssub}[1]{%
  \stepcounter{popsub}%
  \par\vspace{9pt}\noindent{\popxbold\fontsize{11.5}{13}\selectfont\color{popink}{\color{poporange}\thepopsec.\thepopsub}\hspace{6pt}#1}\par\nopagebreak\vspace{4pt}
}

% ============================================================
% Boîtes pédagogiques — même recette : fond blanc, bordure épaisse colorée,
% ombre dure encre, badge plein. Titre optionnel [#1].
% ============================================================
\tcbset{popbase/.style={breakable,enhanced,colback=white,boxrule=2.3pt,arc=9pt,
  shadow={3pt}{-3pt}{0pt}{popink},left=14pt,right=14pt,top=11pt,bottom=11pt,before skip=11pt,after skip=15pt,
  fontupper=\popsemi\fontsize{10.6}{14.5}\selectfont\color{popink},
  fontlower=\popsemi\fontsize{10.6}{14.5}\selectfont\color{popink}}}
% Coche dessinée (le glyphe ✓ n'existe pas dans les polices du document)
\newcommand{\popcheck}[1]{\tikz[baseline=-0.5ex]\draw[#1,line width=1.6pt,line cap=round,line join=round](-0.13,0.0)--(-0.04,-0.09)--(0.13,0.1);}
\providecommand{\coche}{\popcheck{popgreen}}
\providecommand{\resultat}[1]{\par\smallskip\noindent\fcolorbox{popgreen}{popgreenL}{\parbox{\dimexpr\linewidth-2\fboxsep-2\fboxrule\relax}{\textbf{Résultat.} #1}}\par\smallskip}
\newcommand{\popboxhead}[3]{\poppill{#1}{#2}{white}\ifx\relax#3\relax\else\hspace{6pt}{\popxbold\fontsize{10.6}{12}\selectfont\color{popink}#3}\fi\par\vspace{7pt}}

\newtcolorbox{aretenir}[1][]{popbase,colframe=popgreen,before upper={\popboxhead{À retenir}{popgreen}{#1}}}
% Texte en anglais (document de lecture, dialogue, extrait) : cadre
% d'étude. Un titre optionnel (source, auteur) s'affiche dans l'en-tête.
\newtcolorbox{textebox}[1][]{popbase,colframe=popblue,before upper={\popboxhead{Text}{popblue}{#1}\begin{otherlanguage}{english}},after upper={\end{otherlanguage}}}
% Marques ✓ / ✗ (forme correcte / fautive) souvent utilisées par les modèles
\providecommand{\cross}{\textcolor{poppink}{\ensuremath{\times}}}
\providecommand{\xmark}{\cross}\providecommand{\crossmark}{\cross}\providecommand{\cmark}{\ensuremath{\checkmark}}
% Ligne de réponse / justification à compléter (inventée par les modèles)
\providecommand{\justification}[1]{\par\noindent\textit{Justification~:} \dotfill\par}
% \textipa (package tipa non chargé) : la police ne couvre pas l'API -> simple passage
\providecommand{\textipa}[1]{#1}
% Soulignement (ulem non chargé) : \uline, \ul
\providecommand{\uline}[1]{\underline{#1}}\providecommand{\ul}[1]{\underline{#1}}
% \glossaire{\item ...} inventée par les modèles : liste de glossaire
\providecommand{\glossaire}[1]{\begin{itemize}#1\end{itemize}}
% \alert (beamer) inventée par les modèles : texte mis en évidence
\providecommand{\alert}[1]{\textbf{\textcolor{poppink}{#1}}}
% variantes ASCII d'ordinaux écrites en macro
\providecommand{\iere}{\textsuperscript{re}}\providecommand{\ieme}{\textsuperscript{e}}\providecommand{\ere}{\textsuperscript{re}}\providecommand{\eme}{\textsuperscript{e}}\providecommand{\er}{\textsuperscript{er}}\providecommand{\ie}{\textsuperscript{e}}
% Ordinaux français écrits en macro par les modèles : 1\ière, 3\ième
\newcommand{\ière}{\textsuperscript{re}}\newcommand{\ième}{\textsuperscript{e}}
% \exemple (inventée par les modèles) : étiquette « Exemple : »
\providecommand{\exemple}{\textbf{Exemple~:}\ }
% \square utilisable aussi hors mode math (cases à cocher)
\AtBeginDocument{\let\squareMath\square\renewcommand{\square}{\ensuremath{\squareMath}}}
% Mot ou phrase en anglais à l'intérieur d'un paragraphe français
\newcommand{\eng}[1]{\ifmmode\text{\textit{#1}}\else\begingroup\selectlanguage{english}\itshape #1\endgroup\fi}
\let\esp\eng % alias toléré
% flèches d'intonation inventées par les modèles
\providecommand{\textsearrow}{}\renewcommand{\textsearrow}{\ensuremath{\searrow}}\providecommand{\textnearrow}{}\renewcommand{\textnearrow}{\ensuremath{\nearrow}}
% \fr{..} : mot français (inventé par les modèles)
\providecommand{\fr}[1]{\textit{#1}}
\newtcolorbox{exemplebox}[1][]{popbase,colframe=poporange,before upper={\popboxhead{Exemple}{poporange}{#1}}}
\newtcolorbox{definition}[1][]{popbase,colframe=popblue,before upper={\popboxhead{Définition}{popblue}{#1}}}
\newtcolorbox{propriete}[1][]{popbase,colframe=poppurple,before upper={\popboxhead{Propriété}{poppurple}{#1}}}
\newtcolorbox{methode}[1][]{popbase,colframe=popgold,before upper={\popboxhead{Méthode}{popgold}{#1}}}
\newtcolorbox{attention}[1][]{popbase,colframe=poppink,before upper={\popboxhead{Attention}{poppink}{#1}}}
\newtcolorbox{experience}[1][]{popbase,colframe=popblue,colback=popblueL,before upper={\popboxhead{Expérience}{popblue}{#1}}}
% « Le savais-tu ? » : carte violette pleine (contexte, culture, Cameroun)
\newtcolorbox{savaistu}[1][]{popbase,colback=poppurple,colframe=popink,
  fontupper=\popsemi\fontsize{10.4}{14.2}\selectfont\color{white},
  before upper={\poppill{Le savais-tu\,?}{white}{poppurple}\ifx\relax#1\relax\else\hspace{6pt}{\popxbold\color{white}#1}\fi\par\vspace{7pt}}}
% Exercice résolu : énoncé en haut, \tcblower puis la solution sur fond crème
\newcounter{popexo}\newcounter{popexr}
\newtcolorbox{exoresolu}[1][]{popbase,colframe=poporange,colbacklower=popcream,
  segmentation style={popink,line width=1.2pt,dashed},
  before upper={\stepcounter{popexr}\popboxhead{Exercice résolu \thepopexr}{poporange}{#1}},
  before lower={\poppill{Solution}{popink}{popcream}\par\vspace{6pt}}}
% Exercice (non résolu en place) — la correction est dans la partie Corrigés
\newtcolorbox{exercice}[1][]{popbase,colframe=popink,
  before upper={\stepcounter{popexo}\popboxhead{Exercice \thepopexo}{popink}{#1}}}
% Corrigé
\newtcolorbox{corrige}[1][]{popbase,colframe=popgreen,colback=popgreenL,
  before upper={\popboxhead{Corrigé}{popgreen}{#1}}}
% Objectifs / prérequis / bilan
\newtcolorbox{objectifs}{popbase,colframe=popink,colback=poporangeL,
  before upper={\popboxhead{Objectifs de la leçon}{popink}{}}}
\newtcolorbox{prerequis}{popbase,colframe=popink,colback=popblueL,
  before upper={\popboxhead{Prérequis}{popblue}{}}}
\newtcolorbox{bilan}{popbase,colframe=popink,colback=white,boxrule=2.8pt,
  before upper={\popboxhead{Fiche bilan}{poporange}{L'essentiel à connaître pour l'examen}}}
% Figure encadrée avec légende
\newtcolorbox{popfigure}[1][]{enhanced,breakable=false,colback=white,colframe=popline,boxrule=1.4pt,arc=8pt,
  left=10pt,right=10pt,top=10pt,bottom=8pt,before skip=10pt,after skip=12pt,halign=center,
  lower separated=false,halign lower=center,
  fontlower=\popsemi\fontsize{9}{11.5}\selectfont\color{popmuted},
  #1}
\newcommand{\legende}[1]{\par\vspace{6pt}{\popsemi\fontsize{9}{11.5}\selectfont\color{popmuted}#1}}

% ============================================================
% Signature OnBuch+
% ============================================================
\newcommand{\poplogoinline}[1]{{\poplogo\fontsize{#1}{#1}\selectfont\color{popink}OnBuch\color{poporange}+}}
\newcommand{\signatureonbuch}{%
  \begin{tcolorbox}[enhanced,colback=popink,colframe=popink,boxrule=2.2pt,arc=10pt,
      drop shadow=poporange,left=14pt,right=14pt,top=10pt,bottom=10pt,before skip=0pt,after skip=0pt]
    \tikz[baseline=-0.7ex]\node[circle,fill=popgreen,draw=popcream,line width=1.3pt,minimum size=22pt,inner sep=0]{\popcheck{white}};%
    \hspace{9pt}\begin{minipage}[c]{0.62\linewidth}
      {\popxbold\fontsize{12}{13}\selectfont\color{popcream}Signé\ \textbullet\ {\poplogo OnBuch\color{poporange}+}}\par\vspace{2pt}
      {\raggedright\popsemi\fontsize{8}{10.5}\selectfont\color{popline}\MakeUppercase{Équipe pédagogique OnBuch+}\\\MakeUppercase{Conforme au programme officiel MINESEC}\par}
    \end{minipage}\hfill
    {\poplogo\fontsize{22}{22}\selectfont\color{popcream}OnBuch\color{poporange}+}
  \end{tcolorbox}%
}

% ============================================================
% Page de garde
% #1 titre · #2 matière · #3 niveau · #4 module · #5 n° leçon · #6 durée ·
% #7 description / objectifs (texte ou itemize)
% ============================================================
\newcommand{\pagegarde}[7]{%
  \thispagestyle{empty}
  \newgeometry{a4paper,margin=1.7cm,top=1.6cm,bottom=1.4cm}%
  \begin{tikzpicture}[remember picture,overlay]
    \fill[poporange!18] ([xshift=-3.2cm,yshift=-3.6cm]current page.north east) circle (4.6cm);
    \fill[poppurple!14] ([xshift=2.4cm,yshift=7.6cm]current page.south west) circle (3.8cm);
  \end{tikzpicture}%
  {\poplogo\fontsize{30}{30}\selectfont\color{popink}OnBuch\color{poporange}+}\hfill
  \raisebox{0.6ex}{\poppill{Cours signé \textbullet\ Premium}{popink}{popcream}}\par
  \vspace{1pt}\popsquiggle{poppurple}\par
  \vspace{8pt}
  \begin{tcolorbox}[enhanced,colback=poporange,colframe=popink,boxrule=2.8pt,arc=13pt,
      drop shadow=popink,left=18pt,right=18pt,top=12pt,bottom=13pt,after skip=10pt]
    \poppill{#2 \textbullet\ #3}{popink}{popcream}\hspace{5pt}\poppill{#4}{white}{popink}\par\vspace{8pt}
    {\popdisplay\fontsize{14}{14}\selectfont\color{popink}LEÇON #5}\par\vspace{4pt}
    {\raggedright\hyphenpenalty=10000\exhyphenpenalty=10000\popdisplay\fontsize{25}{27}\selectfont\color{popcream}\MakeUppercase{#1}\par}
    \vspace{8pt}
    {\popxbold\fontsize{9.5}{9.5}\selectfont\color{popink}\MakeUppercase{Durée conseillée : #6}}
  \end{tcolorbox}
  \begin{tcolorbox}[enhanced,colback=white,colframe=popink,boxrule=2.2pt,arc=10pt,
      drop shadow=popink,left=16pt,right=16pt,top=10pt,bottom=10pt,after skip=8pt]
    \poppill{Dans cette leçon}{poppurple}{white}\par\vspace{5pt}
    \popsemi\fontsize{9.3}{12.6}\selectfont\color{popink}#7
  \end{tcolorbox}
  \noindent\begin{minipage}[t]{0.487\linewidth}
    \begin{tcolorbox}[enhanced,colback=popgreen,colframe=popink,boxrule=2.2pt,arc=10pt,
        drop shadow=popink,left=11pt,right=11pt,top=10pt,bottom=10pt,height=3.1cm,valign=center]
      \tcbox[colback=white,colframe=popink,boxrule=1.3pt,arc=5pt,boxsep=0pt,nobeforeafter,
        left=3pt,right=3pt,top=3pt,bottom=3pt,valign=center]{\qrcode[height=1.9cm]{https://play.google.com/store/apps/details?id=com.luvvix.onbuch&hl=fr}}%
      \hspace{8pt}\begin{minipage}[c]{0.5\linewidth}
        {\popdisplay\fontsize{11}{12.5}\selectfont\color{white}\MakeUppercase{Télécharge OnBuch}}\par\vspace{4pt}
        {\raggedright\popsemi\fontsize{8.4}{11}\selectfont\color{white}Gratuit \textbullet\ Google Play\\Tuteur IA, annales, concours, résultats\par}
      \end{minipage}
    \end{tcolorbox}
  \end{minipage}\hfill
  \begin{minipage}[t]{0.487\linewidth}
    \begin{tcolorbox}[enhanced,colback=poppurple,colframe=popink,boxrule=2.2pt,arc=10pt,
        drop shadow=popink,left=11pt,right=11pt,top=10pt,bottom=10pt,height=3.1cm,valign=center]
      \tikz[baseline=-0.7ex]\node[circle,fill=white,draw=popink,line width=1.3pt,minimum size=1.6cm,inner sep=0]{\popdisplay\fontsize{24}{24}\selectfont\color{poppurple}+};%
      \hspace{8pt}\begin{minipage}[c]{0.6\linewidth}
        {\popdisplay\fontsize{11}{12.5}\selectfont\color{white}\MakeUppercase{Passe à OnBuch+}}\par\vspace{4pt}
        {\raggedright\popsemi\fontsize{8.4}{11}\selectfont\color{white}Tous les cours signés, Tuteur IA illimité, fiches PDF — onglet OnBuch+ de l'app\par}
      \end{minipage}
    \end{tcolorbox}
  \end{minipage}
  \vspace{8pt}
  \popcontenu
  \vfill
  \signatureonbuch
  \restoregeometry
  \newpage
}

% Rangée « Ce cours contient » (page de garde)
\newcommand{\poptile}[3]{% #1 couleur #2 gros texte #3 légende
  \begin{tcolorbox}[enhanced,colback=#1,colframe=popink,boxrule=2pt,arc=9pt,shadow={2.5pt}{-2.5pt}{0pt}{popink},
      left=6pt,right=6pt,top=9pt,bottom=9pt,halign=center,valign=center,height=1.75cm,width=\linewidth,nobeforeafter]
    {\popdisplay\fontsize{12.5}{13}\selectfont\color{white}\MakeUppercase{#2}}\par\vspace{3pt}
    {\centering\popsemi\fontsize{7.8}{9.5}\selectfont\color{white}#3\par}
  \end{tcolorbox}}
\newcommand{\popcontenu}{%
  \par\noindent{\popxboldls\fontsize{7.5}{7.5}\selectfont\color{popmuted}CE COURS SIGNÉ CONTIENT}\par\vspace{7pt}
  \noindent\begin{minipage}[t]{0.235\linewidth}\poptile{popblue}{Cours}{complet, illustré, pas à pas}\end{minipage}\hfill
  \begin{minipage}[t]{0.235\linewidth}\poptile{poporange}{Méthodes}{et exercices résolus}\end{minipage}\hfill
  \begin{minipage}[t]{0.235\linewidth}\poptile{poppink}{Exercices}{type examen + corrigés}\end{minipage}\hfill
  \begin{minipage}[t]{0.235\linewidth}\poptile{popgreen}{Fiche bilan}{l'essentiel à retenir}\end{minipage}\par}

% Sommaire
\newtcolorbox{sommaire}{enhanced,colback=white,colframe=popink,boxrule=2.2pt,arc=10pt,
  drop shadow=popink,left=18pt,right=18pt,top=13pt,bottom=13pt,before skip=6pt,after skip=20pt,
  before upper={\poppill{Sommaire}{poppurple}{white}\par\vspace{9pt}\popsemi\fontsize{10.6}{14.5}\selectfont\color{popink}}}

% Signature de fin de cours — poussée tout en bas de la dernière page
\newcommand{\findecours}{%
  \par\vfill\nopagebreak
  \begin{center}\popsquiggle{poporange}\end{center}\vspace{4pt}
  \signatureonbuch
}
\AtBeginDocument{\def\llbracket{\mathopen{[\mkern-3.5mu[}}\def\rrbracket{\mathclose{]\mkern-3.5mu]}}} % intervalles d'entiers (glyphe ⟦ absent de la police)
% Glyphes absents de Fira Math : redéfinis à partir de glyphes disponibles
\AtBeginDocument{%
  \def\setminus{\mathbin{\backslash}}\let\smallsetminus\setminus
  \def\vdots{\mathord{\vbox{\baselineskip3.2pt\lineskiplimit0pt\kern4pt\hbox{.}\hbox{.}\hbox{.}}}}%
  \def\ddots{\mathinner{\mkern1mu\raise6.5pt\hbox{.}\mkern2mu\raise3.6pt\hbox{.}\mkern2mu\raise0.7pt\hbox{.}\mkern1mu}}%
  \def\triangle{\mathord{\scalebox{0.9}{$\symup{\Delta}$}}}%
  \def\nabla{\mathord{\raisebox{\depth}{\scalebox{1}[-1]{$\symup{\Delta}$}}}}%
  \def\checkmark{\popcheck{popdark}}%
  \def\star{\mathbin{\ast}}\def\bigstar{\mathord{\ast}}%
  \def\diamond{\mathbin{\scalebox{0.6}{$\Diamond$}}}%
  \def\bigcup{\mathop{\mathchoice{\raisebox{-0.3em}{\scalebox{1.7}{$\cup$}}}{\scalebox{1.25}{$\cup$}}{\cup}{\cup}}\displaylimits}%
  \def\bigcap{\mathop{\mathchoice{\raisebox{-0.3em}{\scalebox{1.7}{$\cap$}}}{\scalebox{1.25}{$\cap$}}{\cap}{\cap}}\displaylimits}%
  \def\lesseqqgtr{\mathrel{\lessgtr}}\def\bigoplus{\mathop{\oplus}\displaylimits}%
}
\providecommand{\ssi}{si et seulement si} % abréviation fréquente
\AtBeginDocument{\providecommand{\wideparen}{\overparen}} % arc de cercle

\begin{document}

\pagegarde{Lesson 19 — Vocabulary: Words and expressions related to requesting for services}{Anglais}{1ère A}{Chapter 2 : Economic Life and Occupations}{ENA19}{2 h}{Cette leçon de vocabulaire présente les mots, expressions et formules polies pour demander un service (banque, hôtel, hôpital, administration, téléphone). Elle développe les champs lexicaux, les collocations, les faux amis et la réutilisation du lexique dans des dialogues et courts textes.
\vspace{4pt}
\begin{multicols}{2}\small
\begin{itemize}
\item À la fin de cette leçon, je sais identifier et définir le vocabulaire lié à la demande de services (request, assistance, appointment, booking...).
\item À la fin de cette leçon, je sais employer les formules de politesse adaptées pour demander un service (Could you...?, I would like to..., Would you mind...?).
\item À la fin de cette leçon, je sais utiliser les collocations correctes (make a request, book a room, fill in a form, pay a fee).
\item À la fin de cette leçon, je sais éviter les faux amis et les calques du français (demander ≠ to demand, assister ≠ to assist).
\item À la fin de cette leçon, je sais former des mots de la même famille (serve, service, servant, helpful, helpless).
\item À la fin de cette leçon, je sais prononcer et orthographier correctement les mots clés du thème.
\end{itemize}\end{multicols}}

\begin{sommaire}
\begin{multicols}{2}
\begin{enumerate}
\item Champ lexical 1 : Demander un service
\item Champ lexical 2 : Formules fonctionnelles de la requête
\item Collocations et expressions idiomatiques
\item Faux amis et pièges des francophones
\item Familles de mots, prononciation et orthographe
\item Le lexique en contexte : texte et dialogues
\item Activité d'intégration
\item Méthodes et exercices résolus
\item Exercices
\item Corrigés des exercices
\item Fiche bilan
\end{enumerate}
\end{multicols}
\end{sommaire}

\begin{objectifs}
\begin{itemize}
\item À la fin de cette leçon, je sais identifier et définir le vocabulaire lié à la demande de services (request, assistance, appointment, booking...).
\item À la fin de cette leçon, je sais employer les formules de politesse adaptées pour demander un service (Could you...?, I would like to..., Would you mind...?).
\item À la fin de cette leçon, je sais utiliser les collocations correctes (make a request, book a room, fill in a form, pay a fee).
\item À la fin de cette leçon, je sais éviter les faux amis et les calques du français (demander ≠ to demand, assister ≠ to assist).
\item À la fin de cette leçon, je sais former des mots de la même famille (serve, service, servant, helpful, helpless).
\item À la fin de cette leçon, je sais prononcer et orthographier correctement les mots clés du thème.
\item À la fin de cette leçon, je sais réemployer ce lexique dans des phrases, des dialogues et un court texte écrit.
\item À la fin de cette leçon, je sais réagir et répondre à une demande de service (acceptation, refus poli, renseignement).
\end{itemize}
\end{objectifs}

\begin{prerequis}
\begin{itemize}
\item Les auxiliaires modaux de base : can, could, may, would (Seconde).
\item Le vocabulaire général de la vie quotidienne : money, bank, hotel, hospital, office.
\item La structure d'une phrase interrogative en anglais (inversion sujet-auxiliaire).
\item Les formules de salutation et de politesse simples (Hello, please, thank you).
\end{itemize}
\end{prerequis}

\newpage

\coursec{Champ lexical 1 : Demander un service}

\textbf{Situation d'accroche.} Imagine : tu arrives à la banque de Yaoundé pour ouvrir un compte, mais le guichetier ne parle que l'anglais. Comment demander poliment un service ? Au Nigeria ou au Ghana, savoir formuler une requête — \eng{Could you help me, please?} — ouvre toutes les portes : hôtel, hôpital, agence de voyage, administration. Dans cette leçon, tu vas apprendre les mots et expressions exacts pour demander un service en anglais, comme un vrai anglophone.

\textbf{Problématique.} Quels mots, quelles expressions et quelles formules permettent de demander un service en anglais, correctement et poliment, sans calquer le français ?

\courssub{Observer d'abord : un texte pour découvrir le lexique}

Avant toute liste de mots, lisons une courte scène de la vie quotidienne. Relève tous les mots qui se rapportent à la demande de service.

\begin{textebox}[At the bank in Yaoundé]
Last Monday, Amina went to a bank on Kennedy Avenue to open a savings account. At the entrance, a customer care officer greeted her and asked how he could help. Amina made a polite request: she wanted to open an account and apply for a bank card. The officer gave her a form to fill in and told her there was a small fee for the card. She also needed to book an appointment with the accounts manager, so the receptionist checked the diary and offered her Thursday morning. Amina thanked him and asked for a receipt for the fee she had paid. The next day, her brother complained at the same bank because his card did not work, but the agent at the counter solved the problem in ten minutes. Both of them agreed that the service was fast and the staff were very helpful. When you need assistance, do not shout or demand things rudely: simply ask for help politely, and people will be happy to serve you.
\end{textebox}

\textbf{Glossaire.}

\begin{center}
\begin{tabularx}{\linewidth}{|l|l|X|}
\hline
\rowcolor{popblueL} Mot anglais & Nature & Traduction française \\
\hline
customer care officer & nom & agent du service client \\
\hline
request & nom & demande (formelle) \\
\hline
to apply for & verbe & faire une demande de, solliciter \\
\hline
form & nom & formulaire \\
\hline
fee & nom & frais \\
\hline
to book & verbe & réserver \\
\hline
appointment & nom & rendez-vous \\
\hline
receipt & nom & reçu \\
\hline
to complain & verbe & se plaindre \\
\hline
counter & nom & guichet, comptoir \\
\hline
staff & nom & personnel \\
\hline
assistance & nom & aide, assistance \\
\hline
\end{tabularx}
\end{center}

\textbf{Questions de compréhension.}
\begin{enumerate}
\item Why did Amina go to the bank?
\item What did the customer care officer give her?
\item Who did she book an appointment with?
\item Why did her brother complain the next day?
\end{enumerate}

\begin{corrige}[Questions de compréhension]
1. \eng{She went to the bank to open a savings account.}

2. \eng{He gave her a form to fill in.}

3. \eng{She booked an appointment with the accounts manager.}

4. \eng{He complained because his bank card did not work.}
\end{corrige}

\courssub{Le mot central : \eng{request}}

\begin{definition}[Request]
\eng{Request} (nom et verbe) signifie \cle{a polite or formal demand for something} : une demande polie ou formelle de quelque chose. C'est le mot roi de cette leçon.
\end{definition}

\begin{exemplebox}[Observer]
\eng{She made a request for a new passport.} = Elle a fait une demande de nouveau passeport.

\eng{I would like to request a meeting with the manager.} = Je voudrais demander un entretien avec le directeur.

\eng{Your request has been approved.} = Votre demande a été approuvée.
\end{exemplebox}

\begin{propriete}[Formes et emplois de \eng{request}]
\begin{itemize}
\item \textbf{Nom} : \eng{a request} ; construction : \eng{to make a request for something} (faire une demande de quelque chose). Exemple : \eng{He made a request for information.} = Il a fait une demande d'informations.
\item \textbf{Verbe} : \eng{to request something} ou \eng{to request that + sujet + base verbale} (registre formel). Exemple : \eng{We request that every customer fill in this form.} = Nous demandons que chaque client remplisse ce formulaire.
\item \textbf{Prononciation} : « ri-kouèste » ; l'accent tombe sur la deuxième syllabe, pour le nom comme pour le verbe.
\end{itemize}
\end{propriete}

\begin{attention}[\eng{Request} n'est pas \eng{question}]
Un francophone confond souvent les deux mots. \eng{A request} est une \cle{demande} (on veut obtenir quelque chose) ; \eng{a question} est une \cle{interrogation} (on veut savoir quelque chose).

\eng{I have a request: could you open the window?} = J'ai une demande : pourriez-vous ouvrir la fenêtre ?

\eng{I have a question: what time does the bank close?} = J'ai une question : à quelle heure ferme la banque ?

On ne dit jamais \eng{to make a question} ; on dit \eng{to ask a question}.
\end{attention}

\courssub{Les verbes clés de la demande de service}

Observons d'abord ces phrases, puis dégageons les emplois.

\begin{exemplebox}[Les verbes en action]
\eng{I would like to book a room for two nights.} = Je voudrais réserver une chambre pour deux nuits.

\eng{He applied for a loan at the bank.} = Il a demandé un prêt à la banque.

\eng{Could I order a taxi, please?} = Pourrais-je commander un taxi, s'il vous plaît ?

\eng{We hired a car for the weekend.} = Nous avons loué une voiture pour le week-end.

\eng{The customer complained about the slow service.} = Le client s'est plaint de la lenteur du service.
\end{exemplebox}

\begin{propriete}[Emploi et construction des verbes clés]
\begin{center}
\begin{tabularx}{\linewidth}{|l|X|X|}
\hline
\rowcolor{popblueL} Verbe & Sens et construction & Exemple \\
\hline
to request & demander (formel) : \eng{request sth} & \eng{She requested a receipt.} \\
\hline
to ask for & demander (courant) : \eng{ask for sth} & \eng{He asked for help.} \\
\hline
to require & avoir besoin de, exiger (formel) : \eng{require sth} & \eng{This form requires a signature.} \\
\hline
to need & avoir besoin de : \eng{need sth / need to do} & \eng{I need a pen.} \\
\hline
to apply for & solliciter (emploi, visa, prêt) : \eng{apply for sth} & \eng{She applied for a visa.} \\
\hline
to book & réserver (chambre, billet) : \eng{book sth} & \eng{We booked two tickets.} \\
\hline
to order & commander (repas, taxi) : \eng{order sth} & \eng{He ordered a meal.} \\
\hline
to hire & louer, engager : \eng{hire sth/sb} & \eng{They hired a guide.} \\
\hline
to complain & se plaindre : \eng{complain about sth, complain to sb} & \eng{She complained to the manager about the noise.} \\
\hline
\end{tabularx}
\end{center}
\end{propriete}

\begin{attention}[Pièges de construction]
\begin{itemize}
\item \eng{To ask for} : la préposition \eng{for} est obligatoire devant la chose demandée. On dit \eng{He asked for a pen} et non \eng{He asked a pen}. En revanche, \eng{to ask somebody} (sans \eng{for}) signifie interroger quelqu'un : \eng{Ask the officer.} = Demande à l'agent.
\item \eng{To apply for} : toujours avec \eng{for}. \eng{She applied for a job} et non \eng{she applied a job}. Attention à l'orthographe du passé : \eng{applied} (le \eng{y} devient \eng{i} devant \eng{-ed}).
\item \eng{To complain} : jamais de complément direct. On dit \eng{complain about the service} (se plaindre du service), \eng{complain to the manager} (se plaindre au directeur).
\end{itemize}
\end{attention}

\courssub{Les noms clés : objets et documents du service}

\begin{definition}[Noms essentiels]
\begin{itemize}
\item \eng{a service} : un service. \eng{This hotel offers excellent service.} = Cet hôtel offre un excellent service.
\item \eng{assistance} : l'aide, l'assistance (nom indénombrable, jamais de pluriel). \eng{Thank you for your assistance.} = Merci de votre aide.
\item \eng{help} : l'aide (indénombrable). \eng{Do you need any help?} = Avez-vous besoin d'aide ?
\item \eng{an appointment} : un rendez-vous. \eng{I have an appointment at ten o'clock.} = J'ai un rendez-vous à dix heures.
\item \eng{a booking / a reservation} : une réservation. \eng{I made a booking online.} = J'ai fait une réservation en ligne.
\item \eng{an application} : une demande, une candidature. \eng{Her application was accepted.} = Sa demande a été acceptée.
\item \eng{a form} : un formulaire. \eng{Please fill in this form.} = Veuillez remplir ce formulaire.
\item \eng{a fee} : des frais (toujours au singulier pour un tarif précis). \eng{There is a fee of 500 francs.} = Il y a des frais de 500 francs.
\item \eng{a receipt} : un reçu. \eng{Keep your receipt.} = Gardez votre reçu.
\end{itemize}
\end{definition}

\begin{attention}[Indénombrables et orthographe]
\eng{Assistance} et \eng{help} sont \cle{indénombrables} : jamais de \eng{an assistance} ni de \eng{helps} au pluriel. On dit \eng{some help}, \eng{a lot of assistance}. Attention aussi à l'orthographe de \eng{receipt} : le \eng{p} ne se prononce pas (« ri-site »), comme dans \eng{reception}.
\end{attention}

\courssub{Les personnes et les lieux de service}

\begin{definition}[Les personnes]
\begin{itemize}
\item \eng{a customer} : un client (d'un magasin, d'une banque). \eng{The customer is always right.} = Le client a toujours raison.
\item \eng{a client} : un client (d'un professionnel : avocat, agence). \eng{The lawyer met her client.} = L'avocate a rencontré son client.
\item \eng{a receptionist} : un(e) réceptionniste. \eng{The receptionist gave me the key.} = La réceptionniste m'a donné la clé.
\item \eng{an officer} : un agent, un fonctionnaire. \eng{The immigration officer checked my passport.} = L'agent d'immigration a vérifié mon passeport.
\item \eng{an agent} : un agent (commercial, de voyage). \eng{The travel agent booked our flight.} = L'agent de voyages a réservé notre vol.
\item \eng{a customer care officer} : un agent du service client. \eng{The customer care officer answered my questions.} = L'agent du service client a répondu à mes questions.
\end{itemize}
\end{definition}

\begin{definition}[Les lieux de services]
\begin{itemize}
\item \eng{a bank} : une banque. \eng{The bank opens at eight.} = La banque ouvre à huit heures.
\item \eng{a post office} : un bureau de poste. \eng{I sent a parcel at the post office.} = J'ai envoyé un colis à la poste.
\item \eng{a hotel} : un hôtel. \eng{We stayed at a hotel in Buea.} = Nous avons séjourné dans un hôtel à Buea.
\item \eng{a hospital} : un hôpital. \eng{She works at the central hospital.} = Elle travaille à l'hôpital central.
\item \eng{an embassy} : une ambassade. \eng{The embassy issues visas.} = L'ambassade délivre des visas.
\item \eng{a travel agency} : une agence de voyages. \eng{The travel agency planned our trip.} = L'agence de voyages a organisé notre voyage.
\item \eng{a customer care centre} : un centre de service client. \eng{Call the customer care centre.} = Appelez le centre de service client.
\end{itemize}
\end{definition}

\courssub{Carte mentale du champ lexical}

\begin{popfigure}
\begin{center}
\begin{tikzpicture}[mindmap, grow cyclic, every node/.style={concept}, root concept/.append style={concept color=popblue, font=\small\bfseries}, level 1/.append style={level distance=3.2cm, sibling angle=90}, level 2/.append style={level distance=2.4cm, sibling angle=45, font=\scriptsize}]
\node[root concept]{REQUESTING SERVICES}
  child[concept color=poporange]{ node {Verbs}
    child { node {request, ask for} }
    child { node {book, order, hire} }
    child { node {apply for, complain} }
  }
  child[concept color=popgreen]{ node {Nouns}
    child { node {request, service} }
    child { node {appointment, receipt} }
    child { node {form, fee} }
  }
  child[concept color=poppurple]{ node {People}
    child { node {customer, client} }
    child { node {receptionist, agent} }
    child { node {officer} }
  }
  child[concept color=poppink]{ node {Places}
    child { node {bank, hotel} }
    child { node {embassy, hospital} }
    child { node {post office} }
  };
\end{tikzpicture}
\end{center}
\legende{Carte mentale du champ lexical de la demande de service : quatre familles de mots.}
\end{popfigure}

\courssub{Tableau récapitulatif : 20 mots essentiels}

\begin{aretenir}[Vingt mots à connaître par cœur]
\begin{center}
\begin{tabularx}{\linewidth}{|l|l|X|}
\hline
\rowcolor{popblueL} English & Français & Exemple court \\
\hline
request & demande & \eng{a written request} \\
\hline
to ask for & demander & \eng{ask for information} \\
\hline
to require & exiger, nécessiter & \eng{a visa is required} \\
\hline
to need & avoir besoin de & \eng{I need help} \\
\hline
to apply for & solliciter & \eng{apply for a job} \\
\hline
to book & réserver & \eng{book a room} \\
\hline
to order & commander & \eng{order a taxi} \\
\hline
to hire & louer, engager & \eng{hire a car} \\
\hline
to complain & se plaindre & \eng{complain about the bill} \\
\hline
service & service & \eng{good service} \\
\hline
assistance & aide & \eng{ask for assistance} \\
\hline
appointment & rendez-vous & \eng{make an appointment} \\
\hline
booking & réservation & \eng{confirm a booking} \\
\hline
application & demande, candidature & \eng{an application form} \\
\hline
fee & frais & \eng{pay a fee} \\
\hline
receipt & reçu & \eng{keep the receipt} \\
\hline
customer & client & \eng{a regular customer} \\
\hline
receptionist & réceptionniste & \eng{ask the receptionist} \\
\hline
embassy & ambassade & \eng{go to the embassy} \\
\hline
travel agency & agence de voyages & \eng{a reliable travel agency} \\
\hline
\end{tabularx}
\end{center}
\end{aretenir}

\courssub{Exercice résolu et piège à éviter}

\begin{exoresolu}[Complète avec le mot juste]
Énoncé. Complète chaque phrase avec le mot correct choisi parmi : \eng{request, booked, applied, receipt, appointment, assistance}.

1. She \eng{................} for a scholarship last year.

2. I need to make an \eng{................} with the doctor.

3. We \eng{................} two seats for the concert in Douala.

4. Please keep your \eng{................} after paying the fee.

5. He made a formal \eng{................} for a new identity card.

6. Do you need any \eng{................} with your luggage?

\tcblower
Solution détaillée.

1. \eng{applied} : construction \eng{to apply for something} (solliciter une bourse). « She applied for a scholarship last year. » = Elle a sollicité une bourse l'année dernière.

2. \eng{appointment} : \eng{to make an appointment with somebody} (prendre rendez-vous avec quelqu'un).

3. \eng{booked} : \eng{to book seats} (réserver des places) ; passé régulier en \eng{-ed}.

4. \eng{receipt} : on garde un reçu après un paiement ; attention au \eng{p} muet.

5. \eng{request} : \eng{to make a request for something} (faire une demande formelle).

6. \eng{assistance} : nom indénombrable, précédé de \eng{any}, jamais de \eng{an}.
\end{exoresolu}

\begin{attention}[Le piège majeur : \eng{demand} est un faux ami]
\eng{Demander un service} ne se traduit \textbf{jamais} par \eng{to demand a service} ! En anglais, \eng{to demand} signifie \cle{exiger}, avec une idée de force ou de colère : c'est beaucoup trop fort et impoli.

\eng{The angry customer demanded a refund.} = Le client furieux a exigé un remboursement.

Pour « demander un service » poliment, on dit : \eng{to ask for a service} ou \eng{to request assistance}.

\eng{I would like to request assistance with my account.} = Je voudrais demander de l'aide pour mon compte.
\end{attention}

\begin{savaistu}[Customer service ou customer care ?]
Au Royaume-Uni et aux États-Unis, on dit \eng{customer service} ou \eng{customer care} pour désigner le service client. Au Nigeria et au Cameroun anglophone, on entend très souvent l'expression \eng{customer care officer} pour désigner l'agent qui accueille et renseigne les clients dans les banques, les compagnies de téléphone et les administrations. À Bamenda ou à Buea, c'est le premier mot que tu entendras en entrant dans une agence !
\end{savaistu}

\begin{exercice}[À toi de jouer]
1. Traduis en anglais : (a) Je voudrais réserver une chambre pour trois nuits. (b) Elle a fait une demande de visa à l'ambassade. (c) Le client s'est plaint des frais.

2. Choisis le bon mot : \eng{request} ou \eng{question} ? (a) I have a \eng{................} : where is the post office? (b) She made a \eng{................} for a new bank card.

3. Trouve l'erreur et corrige-la : \eng{He demanded politely a receipt at the hotel.}
\end{exercice}

\begin{corrige}[Exercice « À toi de jouer »]
1. (a) \eng{I would like to book a room for three nights.} (b) \eng{She made a request for a visa at the embassy.} (c) \eng{The customer complained about the fees.}

2. (a) \eng{question} (on veut une information) ; (b) \eng{request} (on veut obtenir quelque chose).

3. \eng{He asked politely for a receipt at the hotel.} — \eng{To demand} est trop fort ; pour une demande polie, on utilise \eng{to ask for} ou \eng{to request}.
\end{corrige}

\coursec{Champ lexical 2 : Formules fonctionnelles de la requête}

Dans la section précédente, nous avons découvert le vocabulaire général de la demande de service : \eng{to request}, \eng{to ask for}, \eng{assistance}, \eng{a favour}, \eng{a service}. Mais connaître les mots ne suffit pas : en anglais, la \cle{politesse} passe par des \cle{formules} précises, codifiées, qu'il faut employer exactement. Dire \eng{“Give me a room!”} à la réception d'un hôtel de Douala serait perçu comme extrêmement impoli, alors que la traduction française « Donnez-moi une chambre » reste acceptable. Cette section te donne toutes les \cle{formules fonctionnelles} pour demander, répondre, refuser et remercier — à l'oral comme à l'écrit.

\courssub{Observer avant de mémoriser : un dialogue à l'hôtel}

\begin{textebox}[At the hotel — a mini-dialogue]
— Good morning! I would like to book a single room, please.\\
— Certainly, sir. For how many nights?\\
— Two nights, please.\\
— Could you fill in this form?\\
— Of course. Thank you very much.\\
— You're welcome, sir. Enjoy your stay!
\end{textebox}

\textbf{Glossaire}

\begin{center}
\begin{tabularx}{\linewidth}{|X|X|X|}\hline
\rowcolor{popblueL} English & Nature & Français \\ \hline
to book & verbe & réserver \\ \hline
a single room & nom & une chambre simple (une personne) \\ \hline
to fill in (a form) & verbe à particule & remplir (un formulaire) \\ \hline
a stay & nom & un séjour \\ \hline
\end{tabularx}
\end{center}

Observe bien ce dialogue : aucune phrase n'est un ordre direct. Le client dit \eng{“I would like to...”} et non \eng{“I want...”} ; le réceptionniste dit \eng{“Could you fill in...?”} et non \eng{“Fill in...!”}. C'est toute la logique de la politesse anglaise.

\courssub{Demander poliment un service}

\begin{propriete}[La règle d'or de la politesse anglaise]
Plus la phrase est \cle{longue} et \cle{indirecte}, plus elle est \cle{polie} en anglais. L'impératif direct est réservé aux urgences, aux ordres ou aux amis très proches.
\begin{itemize}
\item \eng{“Help me!”} — direct, familier ou urgent.
\item \eng{“Can you help me?”} — neutre, courant entre connaissances.
\item \eng{“Could you help me, please?”} — poli, usage général.
\item \eng{“Could you possibly help me?”} — très poli.
\item \eng{“I was wondering if you could help me.”} — extrêmement poli et formel.
\end{itemize}
\end{propriete}

\begin{popfigure}
\begin{tikzpicture}[scale=1]
\draw[-{Stealth[length=3mm]}, very thick, popblue] (0,0) -- (12,0);
\node[below, popdark] at (0.2,-0.1) {\small impoli};
\node[below, popdark] at (11.8,-0.1) {\small très poli};
\foreach \x/\txt in {0.6/{“Give me\\ a pen.”}, 3.2/{“Can you\\ help me?”}, 5.8/{“Could you\\ help me,\\ please?”}, 8.6/{“Would you mind\\ helping me?”}, 11.2/{“I was wondering if\\ you could help me.”}} {
  \fill[poporange] (\x,0) circle (2.5pt);
  \node[above, align=center, popdark] at (\x,0.15) {\footnotesize \txt};
}
\end{tikzpicture}
\legende{L'échelle de la politesse : plus on monte, plus la formule est longue et indirecte.}
\end{popfigure}

\begin{definition}[Les cinq formules de base pour demander un service]
\begin{itemize}
\item \eng{Could you (please) + base verbale...?} — « Pourriez-vous... ? » : \eng{“Could you please open the door?”} « Pourriez-vous ouvrir la porte, s'il vous plaît ? »
\item \eng{Would you mind + verbe en -ING...?} — « Cela vous dérangerait-il de... ? » : \eng{“Would you mind filling in this form?”} « Cela vous dérangerait-il de remplir ce formulaire ? »
\item \eng{I would like to + base verbale / I would like + nom} — « Je voudrais... » : \eng{“I would like to open an account.”} « Je voudrais ouvrir un compte. »
\item \eng{May I + base verbale...?} — « Puis-je... ? » (formel) : \eng{“May I use your phone?”} « Puis-je utiliser votre téléphone ? »
\item \eng{Can I have + nom...?} — « Est-ce que je peux avoir... ? » : \eng{“Can I have the bill, please?”} « L'addition, s'il vous plaît ? »
\end{itemize}
\end{definition}

\begin{attention}[Après « Would you mind », le verbe est en -ING]
C'est l'erreur classique du francophone. On dit \eng{“Would you mind opening the window?”} et jamais \eng{“Would you mind to open the window?”}. De même, pour répondre, attention à la logique : \eng{“Would you mind waiting?”} — \eng{“No, not at all.”} signifie « Non, cela ne me dérange pas » (donc j'accepte), tandis que \eng{“Yes, I would.”} signifie un refus. Le francophone se trompe souvent dans la réponse !
\end{attention}

\begin{exemplebox}[Exemples traduits]
\begin{itemize}
\item \eng{“Would you mind filling in this form?”} « Cela vous dérangerait-il de remplir ce formulaire ? »
\item \eng{“Could you possibly lend me your dictionary?”} « Pourrais-tu éventuellement me prêter ton dictionnaire ? »
\item \eng{“I would like to see the manager, please.”} « Je voudrais voir le directeur, s'il vous plaît. »
\item \eng{“May I ask you a question?”} « Puis-je vous poser une question ? »
\item \eng{“Can I have two tickets for Bamenda, please?”} « Deux billets pour Bamenda, s'il vous plaît. »
\end{itemize}
\end{exemplebox}

\courssub{Demander un renseignement}

\begin{definition}[Formules pour demander une information]
\begin{itemize}
\item \eng{Could you tell me + question indirecte...?} — « Pourriez-vous me dire... ? » : \eng{“Could you tell me where the bank is?”} « Pourriez-vous me dire où se trouve la banque ? »
\item \eng{I'd like some information about...} — « J'aimerais des renseignements sur... » : \eng{“I'd like some information about your prices.”} « J'aimerais des renseignements sur vos tarifs. »
\item \eng{How much does it cost?} — « Combien cela coûte-t-il ? » ; variante : \eng{“How much is it?”} « C'est combien ? »
\end{itemize}
\end{definition}

\begin{attention}[L'ordre des mots dans la question indirecte]
Après \eng{Could you tell me}, il n'y a \textbf{plus d'inversion} : on revient à l'ordre sujet + verbe.
\begin{itemize}
\item Question directe : \eng{“Where is the post office?”}
\item Question indirecte : \eng{“Could you tell me where the post office \textbf{is}?”} (et non \eng{“where is the post office”})
\end{itemize}
De même : \eng{“Could you tell me how much it costs?”} et non \eng{“how much does it cost”}.
\end{attention}

\courssub{Répondre : accepter, refuser, remercier}

\begin{propriete}[Réponses positives et refus polis]
\textbf{Accepter :} \eng{“Certainly.”} « Bien sûr. » / \eng{“Of course.”} « Naturellement. » / \eng{“With pleasure.”} « Avec plaisir. » / \eng{“Sure, no problem.”} « Bien sûr, pas de problème. » (familier)

\textbf{Refuser poliment :} on ne dit jamais un simple \eng{“No”}. On adoucit toujours :
\begin{itemize}
\item \eng{“I'm afraid not.”} « Je crains que non. »
\item \eng{“I'm sorry, but...”} « Je suis désolé, mais... »
\item \eng{“Unfortunately, we can't...”} « Malheureusement, nous ne pouvons pas... »
\end{itemize}

\textbf{Remercier et conclure :} \eng{“Thank you very much for your help.”} « Merci beaucoup pour votre aide. » / \eng{“I really appreciate it.”} « J'apprécie vraiment. » / \eng{“Thanks a lot.”} « Merci beaucoup. » (familier)
\end{propriete}

\begin{exemplebox}[Exemples traduits]
\begin{itemize}
\item \eng{“I'm afraid the manager is not available.”} « Je crains que le directeur ne soit pas disponible. »
\item \eng{“I'm sorry, but the service is closed on Sundays.”} « Je suis désolé, mais le service est fermé le dimanche. »
\item \eng{“Unfortunately, we can't repair your phone today.”} « Malheureusement, nous ne pouvons pas réparer votre téléphone aujourd'hui. »
\item \eng{“— Could you help me? — With pleasure!”} « — Pourriez-vous m'aider ? — Avec plaisir ! »
\end{itemize}
\end{exemplebox}

\courssub{Registre de langue : formel ou informel ?}

\begin{center}
\begin{tabularx}{\linewidth}{|X|X|X|}\hline
\rowcolor{popblueL} Situation & Formel (banque, administration) & Informel (amis, famille) \\ \hline
Demander & \eng{I would like to...} / \eng{Could you possibly...?} & \eng{Can you...?} / \eng{Can I have...?} \\ \hline
Accepter & \eng{Certainly.} / \eng{Of course.} & \eng{Sure!} / \eng{No problem.} \\ \hline
Refuser & \eng{I'm afraid not.} / \eng{Unfortunately,...} & \eng{Sorry, I can't.} \\ \hline
Remercier & \eng{Thank you very much for your help.} & \eng{Thanks a lot!} / \eng{Cheers!} \\ \hline
\end{tabularx}
\end{center}

\begin{methode}[Formuler une requête en 3 étapes]
\begin{enumerate}
\item \textbf{Attirer l'attention} : \eng{“Excuse me,”} / \eng{“Good morning,”} — « Excusez-moi, » / « Bonjour, »
\item \textbf{Formuler la demande} : \eng{“Could you...?”} / \eng{“I would like to...”} — « Pourriez-vous... ? » / « Je voudrais... »
\item \textbf{Remercier} : \eng{“Thank you so much.”} / \eng{“I really appreciate it.”} — « Merci beaucoup. » / « J'apprécie vraiment. »
\end{enumerate}
Exemple complet : \eng{“Excuse me, could you tell me where the nearest pharmacy is, please? Thank you so much!”} « Excusez-moi, pourriez-vous me dire où se trouve la pharmacie la plus proche, s'il vous plaît ? Merci beaucoup ! »
\end{methode}

\begin{exoresolu}[Reformuler pour être poli]
Énoncé : reformule ces demandes trop directes avec la formule indiquée. 1) \eng{“Open the window!”} (\eng{Would you mind}) 2) \eng{“I want a room.”} (\eng{I would like}) 3) \eng{“Where is the bank?”} (\eng{Could you tell me})
\tcblower
Solution détaillée :
\begin{enumerate}
\item \eng{“Would you mind opening the window?”} — verbe en \textbf{-ING} après \eng{Would you mind}.
\item \eng{“I would like a room, please.”} — \eng{I would like} + nom, plus poli que \eng{I want}.
\item \eng{“Could you tell me where the bank is?”} — question indirecte : sujet (\eng{the bank}) avant le verbe (\eng{is}), sans inversion.
\end{enumerate}
\end{exoresolu}

\begin{exercice}[À toi de jouer]
Mets ces répliques dans l'ordre pour reconstituer un dialogue à la banque : a) \eng{“Certainly, madam. Could you fill in this form?”} b) \eng{“Good morning! I would like to change some money, please.”} c) \eng{“Of course. Thank you very much for your help.”} d) \eng{“Good morning, madam. How can I help you?”}
\end{exercice}

\begin{corrige}[Exercice — À toi de jouer]
Ordre correct : d, b, a, c. L'employé accueille le client (\eng{“How can I help you?”}), le client formule sa demande (\eng{“I would like to...”}), l'employé accepte et demande de remplir un formulaire, le client accepte et remercie.
\end{corrige}

\begin{aretenir}[L'essentiel de la section]
\begin{itemize}
\item En anglais, la politesse = longueur + indirect : \eng{Could you...?} $>$ \eng{Can you...?} $>$ impératif.
\item \eng{Would you mind} + verbe en \textbf{-ING}.
\item Question indirecte après \eng{Could you tell me} : pas d'inversion.
\item On refuse toujours avec une formule adoucie : \eng{“I'm afraid not.”}, \eng{“I'm sorry, but...”}
\item Structure d'une requête : attirer l'attention $\rightarrow$ demander $\rightarrow$ remercier.
\end{itemize}
\end{aretenir}

\coursec{Collocations et expressions idiomatiques}

Dans la section précédente, nous avons appris à formuler une requête polie grâce aux formules fonctionnelles (\eng{Could you...?}, \eng{I would like to...}, \eng{Would you mind...?}). Mais une formule polie ne suffit pas : il faut encore la compléter avec les bons mots. Or, en anglais, les verbes et les noms ne s'associent pas librement : on ne peut pas employer n'importe quel verbe avec n'importe quel nom. C'est précisément l'objet de cette section : les \cle{collocations}, c'est-à-dire les associations habituelles et presque obligatoires entre mots, ainsi que les \cle{expressions idiomatiques} liées au monde des services.

\courssub{Observation : découvrir les collocations dans un texte}

Lisons d'abord ce court texte original. Souligne mentalement tous les groupes « verbe + nom » qui parlent de services : ce sont eux que nous étudierons ensuite.

\begin{textebox}[A busy morning at the Reunification Hotel, Douala]
Mrs Etoundi arrived at the Reunification Hotel in Douala at eight o'clock in the morning. She had \textbf{made a reservation} online two weeks earlier, but when she reached the reception desk, the clerk could not find her name. “I am terribly sorry, madam,” he said. “Could you \textbf{fill in this form} again, please? Our computer system is \textbf{out of service} this morning.”

Mrs Etoundi did not \textbf{make a complaint}. Instead, she smiled and asked: “Could you \textbf{do me a favour}? I have \textbf{made an appointment} with a business partner at ten o'clock, and I really need my room before then.”

The receptionist promised to \textbf{go the extra mile}. He \textbf{booked a room} for her immediately, upgraded her free of charge, and offered her breakfast. “We always try to \textbf{provide the best service} possible,” he explained. “\textbf{Customer satisfaction} is our priority.”

When she reached her room, Mrs Etoundi discovered that she could \textbf{order room service} at any time of the day. She paid the deposit \textbf{by card}, thanked the receptionist warmly, and said: “You have really \textbf{been of service} to me this morning.” — “\textbf{At your service}, madam!” he replied with a bow.
\end{textebox}

\textbf{Glossaire :}

\begin{center}
\begin{tabularx}{\linewidth}{|l|l|X|}
\hline
\rowcolor{popblueL} Mot anglais & Nature & Traduction française \\
\hline
clerk & nom & employé(e) de réception (prononcé « clârk » en anglais britannique) \\
\hline
out of service & locution & hors service, en panne \\
\hline
to upgrade & verbe & surclasser, donner une meilleure chambre \\
\hline
free of charge & locution & gratuitement \\
\hline
deposit & nom & acompte, caution \\
\hline
a bow & nom & une révérence, une inclination (prononcé « bao ») \\
\hline
\end{tabularx}
\end{center}

\textbf{Questions de compréhension :}
\begin{enumerate}
\item Why could the clerk not find Mrs Etoundi's name?
\item What favour did Mrs Etoundi ask for?
\item How did the receptionist “go the extra mile”?
\item How did Mrs Etoundi pay the deposit?
\end{enumerate}

\textbf{Réponses attendues :} 1. Because the computer system was out of service. 2. She asked to have her room before ten o'clock. 3. He booked a room immediately, upgraded her free of charge and offered her breakfast. 4. She paid by card.

Ce texte contient presque toutes les collocations de cette leçon. Observons-les maintenant de près, verbe par verbe.

\courssub{Collocations avec \eng{make} et \eng{do}}

\begin{definition}[Qu'est-ce qu'une collocation ?]
Une \cle{collocation} est une association habituelle, quasi fixe, de deux mots ou plus (souvent verbe + nom) qui « sonnent juste » ensemble pour un locuteur natif. On dit \eng{make a request} et non \eng{do a request} ; on dit \eng{do a favour} et non \eng{make a favour}. Il n'y a souvent aucune logique : il faut les apprendre par cœur, comme des blocs indivisibles.
\end{definition}

\begin{propriete}[Les collocations essentielles avec \eng{make} et \eng{do}]
\begin{itemize}
\item \eng{make a request} = faire une demande. \eng{She made a request for a quieter room.} « Elle a fait une demande pour une chambre plus calme. »
\item \eng{make an appointment} = prendre rendez-vous. \eng{I made an appointment with the dentist.} « J'ai pris rendez-vous chez le dentiste. »
\item \eng{make a reservation / a booking} = faire une réservation. \eng{We made a reservation at a restaurant in Bonanjo.} « Nous avons fait une réservation dans un restaurant à Bonanjo. »
\item \eng{make a complaint} = porter plainte, se plaindre (officiellement). \eng{The customer made a complaint about the dirty sheets.} « Le client s'est plaint des draps sales. »
\item \eng{do someone a favour} = rendre service à quelqu'un. \eng{Could you do me a favour?} « Pourrais-tu me rendre service ? »
\end{itemize}
\textbf{Règle générale (approximative) :} \eng{make} s'emploie pour ce qu'on \emph{crée, produit} (une demande, une réservation, une plainte) ; \eng{do} s'emploie pour une \emph{action, un service, un travail} (\eng{do a favour}, \eng{do business}, \eng{do one's job}).
\end{propriete}

\courssub{Collocations avec \eng{book}, \eng{fill}, \eng{pay}}

\begin{exemplebox}[Les verbes de la démarche administrative]
\begin{itemize}
\item \eng{book a room / book a ticket} = réserver une chambre / un billet. \eng{I booked a ticket to Buea online.} « J'ai réservé un billet pour Buea en ligne. »
\item \eng{fill in a form} (GB) / \eng{fill out a form} (US) = remplir un formulaire. \eng{Please fill in this registration form.} « Veuillez remplir ce formulaire d'inscription. »
\item \eng{pay a fee} = payer des frais (inscription, service). \eng{Students pay a small fee for the school canteen.} « Les élèves paient de petits frais pour la cantine scolaire. »
\item \eng{pay in cash / pay by card / pay by mobile money} = payer en espèces / par carte / par mobile money. \eng{Can I pay by card, or only in cash?} « Puis-je payer par carte, ou seulement en espèces ? »
\end{itemize}
\end{exemplebox}

\begin{attention}[\eng{fill in} et non \eng{fill}]
En anglais britannique, on dit \eng{to fill \textbf{in} a form} ; la particule \eng{in} est obligatoire. Dire \eng{fill a form} est une erreur typique des francophones (calque de « remplir un formulaire »). À l'américain, on dit \eng{to fill \textbf{out} a form}. Attention aussi : \eng{to book} signifie « réserver », pas « réserver un livre » ; un « livre » se dit \eng{a book}, mais le verbe n'a rien à voir !
\end{attention}

\courssub{Collocations avec \eng{provide}, \eng{offer}, \eng{render}}

\begin{propriete}[Parler du côté du prestataire]
\begin{itemize}
\item \eng{provide a service} = fournir un service (le verbe le plus courant). \eng{The bank provides a 24-hour service.} « La banque fournit un service 24 heures sur 24. » Structures : \eng{provide something \textbf{for} somebody} ou \eng{provide somebody \textbf{with} something}. \eng{The hotel provides guests with free Wi-Fi.} « L'hôtel fournit le Wi-Fi gratuit aux clients. »
\item \eng{offer assistance / offer help} = offrir de l'aide. \eng{The staff offered assistance to the elderly passengers.} « Le personnel a offert son aide aux passagers âgés. »
\item \eng{render a service} = rendre un service (registre \emph{formel}, administratif). \eng{The services rendered by the hospital are free of charge.} « Les services rendus par l'hôpital sont gratuits. »
\end{itemize}
\end{propriete}

\courssub{Expressions idiomatiques autour de \eng{service}}

\begin{exemplebox}[Quatre expressions à connaître par cœur]
\begin{itemize}
\item \eng{at your service} = à votre service. \eng{“At your service, madam!” said the porter.} « “À votre service, madame !” dit le porteur. »
\item \eng{to be of service (to somebody)} = être utile, rendre service. \eng{Can I be of any service to you?} « Puis-je vous être utile ? »
\item \eng{out of service} = hors service, en panne. \eng{The lift is out of service.} « L'ascenseur est hors service. » (Contraire : \eng{in service}, en fonctionnement.)
\item \eng{room service} = service en chambre (hôtel). \eng{Guests can order room service until midnight.} « Les clients peuvent commander le service en chambre jusqu'à minuit. »
\end{itemize}
\end{exemplebox}

\begin{savaistu}[“At your service!”]
\eng{At your service!} est une formule cérémonieuse et légèrement théâtrale qu'on entend dans les hôtels de luxe britanniques, dans les films historiques et chez les majordomes (\eng{butlers}) du cinéma anglais. Elle signifie littéralement « je suis à votre disposition ». Au Cameroun anglophone, on l'entend parfois sur le ton de la plaisanterie entre amis : \eng{“Need a lift to the market? At your service!”} « Besoin d'être déposé au marché ? À ton service ! »
\end{savaistu}

\courssub{Expressions utiles pour le Bac}

Deux expressions reviennent souvent dans les sujets d'expression écrite sur la vie économique :

\begin{definition}[Vocabulaire de la qualité du service]
\begin{itemize}
\item \eng{to go the extra mile} = se surpasser, faire plus que ce qui est demandé (pour un client). \eng{A good hotel always goes the extra mile for its guests.} « Un bon hôtel se surpasse toujours pour ses clients. »
\item \eng{customer satisfaction} = la satisfaction du client. \eng{Customer satisfaction should be the priority of every business.} « La satisfaction du client devrait être la priorité de toute entreprise. »
\end{itemize}
Ces deux expressions, bien placées dans une composition, montrent au correcteur une maîtrise du lexique économique authentique.
\end{definition}

\courssub{Tableau récapitulatif : verbe + nom}

\begin{center}
\begin{tabularx}{\linewidth}{|l|l|X|}
\hline
\rowcolor{popblueL} Collocation & Traduction & Exemple en contexte \\
\hline
make a request & faire une demande & \eng{He made a request for a receipt.} \\
\hline
make an appointment & prendre rendez-vous & \eng{She made an appointment at the embassy.} \\
\hline
make a reservation & faire une réservation & \eng{We made a reservation for two nights.} \\
\hline
make a complaint & porter plainte & \eng{They made a complaint about the noise.} \\
\hline
do somebody a favour & rendre service à qqn & \eng{Do me a favour and call the garage.} \\
\hline
book a room / a ticket & réserver une chambre / un billet & \eng{I booked a room near the stadium.} \\
\hline
fill in a form & remplir un formulaire & \eng{Fill in the form in capital letters.} \\
\hline
pay a fee & payer des frais & \eng{You must pay a registration fee.} \\
\hline
pay in cash / by card & payer en espèces / par carte & \eng{He paid the bill in cash.} \\
\hline
provide a service & fournir un service & \eng{The agency provides a delivery service.} \\
\hline
offer assistance & offrir de l'aide & \eng{A policeman offered assistance.} \\
\hline
render a service (formel) & rendre un service & \eng{Thank you for services rendered.} \\
\hline
\end{tabularx}
\end{center}

\begin{popfigure}
\begin{tikzpicture}[font=\small]
\node[draw=popblue, fill=popblueL, rounded corners, minimum width=3cm, minimum height=0.9cm] (v1) at (0,4) {\textbf{make}};
\node[draw=poporange, fill=poporangeL, rounded corners, minimum width=5cm, minimum height=0.7cm] (n1a) at (5.5,4.8) {a request};
\node[draw=poporange, fill=poporangeL, rounded corners, minimum width=5cm, minimum height=0.7cm] (n1b) at (5.5,4.0) {an appointment};
\node[draw=poporange, fill=poporangeL, rounded corners, minimum width=5cm, minimum height=0.7cm] (n1c) at (5.5,3.2) {a reservation / a complaint};
\draw[-{Stealth}, thick, popblue] (v1) -- (n1a);
\draw[-{Stealth}, thick, popblue] (v1) -- (n1b);
\draw[-{Stealth}, thick, popblue] (v1) -- (n1c);
\node[draw=popgreen, fill=popgreenL, rounded corners, minimum width=3cm, minimum height=0.9cm] (v2) at (0,2) {\textbf{book / fill in / pay}};
\node[draw=poporange, fill=poporangeL, rounded corners, minimum width=5cm, minimum height=0.7cm] (n2a) at (5.5,2.6) {a room / a ticket};
\node[draw=poporange, fill=poporangeL, rounded corners, minimum width=5cm, minimum height=0.7cm] (n2b) at (5.5,1.8) {a form};
\node[draw=poporange, fill=poporangeL, rounded corners, minimum width=5cm, minimum height=0.7cm] (n2c) at (5.5,1.0) {a fee / in cash / by card};
\draw[-{Stealth}, thick, popgreen] (v2) -- (n2a);
\draw[-{Stealth}, thick, popgreen] (v2) -- (n2b);
\draw[-{Stealth}, thick, popgreen] (v2) -- (n2c);
\node[draw=poppurple, fill=poppurpleL, rounded corners, minimum width=3cm, minimum height=0.9cm] (v3) at (0,0) {\textbf{provide / offer / do}};
\node[draw=poporange, fill=poporangeL, rounded corners, minimum width=5cm, minimum height=0.7cm] (n3a) at (5.5,0.4) {a service / assistance};
\node[draw=poporange, fill=poporangeL, rounded corners, minimum width=5cm, minimum height=0.7cm] (n3b) at (5.5,-0.4) {somebody a favour};
\draw[-{Stealth}, thick, poppurple] (v3) -- (n3a);
\draw[-{Stealth}, thick, poppurple] (v3) -- (n3b);
\end{tikzpicture}
\legende{Carte des collocations : à gauche le verbe, à droite le groupe nominal qui l'accompagne correctement.}
\end{popfigure}

\courssub{Exercice résolu et pièges à éviter}

\begin{exoresolu}[Choisis la bonne collocation]
Énoncé — Complète chaque phrase avec le bon verbe : \eng{make}, \eng{do}, \eng{book}, \eng{fill in}, \eng{pay}, \eng{provide}.
\begin{enumerate}
\item Could you \dots\ me a favour and wake me up at six?
\item I would like to \dots\ a table for four people, please.
\item All visitors must \dots\ this form before entering.
\item The guest decided to \dots\ a complaint about the cold water.
\item Do you \dots\ room service in this hotel?
\item You can \dots\ the fee in cash or by card.
\end{enumerate}
\tcblower
Solution détaillée :
\begin{enumerate}
\item \eng{do} — on dit toujours \eng{do somebody a favour} (jamais \eng{make a favour}).
\item \eng{book} — \eng{book a table / a room / a ticket} = réserver.
\item \eng{fill in} — \eng{fill in a form} (GB) ; la particule \eng{in} est obligatoire.
\item \eng{make} — \eng{make a complaint} = porter plainte (on « crée » une plainte).
\item \eng{provide} — \eng{provide a service / room service} = fournir un service.
\item \eng{pay} — \eng{pay a fee in cash or by card}.
\end{enumerate}
\end{exoresolu}

\begin{attention}[Les erreurs typiques des francophones]
\begin{itemize}
\item Ne traduis jamais « rendre un service » mot à mot par \eng{to give back a service} ! Cette expression signifierait « rendre un objet emprunté ». On dit \eng{to do someone a favour} (courant) ou \eng{to render a service} (formel).
\item « Prendre rendez-vous » ne se traduit pas par \eng{take an appointment} (calque du français) mais par \eng{\textbf{make} an appointment}.
\item \eng{to assist} est un faux ami partiel : il signifie « aider » (\eng{assist a customer}), alors que « assister à » (une réunion) se dit \eng{to \textbf{attend}}. \eng{I attended the meeting.} « J'ai assisté à la réunion. »
\item Attention aux prépositions : \eng{pay \textbf{in} cash} mais \eng{pay \textbf{by} card} ; \eng{ask \textbf{for} a service}, \eng{thank somebody \textbf{for} their help}.
\end{itemize}
\end{attention}

\begin{exercice}[À toi de jouer]
Traduis en anglais en utilisant les collocations de la leçon :
\begin{enumerate}
\item Pourrais-tu me rendre service et réserver une chambre à Bamenda ?
\item L'ascenseur est hors service ; veuillez remplir ce formulaire à la réception.
\item Cette entreprise se surpasse toujours pour la satisfaction de ses clients.
\end{enumerate}
\end{exercice}

\begin{corrige}[Exercice : traduction]
\begin{enumerate}
\item \eng{Could you do me a favour and book a room in Bamenda?}
\item \eng{The lift is out of service; please fill in this form at the reception desk.}
\item \eng{This company always goes the extra mile for customer satisfaction.}
\end{enumerate}
\end{corrige}

\begin{aretenir}[L'essentiel de la section]
Une collocation s'apprend comme un bloc indivisible : \eng{make a request / an appointment / a reservation / a complaint}, \eng{do somebody a favour}, \eng{book a room}, \eng{fill in a form}, \eng{pay a fee in cash or by card}, \eng{provide a service}, \eng{offer assistance}, \eng{render a service} (formel). Les expressions idiomatiques \eng{at your service}, \eng{to be of service}, \eng{out of service}, \eng{room service}, \eng{to go the extra mile} et \eng{customer satisfaction} enrichiront immédiatement tes productions orales et écrites. Méfie-toi des calques du français : on ne « prend » pas un rendez-vous, on le « fait » (\eng{make}) ; on ne « rend » pas un service en le « redonnant », on fait une faveur (\eng{do a favour}).
\end{aretenir}

\coursec{Faux amis et pièges des francophones}

Dans la section précédente, nous avons appris les collocations et les expressions idiomatiques de la requête : \eng{to make a request}, \eng{to book a table}, \eng{to place an order}. Mais attention : connaître les bons mots ne suffit pas. Encore faut-il ne pas les confondre avec des mots français qui leur ressemblent ! C'est le piège des \cle{faux amis} (\eng{false friends} ou \eng{false cognates}) : des mots anglais qui ressemblent à des mots français mais n'ont pas le même sens. Cette section est peut-être la plus importante de la leçon pour un élève francophone : c'est ici que l'on élimine les erreurs qui trahissent immédiatement un locuteur non natif, à l'oral comme dans les épreuves écrites de Première et du Baccalauréat.

\courssub{Qu'est-ce qu'un faux ami ?}

\begin{definition}[Faux ami]
Un \cle{faux ami} (\eng{false friend}) est un mot anglais qui ressemble par sa forme à un mot français, mais dont le sens est différent. Exemple : \eng{to demand} ressemble à « demander », mais signifie « exiger ». Les faux amis provoquent des contresens parfois comiques, parfois graves : dire \eng{I demand a room} à un réceptionniste, c'est être agressif sans le vouloir !
\end{definition}

\begin{attention}[Pourquoi les francophones sont-ils vulnérables ?]
L'anglais et le français partagent des milliers de mots d'origine latine (\eng{information}, \eng{reservation}, \eng{demand}, \eng{assist}...). La plupart sont de \emph{vrais} amis (\eng{hospital} = hôpital, \eng{service} = service). Mais certains ont évolué différemment dans les deux langues. Le danger vient précisément de cette ressemblance : le mot « vient tout seul », et l'élève ne se méfie pas. Règle d'or : \textbf{plus un mot ressemble au français, plus il faut le vérifier}.
\end{attention}

\courssub{Demander n'est pas « to demand »}

Observons d'abord ces deux phrases :

\begin{exemplebox}[Observation]
\eng{The customer demanded a refund.} « Le client a exigé un remboursement. »

\eng{The customer asked for a refund.} « Le client a demandé un remboursement. »
\end{exemplebox}

Les deux phrases décrivent la même scène, mais le ton est totalement différent. \eng{To demand} exprime une exigence, presque un ultimatum ; \eng{to ask for} exprime une requête normale et polie.

\begin{propriete}[To demand, to ask (for), to request]
\begin{itemize}
\item \eng{To demand something} = \textbf{exiger} quelque chose (ton fort, autoritaire). Structure : \eng{to demand + nom} ou \eng{to demand to + infinitif}.
\item \eng{To ask for something} = \textbf{demander} quelque chose (sens courant, neutre). Attention à la préposition \eng{for} : on demande \emph{for} quelque chose.
\item \eng{To ask somebody to do something} = demander à quelqu'un de faire quelque chose.
\item \eng{To request something} = demander officiellement, par écrit (registre formel, administratif).
\end{itemize}
\end{propriete}

\begin{exemplebox}[Exemples traduits]
\eng{He demanded an explanation.} « Il a exigé une explication. »

\eng{She asked for information.} « Elle a demandé des renseignements. »

\eng{I asked the receptionist to wake me up at six.} « J'ai demandé au réceptionniste de me réveiller à six heures. »

\eng{Passengers are requested to show their tickets.} « Les passagers sont priés de présenter leurs billets. »
\end{exemplebox}

\begin{attention}[Erreur fréquente]
Ne dites jamais \eng{I demand a glass of water} au restaurant : vous donneriez l'impression de menacer le serveur ! Dites \eng{Could I have a glass of water, please?} ou \eng{I would like some water, please.} Réservez \eng{to demand} aux situations de protestation ou d'autorité : \eng{The workers demanded higher wages.} « Les ouvriers ont exigé des salaires plus élevés. »
\end{attention}

\courssub{Assister n'est pas « to assist »}

\begin{propriete}[To assist / to attend]
\begin{itemize}
\item \eng{To assist somebody} = \textbf{aider} quelqu'un (registre formel ; l'équivalent courant est \eng{to help}).
\item « Assister à » un événement (être présent) = \eng{to attend + nom} : \eng{to attend a meeting}, \eng{to attend a ceremony}, \eng{to attend school} (aller à l'école).
\end{itemize}
Attention : \eng{to attend} se construit \textbf{sans préposition} : on dit \eng{I attended the meeting}, jamais \eng{I attended to the meeting}.
\end{propriete}

\begin{exemplebox}[Exemples traduits]
\eng{The nurse assisted the doctor during the operation.} « L'infirmière a aidé le médecin pendant l'opération. »

\eng{We attended the meeting yesterday.} « Nous avons assisté à la réunion hier. »

\eng{Hundreds of fans attended the match at the Ahmadou Ahidjo Stadium.} « Des centaines de supporters ont assisté au match au stade Ahmadou Ahidjo. »
\end{exemplebox}

\begin{attention}[Le piège en situation de service]
À l'hôtel, \eng{Can I assist you?} signifie « Puis-je vous \emph{aider} ? » — c'est la phrase polie du personnel. Si un client répond \eng{No thanks, I don't want to assist} en croyant dire « je ne veux pas assister », le malentendu est total ! De même, un \eng{shop assistant} est un \emph{vendeur} (quelqu'un qui aide les clients), pas un « assistant » au sens français de collaborateur de direction.
\end{attention}

\courssub{Réservation, renseignement, commande : trois pièges classiques}

\textbf{Une réservation.} Le mot existe en anglais : \eng{a reservation} (surtout aux États-Unis). En Grande-Bretagne, on préfère \eng{a booking}, et le verbe est \eng{to book}. Le piège est double : la prononciation et l'orthographe.

\begin{itemize}
\item Prononciation : \eng{reservation} se prononce « ré-zeu-VEI-cheun » (accent sur la troisième syllabe), et non « ré-ser-va-sion » à la française.
\item Orthographe : pas d'accent aigu en anglais ! Écrire \eng{a réservation} est une erreur.
\end{itemize}

\begin{exemplebox}[Exemples traduits]
\eng{I would like to make a reservation for two nights.} (US) « Je voudrais faire une réservation pour deux nuits. »

\eng{Have you got a booking?} (GB) « Avez-vous une réservation ? »

\eng{We booked a table for seven o'clock.} « Nous avons réservé une table pour sept heures. »
\end{exemplebox}

\textbf{Un renseignement.} Le mot français « renseignement » n'a pas d'équivalent qui lui ressemble en anglais : on dit \eng{information}. Et c'est là que se cache un piège grammatical majeur.

\begin{attention}[« Information » est INCOMPTABLE]
En anglais, \eng{information} est un nom \textbf{incomptable} (\eng{uncountable noun}) : il ne prend jamais de \eng{-s} et jamais d'article indéfini \eng{an}.
\begin{itemize}
\item Correct : \eng{some information} « des renseignements », \eng{a piece of information} « un renseignement », \eng{an item of information}, \eng{useful information}.
\item Interdit : \eng{an information}, \eng{informations}.
\end{itemize}
Même famille de pièges : \eng{news} (invariable : \eng{The news is good.}), \eng{advice} (incomptable : \eng{a piece of advice}), \eng{luggage} (incomptable : \eng{a piece of luggage}).
\end{attention}

\begin{exemplebox}[Exemples traduits]
\eng{Could you give me some information about bus times?} « Pourriez-vous me donner des renseignements sur les horaires des bus ? »

\eng{That is a useful piece of information.} « C'est un renseignement utile. »

\eng{For further information, ask at the reception desk.} « Pour de plus amples renseignements, adressez-vous à la réception. »
\end{exemplebox}

\textbf{Une commande.} Au restaurant, « commander » se dit \eng{to order}, et « une commande » se dit \eng{an order}. Le verbe \eng{to command} existe, mais il signifie « commander » au sens militaire ou hiérarchique (commander une armée, un navire), ou « maîtriser » (\eng{a good command of English} = une bonne maîtrise de l'anglais).

\begin{exemplebox}[Exemples traduits]
\eng{Are you ready to order, sir?} « Êtes-vous prêt à commander, monsieur ? »

\eng{The waiter took our order.} « Le serveur a pris notre commande. »

\eng{The general commands an army.} « Le général commande une armée. »
\end{exemplebox}

\courssub{Les calques : traduire mot à mot du français}

Les faux amis ne sont pas les seuls pièges. Il y a aussi les \cle{calques} : des structures françaises copiées mot à mot en anglais. En voici deux, très fréquentes chez les élèves de Première.

\begin{propriete}[Deux calques à bannir]
\begin{itemize}
\item « Je veux que tu m'aides » ne se traduit pas par \eng{I want that you help me}. L'anglais n'emploie pas de subordonnée avec \eng{that} après \eng{to want}. Structure correcte : \eng{to want + complément + to + infinitif} : \eng{I want you to help me.}
\item « Je suis d'accord » ne se traduit pas par \eng{I am agree}. \eng{Agree} est un \textbf{verbe}, pas un adjectif : \eng{I agree.} Au passé : \eng{I agreed.} À la forme négative : \eng{I don't agree.} ou \eng{I disagree.}
\end{itemize}
\end{propriete}

\begin{exemplebox}[Exemples traduits]
\eng{I want you to help me with my luggage.} « Je veux que tu m'aides avec mes bagages. »

\eng{She wants the mechanic to repair her car.} « Elle veut que le mécanicien répare sa voiture. »

\eng{I agree with you: the service was excellent.} « Je suis d'accord avec toi : le service était excellent. »
\end{exemplebox}

\begin{methode}[Comment éviter les calques]
\begin{enumerate}
\item Avant d'écrire ou de parler, demande-toi : « Est-ce que je traduis mot à mot du français ? »
\item Si oui, arrête-toi et reformule : cherche la structure anglaise authentique (verbe + complément + infinitif, verbe sans auxiliaire \eng{be}, etc.).
\item Vérifie le mot suspect dans le tableau des faux amis ci-dessous.
\item Mémorise les expressions \emph{en bloc} (\eng{I agree}, \eng{I want you to...}, \eng{a piece of information}) plutôt que mot par mot : une expression apprise en entier ne peut pas être calquée.
\item Relis ta phrase en te posant la question : « Un Anglais dirait-il cela ainsi ? »
\end{enumerate}
\end{methode}

\courssub{Tableau récapitulatif des faux amis de la requête}

\begin{center}
\begin{tabularx}{\linewidth}{|X|X|X|}
\hline
\rowcolor{popblueL} Français & Faux ami anglais & Traduction correcte \\
\hline
demander & \eng{to demand} (= exiger) & \eng{to ask (for)} / \eng{to request} \\
\hline
assister à & \eng{to assist} (= aider) & \eng{to attend} \\
\hline
un renseignement & — (pas de mot ressemblant) & \eng{a piece of information} \\
\hline
des renseignements & \eng{informations} (n'existe pas) & \eng{information} (invariable) \\
\hline
commander (au restaurant) & \eng{to command} (= commander une armée) & \eng{to order} \\
\hline
une commande & \eng{a command} (= un ordre militaire) & \eng{an order} \\
\hline
une réservation & \eng{a réservation} (accent français !) & \eng{a booking} (GB) / \eng{a reservation} (US) \\
\hline
je suis d'accord & \eng{I am agree} & \eng{I agree} \\
\hline
je veux que tu viennes & \eng{I want that you come} & \eng{I want you to come} \\
\hline
actuellement & \eng{actually} (= en fait) & \eng{at present} / \eng{currently} \\
\hline
\end{tabularx}
\end{center}

\begin{popfigure}
\begin{tikzpicture}[
  col/.style={font=\bfseries\small, text=white, align=center},
  cell/.style={font=\small, align=left, text width=3.4cm, inner sep=4pt}
]
\node[col, fill=popblue, minimum width=3.8cm] (h1) at (0,0) {Français};
\node[col, fill=poporange, minimum width=3.8cm] (h2) at (4.1,0) {Faux ami};
\node[col, fill=popgreen, minimum width=3.8cm] (h3) at (8.2,0) {Traduction correcte};
\node[cell, fill=popblueL, minimum width=3.8cm] at (0,-0.9) {demander};
\node[cell, fill=poporangeL, minimum width=3.8cm] at (4.1,-0.9) {\eng{to demand} = exiger};
\node[cell, fill=popgreenL, minimum width=3.8cm] at (8.2,-0.9) {\eng{to ask (for)}};
\node[cell, fill=popblueL, minimum width=3.8cm] at (0,-2.0) {assister à};
\node[cell, fill=poporangeL, minimum width=3.8cm] at (4.1,-2.0) {\eng{to assist} = aider};
\node[cell, fill=popgreenL, minimum width=3.8cm] at (8.2,-2.0) {\eng{to attend}};
\node[cell, fill=popblueL, minimum width=3.8cm] at (0,-3.1) {commander (repas)};
\node[cell, fill=poporangeL, minimum width=3.8cm] at (4.1,-3.1) {\eng{to command} = commander (armée)};
\node[cell, fill=popgreenL, minimum width=3.8cm] at (8.2,-3.1) {\eng{to order}};
\node[cell, fill=popblueL, minimum width=3.8cm] at (0,-4.2) {des renseignements};
\node[cell, fill=poporangeL, minimum width=3.8cm] at (4.1,-4.2) {\eng{informations} (inexistant)};
\node[cell, fill=popgreenL, minimum width=3.8cm] at (8.2,-4.2) {\eng{information} (invariable)};
\end{tikzpicture}
\legende{Les quatre faux amis essentiels du vocabulaire de la requête : la colonne orange est le piège, la colonne verte la solution.}
\end{popfigure}

\begin{savaistu}[Pourquoi tant de mots se ressemblent-ils ?]
Après la conquête de l'Angleterre par les Normands en 1066, le français est resté la langue de la cour et de l'administration anglaise pendant près de trois siècles. Des milliers de mots français sont alors entrés dans l'anglais : \eng{demand}, \eng{assist}, \eng{command}, \eng{reservation}, \eng{service}... Mais au fil des siècles, certains ont changé de sens d'un côté ou de l'autre de la Manche. C'est ainsi que \eng{to demand} est devenu plus fort que « demander », et que \eng{actually} (qui vient de « actuel ») signifie aujourd'hui « en fait, en réalité ». L'histoire explique le piège : les deux langues sont des sœurs... qui ne se ressemblent pas toujours !
\end{savaistu}

\courssub{Le piège en contexte : une scène à l'hôtel}

\begin{textebox}[At the hotel reception — read and spot the mistakes]
Yesterday evening, a French-speaking tourist arrived at a small hotel in Buea. Tired after his journey, he said to the receptionist: “I demand a room for tonight.” The receptionist looked surprised and a little offended. The tourist continued: “Also, I need some informations about the town. Can you assist me?” The receptionist answered politely: “Of course, sir. I can help you. But first, do you have a reservation?” The tourist replied: “No, but I am agree to pay in advance.”

Later, at the hotel restaurant, he told the waiter: “I want that you bring me the menu, and I will command the fish.” The waiter smiled and said: “Certainly, sir. I will bring you the menu, and you can order the fish.”

The next morning, the tourist attended breakfast at seven o'clock — or rather, he was present at breakfast. He finally understood that English words which look French can be dangerous traps. Before leaving, he asked the receptionist for a piece of advice: “How can I improve my English?” She laughed and said: “Never translate word for word!”

\emph{Glossaire} : offended (adj.) = offensé, vexé ; in advance (loc. adv.) = à l'avance ; rather (adv.) = plutôt ; trap (n.) = piège ; to improve (v.) = améliorer ; word for word (loc.) = mot à mot.
\end{textebox}

\textbf{Questions de compréhension}
\begin{enumerate}
\item Why was the receptionist surprised and offended when the tourist arrived?
\item What mistake did the tourist make with the word \eng{informations}?
\item What is the difference between \eng{to assist} and \eng{to attend} in this text?
\item How should the tourist have said “I want that you bring me the menu”?
\item What advice does the receptionist give at the end?
\end{enumerate}

\begin{exoresolu}[Corrige les erreurs du touriste]
Énoncé. Réécris correctement les cinq phrases fautives du texte : (1) \eng{I demand a room for tonight.} (2) \eng{I need some informations about the town.} (3) \eng{I am agree to pay in advance.} (4) \eng{I want that you bring me the menu.} (5) \eng{I will command the fish.}
\tcblower
Solution détaillée.
\begin{enumerate}
\item \eng{I would like a room for tonight, please.} — \eng{To demand} = exiger ; pour une requête polie, on utilise \eng{I would like...} ou \eng{Could I have...?}
\item \eng{I need some information about the town.} — \eng{Information} est incomptable : jamais de \eng{-s}.
\item \eng{I agree to pay in advance.} — \eng{Agree} est un verbe : pas d'auxiliaire \eng{be}.
\item \eng{I want you to bring me the menu.} — Après \eng{to want}, structure : complément + \eng{to} + infinitif, jamais \eng{that}.
\item \eng{I will order the fish.} — Au restaurant, « commander » = \eng{to order} ; \eng{to command} = commander une armée.
\end{enumerate}
\end{exoresolu}

\begin{exercice}[À toi de jouer]
Traduis en anglais correct, en évitant les faux amis et les calques :
\begin{enumerate}
\item Je voudrais demander un renseignement sur les horaires des cars pour Douala.
\item Nous avons assisté à une cérémonie au palais des congrès.
\item Le client a exigé de parler au directeur.
\item Elle veut que je l'aide à porter ses bagages.
\item Êtes-vous prêts à commander, mesdames ?
\item Je ne suis pas d'accord : le service était très lent.
\end{enumerate}
\end{exercice}

\begin{corrige}[Exercice « À toi de jouer »]
\begin{enumerate}
\item \eng{I would like to ask for some information about coach times to Douala.} (\eng{information} invariable ; \eng{to ask for} = demander.)
\item \eng{We attended a ceremony at the conference centre.} (\eng{to attend} = assister à, sans préposition.)
\item \eng{The customer demanded to speak to the manager.} (\eng{to demand} = exiger ; ici le sens fort est voulu.)
\item \eng{She wants me to help her carry her luggage.} (structure \eng{want + complément + to + infinitif} ; \eng{luggage} incomptable.)
\item \eng{Are you ready to order, ladies?} (\eng{to order} = commander au restaurant.)
\item \eng{I don't agree: the service was very slow.} (pas de \eng{I am agree}.)
\end{enumerate}
\end{corrige}

\begin{aretenir}[L'essentiel de la section]
\begin{itemize}
\item \eng{To demand} = exiger ; « demander » = \eng{to ask (for)} ou \eng{to request}.
\item \eng{To assist} = aider ; « assister à » = \eng{to attend} (sans préposition).
\item \eng{Information} est incomptable : \eng{some information}, \eng{a piece of information}, jamais \eng{informations}.
\item « Commander » au restaurant = \eng{to order} ; \eng{to command} = commander une armée.
\item « Une réservation » = \eng{a booking} (GB) / \eng{a reservation} (US), sans accent aigu.
\item Jamais \eng{I want that you...} ni \eng{I am agree} : dites \eng{I want you to...} et \eng{I agree}.
\item Stratégie : méfie-toi des mots qui ressemblent au français, et apprends les expressions en bloc.
\end{itemize}
\end{aretenir}

\coursec{Familles de mots, prononciation et orthographe}

Dans la section précédente, nous avons débusqué les \cle{faux amis} et les calques du français qui guettent l'élève francophone lorsqu'il demande un service en anglais. Mais connaître un mot isolé ne suffit pas : pour enrichir durablement son vocabulaire, il faut apprendre les mots \emph{par familles}. En effet, à partir d'une seule racine comme \eng{serve}, l'anglais fabrique tout un réseau de mots liés par des \cle{suffixes}. C'est une méthode économique et puissante : retenir une racine, c'est en réalité retenir cinq ou six mots. Dans cette section, nous allons observer comment les suffixes transforment les mots, mémoriser la prononciation de quelques mots pièges, et corriger les fautes d'orthographe les plus fréquentes.

\courssub{Observer avant de mémoriser : un texte pour découvrir les familles de mots}

Lisons d'abord ce court texte original. Soulignez mentalement tous les mots qui se ressemblent : ils appartiennent à des \cle{familles de mots}.

\begin{textebox}[At the microfinance office in Bamenda]
Every morning, Mr Ngwa opens his small microfinance office near the market in Bamenda. He is proud to \textbf{serve} his community, and the \textbf{service} he offers is fast and reliable. His new \textbf{assistant}, Linda, is very \textbf{helpful}: she greets every customer with a smile and helps them fill in the forms. Last week, a young \textbf{applicant} came to \textbf{apply} for a small loan to start a poultry business. His \textbf{application} was complete, so Linda \textbf{booked} an \textbf{appointment} for him with the manager. The \textbf{booking} was confirmed by text message, and the young man received a \textbf{receipt} for the registration fee. “In this office, nobody is \textbf{helpless},” says Mr Ngwa. “When people \textbf{assist} each other, the whole town progresses.” Linda agrees. She says that good \textbf{assistance} at the right moment can change a family's life.
\end{textebox}

\begin{definition}[Glossaire]
\begin{itemize}
\item \eng{loan} (n.) : un prêt, un emprunt.
\item \eng{poultry business} (n.) : une entreprise avicole (élevage de volailles).
\item \eng{to fill in a form} (v.) : remplir un formulaire.
\item \eng{registration fee} (n.) : les frais d'inscription.
\item \eng{reliable} (adj.) : fiable, digne de confiance.
\item \eng{receipt} (n.) : un reçu, une quittance.
\end{itemize}
\end{definition}

\textbf{Questions de compréhension.}
\begin{enumerate}
\item Where is Mr Ngwa's office and what does he do there?
\item Why is Linda described as “helpful”?
\item What did the young applicant want, and what two documents are mentioned in the text?
\end{enumerate}

\textbf{Réponses attendues.} 1. His office is near the market in Bamenda; he runs a microfinance office and serves his community. 2. Because she greets customers with a smile and helps them fill in the forms. 3. He wanted a small loan to start a poultry business; the documents are the application (form) and the receipt.

\courssub{La famille de SERVE}

Observons les phrases suivantes :

\eng{The waiter serves customers quickly.} « Le serveur sert les clients rapidement. »

\eng{The service in this bank is excellent.} « Le service dans cette banque est excellent. »

\eng{In old novels, the servants live in the basement.} « Dans les romans anciens, les domestiques vivent au sous-sol. »

\eng{This old bicycle is still serviceable.} « Ce vieux vélo est encore utilisable. »

Nous voyons que la racine \cle{serve} donne naissance à un verbe, deux noms et un adjectif. Voici l'arbre lexical complet :

\begin{popfigure}
\begin{center}
\begin{tikzpicture}[
grow=down,
level distance=1.6cm,
level 1/.style={sibling distance=3.4cm},
every node/.style={draw=popblue, rounded corners, fill=popblueL, align=center, font=\small},
edge from parent/.style={draw=poporange, thick, -{Stealth}}
]
\node[fill=poporangeL, draw=poporange, font=\bfseries, align=center]{SERVE\\ \emph{racine : servir}}
child { node[align=center]{to serve (v.)\\ servir} }
child { node[align=center]{service (n.)\\ le service} }
child { node[align=center]{servant (n.)\\ le/la domestique} }
child { node[align=center]{serviceable (adj.)\\ utilisable, en état de marche} };
\end{tikzpicture}
\end{center}
\legende{Arbre lexical de la racine SERVE : une racine, quatre mots.}
\end{popfigure}

\begin{attention}[Servant : un mot à manier avec précaution]
\eng{Servant} signifie « domestique » (employé de maison) et a une connotation historique ou littéraire. Pour parler d'un employé de banque, d'un agent ou d'un fonctionnaire, on dira \eng{employee}, \eng{clerk}, \eng{officer} ou \eng{civil servant} (fonctionnaire de l'État). Ne dites jamais \eng{the servant of the bank} pour « l'employé de la banque » ! Dites : \eng{a bank employee} ou \eng{a bank clerk}.
\end{attention}

\courssub{Les familles de HELP, ASSIST, APPLY et BOOK}

\begin{propriete}[Les suffixes qui changent la nature des mots]
En anglais, on fabrique de nouveaux mots en ajoutant un \cle{suffixe} à une racine :
\begin{itemize}
\item \textbf{-ance / -ence} forment des \textbf{noms} abstraits : \eng{assist} $\rightarrow$ \eng{assistance} (l'aide) ; \eng{attend} $\rightarrow$ \eng{attendance} (la présence).
\item \textbf{-ant} forme des noms de \textbf{personnes} : \eng{assist} $\rightarrow$ \eng{assistant} (l'assistant) ; \eng{apply} $\rightarrow$ \eng{applicant} (le candidat).
\item \textbf{-ful} (« plein de ») et \textbf{-less} (« sans ») forment des \textbf{adjectifs contraires} : \eng{helpful} = serviable, qui aide ; \eng{helpless} = impuissant, sans aide.
\item \textbf{-tion} forme des noms d'action : \eng{apply} $\rightarrow$ \eng{application} (la candidature, la demande).
\item \textbf{-ing} forme des noms d'activité : \eng{book} $\rightarrow$ \eng{booking} (la réservation).
\item \textbf{-able} forme des adjectifs de possibilité : \eng{service} $\rightarrow$ \eng{serviceable} (utilisable) ; \eng{book} $\rightarrow$ \eng{bookable} (réservable).
\end{itemize}
\end{propriete}

\begin{center}
\begin{tabularx}{\linewidth}{|l|X|X|X|}
\hline
\rowcolor{popblueL} \textbf{Racine} & \textbf{Verbe} & \textbf{Noms} & \textbf{Adjectifs} \\
\hline
SERVE & to serve (servir) & service (service) ; servant (domestique) & serviceable (utilisable) \\
\hline
HELP & to help (aider) & help (aide) ; helper (auxiliaire) & helpful (serviable) / helpless (impuissant) \\
\hline
ASSIST & to assist (aider) & assistance (aide) ; assistant (assistant) & --- \\
\hline
APPLY & to apply (postuler) & application (candidature) ; applicant (candidat) & applicable (applicable) \\
\hline
BOOK & to book (réserver) & booking (réservation) ; booklet (livret) & bookable (réservable) \\
\hline
\end{tabularx}
\end{center}

\begin{exemplebox}[Exemples traduits]
\eng{The applicant filled in the application form.} « Le candidat a rempli le formulaire de candidature. »

\eng{The receptionist was very helpful.} « La réceptionniste était très serviable. »

\eng{I need assistance with my luggage, please.} « J'ai besoin d'aide avec mes bagages, s'il vous plaît. »

\eng{She works as a teaching assistant at the University of Buea.} « Elle travaille comme assistante d'enseignement à l'université de Buea. »

\eng{I made a booking for two nights at a hotel in Douala.} « J'ai fait une réservation pour deux nuits dans un hôtel à Douala. »

\eng{Without her glasses, my grandmother feels helpless.} « Sans ses lunettes, ma grand-mère se sent impuissante. »

\eng{The hotel gives every guest a booklet with useful information.} « L'hôtel remet à chaque client un livret d'informations utiles. »
\end{exemplebox}

\begin{attention}[Assist : attention au faux ami interne]
\eng{To assist} signifie « aider » (\eng{Can I assist you?} = « Puis-je vous aider ? »). Pour dire « assister à » une réunion ou un match, l'anglais utilise \eng{to attend} : \eng{I attended the meeting.} « J'ai assisté à la réunion. » Ne confondez pas \eng{assistance} (l'aide) et \eng{attendance} (la présence, l'assistance au sens de « public présent ») !
\end{attention}

\courssub{Prononciation : les mots pièges de la leçon}

L'orthographe anglaise ne dit pas toujours comment prononcer le mot. Voici les cinq mots de cette leçon à prononcer correctement, notés en lettres françaises :

\begin{center}
\begin{tabularx}{\linewidth}{|X|X|X|}
\hline
\rowcolor{popblueL} \textbf{Mot} & \textbf{Prononciation} & \textbf{Piège à éviter} \\
\hline
request & « ri-kouèste » & ne pas dire « ré-kouèste » ; le \emph{re-} initial se dit « ri » \\
\hline
assistance & « a-sis-tens » & accent sur la 2\textsuperscript{e} syllabe, pas « a-ssis-tance » à la française \\
\hline
appointment & « a-pointe-mente » & ne pas prononcer « a-po-int-ment » syllabe par syllabe \\
\hline
receipt & « ri-site » & le \textbf{P est muet} ! \\
\hline
queue & « kiou » & se prononce exactement comme la lettre \textbf{Q} \\
\hline
\end{tabularx}
\end{center}

\begin{attention}[Deux prononciations à retenir absolument]
\eng{Receipt} (le reçu) se prononce « ri-site » : le \textbf{P est muet}, comme dans \eng{psychology} (« saï-ko-lo-dji »). Et \eng{queue} (la file d'attente) se prononce exactement comme la lettre \textbf{Q} : « kiou ». Les quatre lettres \emph{-ueue} ne servent qu'à l'écriture ! À l'oral : \eng{There was a long queue at the bank.} « Il y avait une longue file d'attente à la banque » se dit « zère ouoz e long kiou ».
\end{attention}

\begin{savaistu}[Queue ou line ?]
En Grande-Bretagne, faire la queue est presque un sport national : on fait la queue à la banque, à l'arrêt de bus, au supermarché, et sauter la file est très mal vu ! Le verbe britannique est \eng{to queue up} : \eng{We queued up for an hour.} « Nous avons fait la queue pendant une heure. » Aux États-Unis, on dit plutôt \eng{to stand in line} : \eng{We stood in line for an hour.} Au Cameroun anglophone, c'est l'usage britannique qui domine : \eng{queue} et \eng{queue up}.
\end{savaistu}

\courssub{Orthographe : lettres doubles et lettres muettes}

\begin{propriete}[Le redoublement de la consonne finale]
Quand un verbe court se termine par \textbf{une seule voyelle brève + une seule consonne} et que la syllabe est accentuée, la consonne finale \textbf{double} devant \emph{-ing} et \emph{-ed} :
\begin{itemize}
\item \eng{stop} $\rightarrow$ \eng{stopping}, \eng{stopped} ; \eng{plan} $\rightarrow$ \eng{planning}, \eng{planned} ; \eng{put} $\rightarrow$ \eng{putting}.
\item \textbf{Mais} \eng{book} $\rightarrow$ \eng{booking}, \eng{booked} : le \emph{k} ne double \textbf{pas}, car \emph{oo} est un groupe de \textbf{deux voyelles} formant un son long ; la règle ne s'applique qu'après une voyelle simple et brève.
\item \eng{fill} $\rightarrow$ \eng{filling} : le \emph{l} est déjà double dans la racine, rien ne change.
\item \eng{apply} $\rightarrow$ \eng{applying} : le \emph{y} ne change pas devant \emph{-ing} (mais devant \emph{-ed}, il devient \emph{i} : \eng{applied}).
\end{itemize}
Retenez les formes utiles de cette leçon : \eng{booking} (réservation), \eng{filling in a form} (remplir un formulaire), \eng{applying for a job} (postuler à un emploi).
\end{propriete}

\begin{propriete}[Les lettres muettes]
Certaines lettres s'écrivent mais ne se prononcent pas :
\begin{itemize}
\item \eng{receipt} : \textbf{p} muet (« ri-site »).
\item \eng{doubt} : \textbf{b} muet (« daoute ») — comme dans \eng{debt} (« dète », la dette).
\item \eng{queue} : \emph{ueue} muettes (« kiou »).
\item \eng{booklet} : les deux \emph{o} forment un seul son court (« bouk-lete »).
\end{itemize}
\end{propriete}

\begin{exoresolu}[Familles de mots : complétez avec la bonne forme]
Énoncé. Complétez chaque phrase avec la forme correcte du mot entre parenthèses.

1. In the old house, the \ldots\ldots\ldots cleaned the rooms and cooked the meals. (serve — personne)

2. I would like to \ldots\ldots\ldots for the job of receptionist. (apply)

3. Thank you for your \ldots\ldots\ldots ; I could not have finished alone. (assist)

4. The hotel confirmed our \ldots\ldots\ldots by email. (book)

5. This old generator is still \ldots\ldots\ldots . (serve — adjectif)

\tcblower

Solution détaillée.

1. \eng{servant} — on cherche un nom de personne de la famille de \eng{serve} ; le contexte (une vieille maison, les tâches domestiques) appelle \eng{servant} : \eng{In the old house, the servant cleaned the rooms and cooked the meals.} « Dans la vieille maison, le domestique nettoyait les chambres et préparait les repas. »

2. \eng{apply} — après \eng{would like to}, il faut l'infinitif sans modification : \eng{I would like to apply for the job of receptionist.} Attention à la préposition : \eng{to apply \textbf{for} a job} (postuler à un emploi).

3. \eng{assistance} — après le possessif \eng{your}, il faut un nom : \eng{Thank you for your assistance.} Le suffixe \emph{-ance} transforme le verbe en nom abstrait.

4. \eng{booking} — après le possessif \eng{our}, il faut un nom d'activité : \eng{The hotel confirmed our booking by email.} (usage britannique ; en américain : \eng{reservation}).

5. \eng{serviceable} — après \eng{is still}, il faut un adjectif : \eng{This old generator is still serviceable.} « Ce vieux groupe électrogène est encore utilisable. »
\end{exoresolu}

\begin{exercice}[À vous de jouer]
1. Conjuguez la famille : donnez le nom de personne correspondant à \eng{assist} et à \eng{apply}.

2. Traduisez en anglais : « La réceptionniste serviable m'a aidé à remplir le formulaire. »

3. Dites à voix haute, puis écrivez la prononciation en lettres françaises : \eng{receipt}, \eng{queue}, \eng{request}.

4. Vrai ou faux ? \eng{Helpless} signifie « serviable ». Justifiez.

5. Écrivez la forme en \emph{-ing} de : \eng{stop}, \eng{book}, \eng{apply}.
\end{exercice}

\begin{corrige}[Exercice « À vous de jouer »]
1. \eng{assist} $\rightarrow$ \eng{assistant} ; \eng{apply} $\rightarrow$ \eng{applicant}.

2. \eng{The helpful receptionist helped me (to) fill in the form.}

3. \eng{receipt} = « ri-site » (p muet) ; \eng{queue} = « kiou » ; \eng{request} = « ri-kouèste ».

4. Faux. \eng{Helpless} signifie « impuissant, sans aide » (suffixe \emph{-less} = « sans »). « Serviable » se dit \eng{helpful} (suffixe \emph{-ful} = « plein de »).

5. \eng{stop} $\rightarrow$ \eng{stopping} (le \emph{p} double après une voyelle brève) ; \eng{book} $\rightarrow$ \eng{booking} (pas de redoublement après \emph{oo}) ; \eng{apply} $\rightarrow$ \eng{applying} (le \emph{y} est conservé devant \emph{-ing}).
\end{corrige}

\begin{aretenir}[L'essentiel de la section]
\begin{itemize}
\item Une racine + un suffixe = une famille de mots : \eng{serve / service / servant / serviceable}.
\item \emph{-ance} et \emph{-tion} donnent des noms ; \emph{-ant} donne des noms de personnes ; \emph{-ful} et \emph{-less} donnent des adjectifs contraires ; \emph{-able} donne des adjectifs de possibilité.
\item Prononciations pièges : \eng{receipt} « ri-site » (p muet), \eng{queue} « kiou », \eng{request} « ri-kouèste ».
\item Orthographe : la consonne finale double après une voyelle brève accentuée (\eng{stopping}) mais pas après un groupe de voyelles (\eng{booking}) ; lettres muettes dans \eng{receipt} et \eng{doubt}.
\item \eng{To assist} = aider ; « assister à » = \eng{to attend}.
\end{itemize}
\end{aretenir}

\coursec{Le lexique en contexte : texte et dialogues}

Dans la section précédente, nous avons étudié les \cle{familles de mots}, la \cle{prononciation} et l'\cle{orthographe} du vocabulaire de la demande de services : \eng{to request} et \eng{a request}, \eng{to serve} et \eng{service}, \eng{to assist} et \eng{assistance}. Mais un mot n'est vraiment appris que lorsqu'on sait l'\emph{employer dans une situation réelle}. C'est l'objectif de cette dernière section : observer le lexique de la leçon \emph{en action}, dans un texte et des dialogues authentiques, puis le réemployer toi-même à l'écrit et à l'oral. C'est exactement la compétence exigée à l'examen : comprendre un document de la vie quotidienne et réutiliser son vocabulaire dans une production personnelle.

\courssub{Lecture : « At the Buea Post Office »}

\begin{methode}[Comment lire un dialogue de service]
Avant de lire un dialogue ou un texte de service, suis ces quatre étapes :
\begin{enumerate}
\item \textbf{Identifier le lieu et les personnages} : qui parle ? où ? (un client et un employé, à la poste, à l'hôpital, à la banque\ldots).
\item \textbf{Repérer la demande principale} : que veut le client ? C'est le cœur de la scène.
\item \textbf{Noter les formules de politesse} : \eng{Could you\ldots?}, \eng{I would like\ldots}, \eng{please}, \eng{thank you}.
\item \textbf{Relever les mots nouveaux et deviner leur sens grâce au contexte} avant de consulter le glossaire ou le dictionnaire.
\end{enumerate}
\end{methode}

\begin{textebox}[At the Buea Post Office]
Mr Tabi enters the post office in Buea on Monday morning. He is carrying a small brown parcel under his arm. He waits patiently in the queue, then walks to the counter.

— Good morning, madam. I would like to send this parcel to Bamenda, please.

— Certainly, sir. Could you fill in this form and write the receiver's address clearly?

— Of course. Here it is. How much is the fee?

— It depends on the weight. Let me weigh it\ldots\ That will be 2,500 francs.

— Here you are. Could I have a receipt, please?

— Here is your receipt. Your parcel will arrive in Bamenda within three working days.

— That's perfect. And could you also tell me where I can buy stamps? I need to post two letters as well.

— Stamps are sold at counter number four, just on your left.

— Thank you very much for your help. Have a nice day!

— You're welcome, sir. Thank you for using our services. Goodbye!

Mr Tabi smiles, takes his receipt carefully, and walks towards counter number four. He is satisfied: the service was quick, the clerk was polite, and his parcel is now on its way to his sister in Bamenda.
\end{textebox}

\textbf{Glossaire :}

\begin{center}
\begin{tabularx}{\linewidth}{|l|l|X|}\hline
\rowcolor{popblueL} \textbf{Mot anglais} & \textbf{Nature} & \textbf{Traduction française} \\ \hline
\eng{parcel} & nom & colis \\ \hline
\eng{receiver} & nom & destinataire (celui qui reçoit) \\ \hline
\eng{fee} & nom & frais, tarif \\ \hline
\eng{to weigh} & verbe & peser \\ \hline
\eng{receipt} & nom & reçu (prononce « ri-site », le \emph{p} est muet) \\ \hline
\eng{counter} & nom & guichet, comptoir \\ \hline
\eng{queue} & nom & file d'attente (prononce « kiou ») \\ \hline
\eng{clerk} & nom & employé(e) de guichet (prononce « klark » en GB) \\ \hline
\eng{within three working days} & locution & sous trois jours ouvrables \\ \hline
\end{tabularx}
\end{center}

\courssub{Compréhension du texte}

\begin{exercice}[Comprehension check]
Réponds aux questions suivantes sur le texte \emph{At the Buea Post Office}.

\textbf{A. True or False?} (Vrai ou faux ? Justifie ta réponse par une phrase du texte.)
\begin{enumerate}
\item Mr Tabi wants to send a letter to Bamenda.
\item The clerk weighs the parcel before giving the price.
\end{enumerate}

\textbf{B. Choisis la bonne réponse (QCM).}

The fee for the parcel depends on:
\begin{enumerate}
\item[a)] the receiver's address
\item[b)] the weight of the parcel
\item[c)] the number of stamps
\end{enumerate}

\textbf{C. Réponds par une phrase courte et complète.}
\begin{enumerate}
\item Where can Mr Tabi buy stamps?
\item How long will the parcel take to arrive in Bamenda?
\end{enumerate}
\end{exercice}

\begin{corrige}[Comprehension check]
\textbf{A.} 1. \textbf{False} — He wants to send a \emph{parcel}, not a letter: \eng{“I would like to send this parcel to Bamenda.”} 2. \textbf{True} — \eng{“Let me weigh it\ldots\ That will be 2,500 francs.”}

\textbf{B.} La bonne réponse est \textbf{b)} : \eng{“It depends on the weight.”}

\textbf{C.} 1. \eng{He can buy stamps at counter number four.} 2. \eng{It will arrive within three working days.}
\end{corrige}

\courssub{Repérage dans le texte : formules de requête et collocations}

\begin{exoresolu}[Souligne les formules et les collocations]
\textbf{Énoncé.} Dans le texte \emph{At the Buea Post Office}, souligne \textbf{5 formules de requête} et \textbf{3 collocations} (associations de mots typiques).
\tcblower
\textbf{Solution détaillée.}

\textbf{Les 5 formules de requête :}
\begin{enumerate}
\item \eng{I would like to send this parcel to Bamenda, please.} (demande polie avec \eng{I would like} + infinitif)
\item \eng{Could you fill in this form\ldots?} (demande d'action avec \eng{Could you} + base verbale)
\item \eng{Could I have a receipt, please?} (demande d'objet avec \eng{Could I have\ldots?})
\item \eng{Could you also tell me where I can buy stamps?} (demande d'information ; attention à l'ordre des mots : \eng{where I can buy}, et non \emph{where can I buy}, car c'est une question indirecte)
\item \eng{How much is the fee?} (question directe sur le prix, acceptable entre inconnus dans un contexte de service)
\end{enumerate}

\textbf{Les 3 collocations :}
\begin{enumerate}
\item \eng{to send a parcel} = envoyer un colis (verbe + nom)
\item \eng{to fill in a form} = remplir un formulaire (verbe + particule + nom)
\item \eng{working days} = jours ouvrables (adjectif + nom)
\end{enumerate}
On peut ajouter \eng{to weigh a parcel} (peser un colis) et \eng{to buy stamps} (acheter des timbres).
\end{exoresolu}

\begin{exemplebox}[Exemples traduits]
\begin{itemize}
\item \eng{Could I have a receipt, please?} = « Pourrais-je avoir un reçu, s'il vous plaît ? »
\item \eng{Thank you for using our services.} = « Merci d'avoir utilisé nos services. » (attention : \eng{thank you for} + nom ou verbe en \emph{-ing})
\item \eng{Here you are.} = « Voici. » (dit en tendant l'objet)
\item \eng{It depends on the weight.} = « Cela dépend du poids. » (\eng{to depend on} : la préposition \eng{on} est obligatoire)
\end{itemize}
\end{exemplebox}

\begin{attention}[Here you are ou Here we are ?]
\eng{Here you are} signifie « voici » : on le dit \textbf{quand on donne quelque chose} à quelqu'un. Ne le confonds pas avec \eng{Here we are}, qui signifie « nous voilà (arrivés) ». Autre piège : \eng{receipt} se prononce « ri-site » — le \emph{p} ne se prononce pas, contrairement au français « réception ». Enfin, \eng{to fill in a form} (GB) ou \eng{to fill out a form} (US) : jamais \emph{to fill a form} seul.
\end{attention}

\courssub{Dialogue modèle : prendre rendez-vous chez le médecin par téléphone}

\begin{textebox}[Making an appointment with the doctor]
— Good morning, Dr Ndi's clinic. How may I help you?

— Good morning. I would like to make an appointment with the doctor, please.

— Certainly. May I have your name, please?

— My name is Amina Fotso.

— Thank you, Mrs Fotso. Is it urgent, or can it wait a few days?

— Well, I have had a strong headache since Friday, so I would prefer to come as soon as possible.

— I understand. Could you come on Tuesday at ten o'clock?

— I'm afraid I work on Tuesday morning. Would Wednesday afternoon be possible?

— Let me check\ldots\ Yes, the doctor is free on Wednesday at half past two.

— That's fine, thank you. Could you tell me what I should bring?

— Please bring your identity card and your health insurance card.

— All right. Thank you very much for your help. Goodbye!

— You're welcome, Mrs Fotso. See you on Wednesday. Goodbye!
\end{textebox}

\textbf{Glossaire :} \eng{appointment} (nom) = rendez-vous ; \eng{urgent} (adjectif) = urgent ; \eng{as soon as possible} (locution) = dès que possible ; \eng{I'm afraid\ldots} (formule) = je crains que\ldots / malheureusement\ldots ; \eng{health insurance card} (nom) = carte d'assurance maladie ; \eng{half past two} (locution) = deux heures et demie.

\begin{exemplebox}[Les étapes d'une demande de service]
Observe comment les deux dialogues suivent la même structure en cinq étapes : salutation, requête, détails, confirmation, remerciements. Cette structure est un modèle universel que tu peux réutiliser à la banque, à l'hôpital, à la poste ou au restaurant.
\end{exemplebox}

\begin{popfigure}
\begin{tikzpicture}[node distance=0.5cm, every node/.style={font=\small}]
\node[draw=popblue, fill=popblueL, rounded corners, thick, minimum width=4.6cm, minimum height=0.95cm, align=center] (s1) {\textbf{1. Greeting}\\ \eng{Good morning, madam.}};
\node[draw=popgreen, fill=popgreenL, rounded corners, thick, minimum width=4.6cm, minimum height=0.95cm, align=center, below=of s1] (s2) {\textbf{2. Request}\\ \eng{I would like to send this parcel.}};
\node[draw=poporange, fill=poporangeL, rounded corners, thick, minimum width=4.6cm, minimum height=0.95cm, align=center, below=of s2] (s3) {\textbf{3. Details}\\ \eng{Could you fill in this form?}};
\node[draw=poppurple, fill=poppurpleL, rounded corners, thick, minimum width=4.6cm, minimum height=0.95cm, align=center, below=of s3] (s4) {\textbf{4. Confirmation}\\ \eng{That will be 2,500 francs.}};
\node[draw=popgold, fill=popgoldL, rounded corners, thick, minimum width=4.6cm, minimum height=0.95cm, align=center, below=of s4] (s5) {\textbf{5. Thanks}\\ \eng{Thank you for your help.}};
\draw[-{Stealth}, thick, popdark] (s1) -- (s2);
\draw[-{Stealth}, thick, popdark] (s2) -- (s3);
\draw[-{Stealth}, thick, popdark] (s3) -- (s4);
\draw[-{Stealth}, thick, popdark] (s4) -- (s5);
\end{tikzpicture}
\legende{Organigramme : les cinq étapes d'une demande de service réussie.}
\end{popfigure}

\begin{center}
\begin{tabularx}{\linewidth}{|X|X|X|}\hline
\rowcolor{popblueL} \textbf{Étape} & \textbf{Formules anglaises} & \textbf{Équivalent français} \\ \hline
Greeting & \eng{Good morning. How may I help you?} & Bonjour. Puis-je vous aider ? \\ \hline
Request & \eng{I would like to\ldots\ / Could I have\ldots?} & Je voudrais\ldots\ / Pourrais-je avoir\ldots ? \\ \hline
Details & \eng{Could you fill in this form? / May I have your name?} & Pourriez-vous remplir ce formulaire ? / Votre nom, s'il vous plaît ? \\ \hline
Confirmation & \eng{That will be\ldots\ / Let me check.} & Cela fera\ldots\ / Je vérifie. \\ \hline
Thanks & \eng{Thank you for your help. / You're welcome.} & Merci de votre aide. / Je vous en prie. \\ \hline
\end{tabularx}
\end{center}

\courssub{Réemploi guidé : dialogue à trous}

\begin{exoresolu}[At the bank — complète le dialogue]
\textbf{Énoncé.} Complète ce dialogue avec les mots et expressions de la leçon : \eng{receipt — fee — fill in — I would like — Could you — Here you are}.

— Good morning, sir. \ldots\ldots\ldots\ldots\ldots\ to open a bank account, please.

— Certainly, madam. \ldots\ldots\ldots\ldots\ldots\ \ldots\ldots\ldots\ldots\ldots\ this application form, please?

— Of course. Is there a \ldots\ldots\ldots\ldots\ldots\ for opening the account?

— Yes, there is a small charge of 5,000 francs.

— All right. \ldots\ldots\ldots\ldots\ldots\ . And could I have a \ldots\ldots\ldots\ldots\ldots\ , please?

— Of course. Here it is. Thank you for your visit, madam.
\tcblower
\textbf{Solution détaillée.}
\begin{enumerate}
\item \eng{I would like} to open a bank account — formule de requête polie + infinitif.
\item \eng{Could you} \eng{fill in} this application form? — demande d'action ; collocation \eng{to fill in a form}.
\item Is there a \eng{fee} for opening the account? — \eng{fee} = frais (faux ami : pas « fée » !).
\item \eng{Here you are}. — « voici », en tendant l'argent.
\item Could I have a \eng{receipt}, please? — reçu, prononcé « ri-site ».
\end{enumerate}
\end{exoresolu}

\courssub{Production personnelle}

\begin{exercice}[À toi d'écrire]
Rédige \textbf{5 phrases personnelles} en utilisant chacune un mot nouveau de la leçon : \eng{parcel}, \eng{fee}, \eng{receipt}, \eng{appointment}, \eng{to fill in}. Situe tes phrases dans ton quotidien (Yaoundé, Douala, ton école, le marché\ldots).

\emph{Modèles possibles :}
\begin{enumerate}
\item \eng{Last week, my uncle sent me a parcel from Douala.}
\item \eng{The fee for a bendskin ride in Yaoundé depends on the distance.}
\item \eng{Always keep your receipt when you pay your school fees.}
\item \eng{I made an appointment with the dentist for Saturday morning.}
\item \eng{My father asked me to fill in the form for him.}
\end{enumerate}
\end{exercice}

\begin{experience}[Jeu de rôle : au guichet]
Par groupes de deux : l'élève A est le client, l'élève B est l'employé. Choisissez un lieu (poste, banque, clinique, agence de voyage). Le client doit : saluer, formuler une requête avec \eng{I would like} ou \eng{Could I\ldots?}, répondre à deux questions de détails, demander un reçu ou une confirmation, et remercier. L'employé doit utiliser \eng{How may I help you?}, \eng{Certainly}, \eng{Here you are} et \eng{Thank you for using our services}. Puis échangez les rôles.
\end{experience}

\begin{aretenir}[L'essentiel de la section]
\begin{itemize}
\item Une demande de service suit cinq étapes : \textbf{greeting → request → details → confirmation → thanks}.
\item Les formules clés : \eng{I would like to\ldots}, \eng{Could you\ldots?}, \eng{Could I have\ldots?}, \eng{Here you are}, \eng{Thank you for your help}.
\item Les collocations à mémoriser : \eng{send a parcel}, \eng{fill in a form}, \eng{make an appointment}, \eng{pay a fee}, \eng{keep a receipt}.
\item Pièges : \eng{receipt} se prononce « ri-site » ; \eng{Here you are} (voici) n'est pas \eng{Here we are} (nous voilà) ; \eng{fee} = frais, pas « fée » ; dans une question indirecte, l'ordre est sujet + verbe : \eng{Could you tell me where I can buy stamps?}
\end{itemize}
\end{aretenir}

\newpage

\coursec{Activité d'intégration}

\coursec{Lesson 19 — Vocabulary: Words and expressions related to requesting for services}

\courssub{Champ lexical 1 : Demander un service — Noms et verbes de base}

\begin{definition}[Vocabulaire essentiel : les services courants]
Le tableau ci-dessous regroupe les noms et verbes indispensables pour formuler une demande de service dans un contexte administratif, bancaire, hôtelier, médical ou postal (contexte camerounais et international).
\begin{center}
\begin{tabularx}{\linewidth}{|X|X|X|}\hline
\rowcolor{popblueL} \textbf{Français} & \textbf{English} & \textbf{Collocation type} \\ \hline
un rendez-vous & \eng{an appointment} & \eng{make an appointment} \\ \hline
une réservation / un booking & \eng{a booking} & \eng{book a room / a table} \\ \hline
un compte (bancaire) & \eng{an account} & \eng{open an account} \\ \hline
des frais / une taxe / un droit & \eng{a fee} & \eng{pay a fee} \\ \hline
un formulaire & \eng{a form} & \eng{fill in a form} (UK) / \eng{fill out a form} (US) \\ \hline
un reçu & \eng{a receipt} & \eng{ask for a receipt} / \eng{get a receipt} \\ \hline
un colis & \eng{a parcel} & \eng{send a parcel} \\ \hline
un certificat médical & \eng{a medical certificate} & \eng{request a medical certificate} \\ \hline
de l'argent & \eng{money} & \eng{deposit money} / \eng{withdraw money} \\ \hline
une remise / une réduction & \eng{a discount} & \eng{ask for a discount} \\ \hline
la disponibilité & \eng{availability} & \eng{check availability} \\ \hline
le guichet & \eng{the counter} / \eng{the desk} & \eng{at the service desk} \\ \hline
le personnel & \eng{the staff} (sing.) & \eng{The staff is helpful.} \\ \hline
aider / assister & \eng{to assist} & \eng{How may I assist you?} \\ \hline
\end{tabularx}
\end{center}
\end{definition}

\begin{propriete}[Rappel : \eng{staff} et \eng{assist}]
\begin{itemize}
\item \eng{Staff} est un nom collectif \textbf{singulier} en anglais britannique : \eng{The staff is very efficient.} (Le personnel est très efficace.)
\item \eng{To assist} signifie \textbf{aider, porter assistance}. C'est un \cle{faux ami} : il ne signifie pas « assister à » (un spectacle, une réunion) qui se dit \eng{to attend}.
\end{itemize}
\end{propriete}

\courssub{Champ lexical 2 : Formules fonctionnelles de la requête — Grammaire de la politesse}

\begin{propriete}[Les quatre piliers de la requête polie]
Pour demander un service en anglais standard (registre formel / neutre), on utilise principalement quatre structures. L'ordre des mots et la forme verbale sont stricts.
\begin{center}
\begin{tabularx}{\linewidth}{|X|X|X|}\hline
\rowcolor{popblueL} \textbf{Structure} & \textbf{Exemple} & \textbf{Niveau / Nuance} \\ \hline
\eng{I would like to} + BV (infinitif) & \eng{I would like to book a room.} & Standard, clair, professionnel. \\ \hline
\eng{Could you} + BV (infinitif sans \eng{to}) & \eng{Could you fill in this form?} & Poli, demande d'action à l'interlocuteur. \\ \hline
\eng{Would you mind} + V-\emph{ing} (gérondif) & \eng{Would you mind waiting a moment?} & Très poli, demande d'effort / patience. \\ \hline
\eng{I was wondering if} + sujet + \eng{could} / \eng{would} & \eng{I was wondering if you could help me.} & Indirect, très courtois, « \emph{hedging} ». \\ \hline
\end{tabularx}
\end{center}
\end{propriete}

\begin{propriete}[La question indirecte (\eng{Embedded question})]
Après \eng{Could you tell me...}, \eng{I'd like to know...}, \eng{Do you know...}, \textbf{pas d'inversion du sujet et de l'auxiliaire}. L'ordre est : \textbf{Sujet + Verbe}.
\end{propriete}

\begin{exemplebox}[Ordre des mots : Question directe vs Indirecte]
\begin{itemize}
\item Directe : \eng{Where is the counter?} $\rightarrow$ Indirecte : \eng{Could you tell me \textbf{where the counter is}?}
\item Directe : \eng{How much is the fee?} $\rightarrow$ Indirecte : \eng{I was wondering \textbf{how much the fee is}.}
\item Directe : \eng{Does the train leave at 8?} $\rightarrow$ Indirecte : \eng{Do you know \textbf{if the train leaves at 8}?} (si / whether)
\end{itemize}
\end{exemplebox}


\courssub{Collocations et expressions idiomatiques}

\begin{definition}[Collocations fortes du thème « Services »]
Ces associations mots-à-mots sont figées. Les apprendre par blocs (\cle{chunks}) évite les calques du français (*\eng{make a reservation} pour une chambre $\rightarrow$ \eng{book a room} ; *\eng{do a form} $\rightarrow$ \eng{fill in a form}).
\begin{center}
\begin{tabularx}{\linewidth}{|X|X|}\hline
\rowcolor{popblueL} \textbf{Collocation (Verbe + Nom)} & \textbf{Traduction / Contexte} \\ \hline
\eng{make an appointment} & Prendre rendez-vous (médecin, administration) \\ \hline
\eng{book a room / a table / a ticket} & Réserver (hôtel, restaurant, voyage) \\ \hline
\eng{open / close an account} & Ouvrir / fermer un compte (banque) \\ \hline
\eng{deposit / withdraw money} & Déposer / retirer de l'argent \\ \hline
\eng{pay a fee / a fine / a bill} & Payer des frais / une amende / une facture \\ \hline
\eng{fill in a form} (UK) & Remplir un formulaire \\ \hline
\eng{ask for / get a receipt} & Demander / obtenir un reçu \\ \hline
\eng{send a parcel / a package} & Envoyer un colis \\ \hline
\eng{check availability} & Vérifier les disponibilités \\ \hline
\eng{request a certificate} & Demander un certificat (formel) \\ \hline
\end{tabularx}
\end{center}
\end{definition}

\begin{savaistu}[Expressions idiomatiques utiles au guichet]
\begin{itemize}
\item \eng{“To kill two birds with one stone”} : Faire d'une pierre deux coups (ex. \eng{I'll open an account and deposit money at the same time — kill two birds with one stone.}).
\item \eng{“No problem” / “Not at all” / “Certainly”} : Réponses standards à \eng{Would you mind...?} (Attention : \eng{Would you mind?} $\rightarrow$ \eng{Not at all} = « Je ne m'en formalise pas / Avec plaisir »).
\item \eng{“To be fully booked”} : Complet (hôtel, avion). \eng{I'm afraid the hotel is fully booked.}
\item \eng{“To wait in the queue”} (UK) / \eng{“To wait in line”} (US) : Faire la queue.
\end{itemize}
\end{savaistu}

\courssub{Faux amis et pièges des francophones}

\begin{attention}[Pièges majeurs à éviter — Liste de contrôle]
\begin{itemize}
\item \textbf{\eng{Assist} $\neq$ Assister à}. \eng{Assist} = aider. \eng{Attend} = assister à (un événement). \eng{“May I assist you?”} = « Puis-je vous aider ? »
\item \textbf{\eng{Actually} $\neq$ Actuellement}. \eng{Actually} = en fait, en réalité. \eng{Currently} / \eng{At the moment} = actuellement.
\item \textbf{\eng{Receipt} $\neq$ Recette}. \eng{Receipt} = reçu (preuve de paiement). \eng{Recipe} = recette (cuisine). \eng{Prescription} = ordonnance (médecin).
\item \textbf{\eng{Form} $\neq$ Forme}. \eng{Form} = formulaire (admin). \eng{Shape} / \eng{Form} (philosophie) = forme.
\item \textbf{\eng{Booking} $\neq$ Booking (franglais)}. En français on dit « un booking », en anglais \eng{a booking} ou \eng{a reservation}. Verbe : \eng{to book}.
\item \textbf{\eng{Queue} $\neq$ Queue (file d'attente informatique seulement)}. \eng{Queue} = file d'attente physique. Prononciation : « \textbf{kiou} » /kjuː/.
\item \textbf{\eng{Appointment} $\neq$ Appointement}. \eng{Appointment} = rendez-vous. \eng{Salary} / \eng{Wages} = appointements / salaire.
\item \textbf{Ordre des mots : \eng{I would like \textbf{to} + verbe}} (pas *\eng{I would like \textbf{that} I...}).
\item \textbf{\eng{Would you mind} + \textbf{V-ing}} (pas *\eng{Would you mind to...}).
\item \textbf{\eng{Staff} au singulier} : \eng{The staff \textbf{is} here.} (pas *\eng{are}).
\end{itemize}
\end{attention}

\courssub{Familles de mots, prononciation et orthographe}

\begin{propriete}[Dérivation morphologique (Word Formation)]
Construire le nom, le verbe, l'adjectif ou la personne permet de multiplier le lexique.
\begin{center}
\begin{tabular}{|l|l|l|l|l|}\hline
\rowcolor{popblueL} \textbf{Verbe} & \textbf{Nom (action/état)} & \textbf{Nom (agent)} & \textbf{Adjectif / Participe} & \textbf{Traduction} \\ \hline
\eng{request} & \eng{a request} & \eng{a requester} & \eng{requested} & demander / demande \\ \hline
\eng{book} & \eng{a booking} & \eng{a booker} & \eng{booked} & réserver / réservation \\ \hline
\eng{appoint} & \eng{an appointment} & -- & \eng{appointed} & nominer / rendez-vous \\ \hline
\eng{receive} & \eng{a receipt} & \eng{a receiver} / \eng{recipient} & \eng{received} & recevoir / reçu \\ \hline
\eng{assist} & \eng{assistance} & \eng{an assistant} & \eng{assisted} & aider / assistance \\ \hline
\eng{apply} & \eng{an application} & \eng{an applicant} & \eng{applied} & postuler / candidature \\ \hline
\eng{register} & \eng{registration} & \eng{a registrant} & \eng{registered} & inscrire / inscription \\ \hline
\end{tabular}
\end{center}
\end{propriete}

\begin{exemplebox}[Focus Prononciation : Mots pièges]
Entraîne-toi à lire ces mots à haute voix. La prononciation approximative en lettres françaises est donnée entre guillemets.
\begin{itemize}
\item \eng{Receipt} /rɪˈsiːt/ $\rightarrow$ « \textbf{ri-ssite} » (le \eng{p} est muet, \eng{ei} = /iː/).
\item \eng{Queue} /kjuː/ $\rightarrow$ « \textbf{kiou} » (seul le \eng{q} se prononce).
\item \eng{Appointment} /əˈpɔɪntmənt/ $\rightarrow$ « \textbf{a-pointe-mente} » (accent sur \eng{point}, \eng{a-} faible).
\item \eng{Form} /fɔːm/ $\rightarrow$ « \textbf{forme} » (voyelle longue, \eng{r} non roulé).
\item \eng{Parcel} /ˈpɑːsl/ $\rightarrow$ « \textbf{par-seul} » (accent sur 1ère syllabe, \eng{cel} = /sl/).
\item \eng{Available} /əˈveɪləbl/ $\rightarrow$ « \textbf{a-véi-la-beul} » (accent sur \eng{vai}, \eng{able} = /əbl/).
\end{itemize}
\end{exemplebox}

\begin{popfigure}
\begin{tikzpicture}[mindmap, grow cyclic, every node/.style=concept, concept color=popblue,
  level 1/.append style={level distance=3.5cm, sibling angle=90},
  level 2/.append style={level distance=2.5cm, sibling angle=45}]
\node {Request / Service}
  child[concept color=popgreen] { node {Nouns}
    child { node {Appointment} }
    child { node {Booking} }
    child { node {Receipt} }
    child { node {Fee} }
    child { node {Form} }
  }
  child[concept color=poporange] { node {Verbs}
    child { node {Book} }
    child { node {Fill in} }
    child { node {Pay} }
    child { node {Send} }
    child { node {Assist} }
  }
  child[concept color=poppurple] { node {People}
    child { node {Clerk} }
    child { node {Customer} }
    child { node {Staff} }
    child { node {Applicant} }
  }
  child[concept color=poppink] { node {Places}
    child { node {Counter} }
    child { node {Desk} }
    child { node {Office} }
    child { node {Centre} }
  };
\end{tikzpicture}
\legende{Carte mentale : Familles lexicales du thème « Requesting Services »}
\end{popfigure}

\courssub{Le lexique en contexte : Texte et dialogues}

\courssub{Support authentique : Panneau d'affichage — Buea Community Service Centre}

\begin{textebox}[Notice board — Buea Community Service Centre]
WELCOME TO BUEA COMMUNITY SERVICE CENTRE

Our counters are open from 8 a.m. to 4 p.m., Monday to Friday.

Counter 1 — BANK: open an account, deposit or withdraw money, pay a fee.

Counter 2 — HOTEL BOOKINGS: book a room, check availability, ask for a discount.

Counter 3 — HEALTH CENTRE: make an appointment, see a doctor, request a medical certificate.

Counter 4 — POST OFFICE: send a parcel, buy stamps, fill in a form, ask for a receipt.

Please take a ticket and wait in the queue. Our staff will assist you.
\end{textebox}

\begin{definition}[Glossaire du panneau — Mots difficiles]
\begin{itemize}
\item \eng{counter} (n.) : le guichet ; \eng{to deposit money} : déposer de l'argent ; \eng{to withdraw} : retirer (de l'argent) ;
\item \eng{availability} (n.) : la disponibilité ; \eng{discount} (n.) : la remise, la réduction ;
\item \eng{medical certificate} (n.) : le certificat médical ; \eng{parcel} (n.) : le colis ;
\item \eng{staff} (n. singulier) : le personnel ; \eng{to assist} : aider, assister (faux ami : ce n'est pas « assister à »).
\item \eng{stamps} (n. pl.) : des timbres-poste.
\end{itemize}
\end{definition}

\courssub{Activité d'intégration : Jeu de rôle au Centre de services de Buea}

\courssub{Objectifs de l'activité}

À l'issue de cette activité d'intégration, tu seras capable de :
\begin{itemize}
\item réemployer le \cle{lexique de la demande de service} (\eng{request, appointment, booking, fee, form, receipt}) dans une situation de communication authentique ;
\item utiliser correctement les \cle{formules polies de requête} (\eng{Could you...?}, \eng{I would like to...}, \eng{Would you mind...?}) et les réponses appropriées (acceptation ou refus poli) ;
\item employer au moins trois \cle{collocations} étudiées dans la leçon (\eng{make an appointment}, \eng{fill in a form}, \eng{pay a fee}, \eng{book a room}, \eng{open an account}) ;
\item produire un dialogue oral fluide et intelligible, avec une prononciation soignée des mots difficiles (\eng{receipt} « ri-ssite », \eng{queue} « kiou », \eng{appointment} « a-pointe-mente ») ;
\item écouter un dialogue pour y repérer des formules précises (compréhension de l'oral).
\end{itemize}

\courssub{Situation}

Tu es en vacances à Buea, dans la région du Sud-Ouest. Pendant la semaine, tu dois régler plusieurs démarches dans des services publics et privés : ouvrir un compte à la banque, réserver une chambre d'hôtel pour ton cousin qui arrive de Lagos, prendre rendez-vous à l'hôpital et envoyer un colis à ta sœur qui étudie à Bamenda. À chaque fois, tu dois t'adresser à un employé au guichet (\eng{at the service desk}) et formuler ta demande poliment, en anglais.

\courssub{Consignes}

Travaillez par groupes de deux. Suivez les étapes dans l'ordre.

\begin{enumerate}
\item \textbf{Choose your scene.} Avec ton partenaire, choisissez un des quatre guichets du panneau : la banque, l'hôtel, l'hôpital ou la poste. L'élève A joue le client (\eng{the customer}), l'élève B joue l'employé (\eng{the clerk}).
\item \textbf{Prepare your vocabulary.} Soulignez dans le panneau les mots et expressions utiles pour votre scène. Ajoutez au moins trois collocations de la leçon : \eng{make an appointment}, \eng{fill in a form}, \eng{pay a fee}, \eng{book a room}, \eng{open an account}, \eng{send a parcel}.
\item \textbf{Write your dialogue.} Rédigez un dialogue de 8 à 10 répliques. Le client doit utiliser au moins trois formules polies de requête (\eng{Could you...?}, \eng{I would like to...}, \eng{Would you mind...?}, \eng{I was wondering if...}). L'employé doit répondre par des formules d'acceptation (\eng{Certainly.}, \eng{Of course, madam.}) ou de refus poli (\eng{I'm afraid...}, \eng{I'm sorry, but...}).
\item \textbf{Check your English.} Relisez votre dialogue : vérifiez l'ordre des mots dans les questions indirectes (\eng{Could you tell me where the counter is?} et non \emph{where is the counter}), la politesse du ton et l'orthographe (\eng{receipt}, \eng{appointment}, \eng{available}).
\item \textbf{Rehearse.} Répétez votre scène à voix haute. Soignez la prononciation et l'intonation montante des questions polies.
\item \textbf{Perform and listen.} Jouez la scène devant la classe. Pendant la prestation des autres binômes, cochez sur la grille d'écoute chaque formule que vous entendez.
\item \textbf{Vote.} La classe désigne le binôme le plus clair, le plus poli et le plus naturel : il reçoit le titre de \eng{Best Customer Service of the Day}.
\end{enumerate}

\begin{experience}[Grille d'écoute — Listening grid]
Pendant chaque prestation, coche (\checkmark) les éléments entendus :
\begin{center}
\begin{tabularx}{\linewidth}{|X|c|}\hline
\rowcolor{popblueL} \textbf{Formula heard} & \textbf{Tick} \\ \hline
\eng{Could you...?} / \eng{Can I...?} & \\ \hline
\eng{I would like to...} & \\ \hline
\eng{Would you mind...?} / \eng{I was wondering if...} & \\ \hline
\eng{make an appointment} / \eng{book a room} / \eng{open an account} & \\ \hline
\eng{fill in a form} / \eng{pay a fee} / \eng{ask for a receipt} & \\ \hline
Acceptation : \eng{Certainly.} / \eng{Of course.} / \eng{No problem.} & \\ \hline
Refus poli : \eng{I'm afraid...} / \eng{I'm sorry, but...} & \\ \hline
Formule de clôture : \eng{Thank you for your help.} / \eng{Have a nice day.} & \\ \hline
\end{tabularx}
\end{center}
\end{experience}

\courssub{Ressources à mobiliser}

\begin{aretenir}[Boîte à outils de la requête polie]
\begin{itemize}
\item \textbf{Demander} : \eng{Could you help me, please?} ; \eng{I would like to open an account.} ; \eng{Would you mind filling in this form?} ; \eng{I was wondering if you could send this parcel.}
\item \textbf{Accepter} : \eng{Certainly, sir.} ; \eng{Of course, madam.} ; \eng{No problem at all.}
\item \textbf{Refuser poliment} : \eng{I'm afraid the doctor is not available today.} ; \eng{I'm sorry, but the rooms are fully booked.}
\item \textbf{Remercier et conclure} : \eng{Thank you very much for your help.} ; \eng{You're welcome. Have a nice day!}
\end{itemize}
\end{aretenir}

\begin{attention}[Pièges à éviter dans ton dialogue]
\begin{itemize}
\item Ne dis pas \emph{« I want a room »} : trop direct, voire impoli. Dis \eng{I would like to book a room, please.}
\item Après \eng{Would you mind}, utilise le verbe en \emph{-ing} : \eng{Would you mind waiting a moment?}
\item Question indirecte : pas d'inversion. \eng{Could you tell me how much the fee is?} (et non \emph{how much is the fee}).
\item Faux ami : \eng{to assist} signifie « aider » ; « assister à » se dit \eng{to attend}.
\item \eng{staff} est singulier : \eng{The staff is very helpful.}
\end{itemize}
\end{attention}

\courssub{Proposition de production et corrigé commenté}

\begin{exoresolu}[Dialogue modèle — At the Post Office (Counter 4)]
Énoncé : rédige et joue un dialogue de 8 à 10 répliques au guichet de la poste, avec au moins trois formules polies et trois collocations de la leçon.
\tcblower
\textbf{Modèle de dialogue :}

\eng{Clerk: Good morning, madam. Welcome to Counter 4. How may I assist you?}

\eng{Customer: Good morning. I would like to send a parcel to Bamenda, please.}

\eng{Clerk: Certainly. Could you fill in this form first? Write the name and address of the receiver clearly.}

\eng{Customer: Of course. Could you tell me how much the fee is?}

\eng{Clerk: It depends on the weight. Let me weigh your parcel... That will be two thousand CFA francs.}

\eng{Customer: All right. I was wondering if it could arrive before Friday.}

\eng{Clerk: I'm afraid express delivery is not available today, but it should arrive on Thursday.}

\eng{Customer: That's fine. Would you mind giving me a receipt, please?}

\eng{Clerk: Not at all. Here is your receipt. Keep it carefully.}

\eng{Customer: Thank you very much for your help. Have a nice day!}

\eng{Clerk: You're welcome, madam. Goodbye.}

\textbf{Commentaire des choix de langue :}
\begin{itemize}
\item \textbf{Formules polies} : quatre requêtes différentes sont utilisées (\eng{I would like to}, \eng{Could you tell me}, \eng{I was wondering if}, \eng{Would you mind giving me}) — au-delà du minimum demandé, ce qui rend le dialogue naturel.
\item \textbf{Collocations} : \eng{send a parcel}, \eng{fill in a form}, \eng{the fee}, \eng{give me a receipt} — le lexique de la leçon est bien réemployé en contexte.
\item \textbf{Question indirecte} : \eng{Could you tell me how much the fee is?} respecte l'ordre sujet + verbe, sans inversion.
\item \textbf{Refus poli} : \eng{I'm afraid express delivery is not available today} adoucit la mauvaise nouvelle et propose une solution (\eng{it should arrive on Thursday}).
\item \textbf{Grammaire} : \eng{Would you mind giving} (verbe en \emph{-ing}) ; \eng{should arrive} pour exprimer la probabilité.
\item \textbf{Clôture} : le dialogue se termine par des formules de politesse, comme dans un vrai échange au guichet.
\end{itemize}
\end{exoresolu}

\courssub{Bilan de l'activité}

Grâce à ce jeu de rôle, tu as mis en pratique l'ensemble des compétences de la leçon dans une situation de communication réelle :
\begin{itemize}
\item tu as \textbf{compris} un document authentique (le panneau du centre de services) et les dialogues de tes camarades (grille d'écoute) ;
\item tu as \textbf{réemployé le lexique} de la demande de service : noms (\eng{request, appointment, fee, receipt, parcel}), verbes (\eng{book, fill in, deposit, send}) et collocations ;
\item tu as \textbf{structuré des requêtes polies} avec \eng{Could you}, \eng{I would like}, \eng{Would you mind} et les questions indirectes ;
\item tu as \textbf{produit un dialogue oral} fluide, avec une prononciation et une intonation soignées ;
\item tu as évité les \textbf{faux amis et calques du français} (\eng{assist} / « assister », \eng{I want} trop direct).
\end{itemize}

\begin{exercice}[Pour aller plus loin — à la maison]
Rédige un second dialogue (10 répliques) dans un guichet différent de celui joué en classe : à la banque, à l'hôtel ou à l'hôpital. Utilise au moins quatre formules polies, trois collocations et un refus poli de la part de l'employé. Tu le présenteras au début de la leçon suivante.
\end{exercice}

\begin{corrige}[Pour aller plus loin]
Exemple de réponse attendue (hôtel) : \eng{Receptionist: Good evening, sir. How may I help you? — Customer: Good evening. I would like to book a double room for two nights, please. — Receptionist: Certainly. Could you fill in this registration form? — Customer: Sure. Could you tell me if breakfast is included? — Receptionist: I'm afraid it is not, but you can pay a small extra fee. — Customer: All right. Would you mind showing me the room first? — Receptionist: Not at all. Follow me, please. — Customer: Thank you for your help. — Receptionist: You're welcome, sir.} — Critères de correction : présence des formules polies, collocations correctes, question indirecte sans inversion, refus poli avec \eng{I'm afraid}, clôture courtoise.
\end{corrige}

\newpage

\coursec{Méthodes et exercices résolus}

\courssub{Savoir-faire 1 : Identifier et employer le vocabulaire de la demande de service}

\begin{methode}[Comment identifier et employer le vocabulaire de la requête]
Pour bien utiliser le lexique de la demande de service (\eng{requesting services}), suis ces étapes :

\begin{enumerate}
\item \textbf{Repère la situation de communication.} Qui parle à qui ? Un client à un employé de banque, un passager à un guichetier, un élève à un secrétaire ? La relation entre les personnes détermine le degré de politesse.
\item \textbf{Choisis le verbe ou le nom adapté au service.} Mémorise les couples verbe + nom : \eng{to book a room} (réserver une chambre), \eng{to order a meal} (commander un repas), \eng{to apply for a loan} (demander un prêt), \eng{to request information} (demander des informations), \eng{to fill in a form} (remplir un formulaire), \eng{to deliver a parcel} (livrer un colis).
\item \textbf{Construis la requête avec la formule de politesse correcte.} Structure de base : \eng{Could you + verbe à l'infinitif...?} ou \eng{I would like to + infinitif...}
\item \textbf{Vérifie la préposition.} En anglais, chaque verbe impose sa préposition : \eng{to ask \textbf{for} something}, \eng{to apply \textbf{for} a job}, \eng{to wait \textbf{for} someone}, \eng{to pay \textbf{for} a service}.
\item \textbf{Réemploie le mot dans une phrase personnelle} pour l'ancrer en mémoire : \eng{I would like to book a taxi to the airport.}
\end{enumerate}
\end{methode}

\begin{aretenir}[Les verbes-clés de la demande de service]
\begin{center}
\begin{tabularx}{\linewidth}{|X|X|X|}\hline
\rowcolor{popblueL} English & Français & Exemple \\ \hline
to request & demander (soutenu) & to request a receipt \\ \hline
to book / to reserve & réserver & to book a flight \\ \hline
to order & commander & to order a drink \\ \hline
to apply for & faire une demande de & to apply for a visa \\ \hline
to deliver & livrer & to deliver goods \\ \hline
to repair / to fix & réparer & to fix a phone \\ \hline
to assist & aider, assister & to assist a customer \\ \hline
\end{tabularx}
\end{center}
\end{aretenir}

\begin{exoresolu}[Vocabulary in context — Compléter un dialogue]
Énoncé — Complete the dialogue with the following words: \emph{book, deliver, fill in, receipt, appointment}.

\eng{Receptionist: Good morning, sir. How can I help you?}

\eng{Customer: Good morning. I would like to (1) .......... a room for two nights.}

\eng{Receptionist: Certainly. Could you (2) .......... this form, please?}

\eng{Customer: Of course. Can you also (3) .......... my breakfast to the room tomorrow?}

\eng{Receptionist: No problem. Here is your (4) .......... for the payment.}

\eng{Customer: Thank you. I also need an (5) .......... with the manager this afternoon.}

\tcblower
Solution détaillée :

\begin{enumerate}
\item \textbf{book} — On « réserve » une chambre : la collocation correcte est \eng{to book a room}. Après \eng{would like to}, le verbe reste à l'infinitif sans \eng{to}.
\item \textbf{fill in} — La collocation avec \eng{a form} est \eng{to fill in a form} (remplir un formulaire). Après \eng{Could you}, infinitif sans \eng{to}.
\item \textbf{deliver} — \eng{to deliver breakfast to the room} = livrer le petit-déjeuner dans la chambre. Attention à la préposition \eng{to} (livrer \emph{à} un endroit).
\item \textbf{receipt} — Après un paiement, on reçoit un \eng{receipt} (un reçu). Ne pas confondre avec \eng{recipe} (recette de cuisine) : faux ami classique.
\item \textbf{appointment} — \eng{an appointment with the manager} = un rendez-vous avec le directeur. On dit \eng{to have / to make an appointment}, jamais \eng{to do an appointment}.
\end{enumerate}

Réponse rédigée : \eng{I would like to book a room for two nights. Could you fill in this form, please? Can you also deliver my breakfast to the room tomorrow? Here is your receipt for the payment. I also need an appointment with the manager this afternoon.}

Conclusion : en contexte, le choix du mot dépend de la \cle{collocation} (le mot qui l'accompagne naturellement), pas seulement de la traduction française.
\end{exoresolu}

\courssub{Savoir-faire 2 : Réemployer le lexique à l'oral et à l'écrit}

\begin{methode}[Comment formuler une requête polie à l'oral et à l'écrit]
Pour demander un service correctement, adapte ta formule au degré de formalité :

\begin{enumerate}
\item \textbf{Évalue le degré de formalité.} À un ami : \eng{Can you...?} À un inconnu ou un supérieur : \eng{Could you...?} / \eng{Would you mind...?} Dans une lettre officielle : \eng{I would be grateful if you could...}
\item \textbf{Utilise la structure grammaticale exacte :}
\begin{itemize}
\item \eng{Could you + infinitif} : \eng{Could you repair my bicycle?}
\item \eng{Would you mind + verbe en -ing} : \eng{Would you mind checking my file?}
\item \eng{I would like + nom} ou \eng{I would like to + infinitif} : \eng{I would like some information.} / \eng{I would like to open an account.}
\item \eng{I would be grateful if you could + infinitif} : \eng{I would be grateful if you could send me the documents.}
\end{itemize}
\item \textbf{Ajoute une justification si nécessaire} : \eng{Could you come quickly, please? My car has broken down on the Douala road.}
\item \textbf{Termine par une formule de remerciement} : \eng{Thank you very much.} / \eng{I really appreciate your help.}
\item \textbf{À l'écrit (lettre formelle),} respecte le plan : adresse, objet (\eng{Re: Request for...}), salutation (\eng{Dear Sir or Madam}), corps, formule finale (\eng{Yours faithfully}).
\end{enumerate}
\end{methode}

\begin{attention}[La structure de « Would you mind »]
Après \eng{Would you mind}, le verbe est obligatoirement en \cle{-ing} : \eng{Would you mind helping me?} et jamais \eng{Would you mind to help me?} Pour répondre « oui » (j'accepte), on dit \eng{Not at all} ou \eng{Of course not} — et non « yes », qui signifierait un refus !
\end{attention}

\begin{exoresolu}[Transformation — Du familier au formel]
Énoncé — Rewrite each request to make it more formal, as in the example.

Example: \eng{Can you help me?} $\rightarrow$ \eng{Could you help me, please?}

\begin{enumerate}
\item \eng{Give me a taxi to the station.}
\item \eng{I want information about your prices.}
\item \eng{Send me the documents.}
\item \eng{Help me to carry these bags.}
\end{enumerate}

\tcblower
Solution détaillée :

\begin{enumerate}
\item \eng{Could you call me a taxi to the station, please?} — On remplace l'impératif direct (trop brutal en anglais) par \eng{Could you + infinitif}. La collocation correcte est \eng{to call a taxi} (appeler un taxi), pas « give ».
\item \eng{I would like some information about your prices, please.} — \eng{I want} est trop direct ; le registre poli exige \eng{I would like}. Attention : \eng{information} est \cle{indénombrable} en anglais : jamais de « s » (\eng{informations} est une erreur), on dit \eng{some information}.
\item \eng{I would be grateful if you could send me the documents.} — Formule de lettre officielle. Structure : \eng{I would be grateful if you could + infinitif}. C'est le niveau de politesse maximal.
\item \eng{Would you mind helping me to carry these bags, please?} — Après \eng{Would you mind}, le verbe prend la forme en \eng{-ing} : \eng{helping}. C'est la règle grammaticale à justifier dans la copie.
\end{enumerate}

Conclusion : pour rendre une requête plus formelle, trois réflexes : remplacer l'impératif ou \eng{Can you} par \eng{Could you / Would you mind + -ing / I would be grateful if you could}, employer \eng{I would like} au lieu de \eng{I want}, et ajouter \eng{please}.
\end{exoresolu}

\courssub{Savoir-faire 3 : Éviter les faux amis et les calques du français}

\begin{methode}[Comment déjouer les faux amis du domaine des services]
\begin{enumerate}
\item \textbf{Méfie-toi des mots qui ressemblent au français.} Un mot anglais proche d'un mot français n'a pas toujours le même sens : \eng{to demand} ne veut pas dire « demander » mais « exiger ».
\item \textbf{Vérifie la collocation, pas seulement le mot.} En français on « fait » un rendez-vous ; en anglais on dit \eng{to make an appointment} (et non « to do »). On « prend » un taxi : \eng{to take a taxi}, jamais « to get a taxi » pour ce sens.
\item \textbf{Ne traduis pas mot à mot.} « Je voudrais que vous m'envoyiez... » ne se traduit pas par « I want that you send » mais par \eng{I would like you to send me...} (structure : \eng{to want / would like + quelqu'un + to + infinitif}).
\item \textbf{Contrôle les indénombrables.} \eng{information, advice, luggage, furniture} n'ont jamais de pluriel : \eng{some advice}, \eng{a piece of information}.
\item \textbf{Fais une liste personnelle des faux amis} rencontrés dans les textes et relis-la avant chaque évaluation.
\end{enumerate}
\end{methode}

\begin{aretenir}[Tableau des faux amis du domaine des services]
\begin{center}
\begin{tabularx}{\linewidth}{|X|X|X|}\hline
\rowcolor{popblueL} Faux ami & Vrai sens & « Faux sens » français \\ \hline
to demand & exiger & demander = to ask for / to request \\ \hline
a receipt & un reçu & une recette = a recipe \\ \hline
to assist & aider & assister à = to attend \\ \hline
sensible & raisonnable & sensible = sensitive \\ \hline
actually & en fait & actuellement = currently \\ \hline
\end{tabularx}
\end{center}
\end{aretenir}

\begin{exoresolu}[Error correction — Chasse aux calques]
Énoncé — Each sentence contains one mistake caused by a French calque or a false friend. Correct it.

\begin{enumerate}
\item \eng{I demand a room for tonight, please.}
\item \eng{Can you give me some advices about hotels in Buea?}
\item \eng{I want that you repair my phone today.}
\item \eng{She assisted the meeting of the workers' union.}
\item \eng{Please give me the recipe of my payment.}
\end{enumerate}

\tcblower
Solution détaillée :

\begin{enumerate}
\item \eng{I \textbf{would like} a room for tonight, please.} — \eng{To demand} signifie « exiger » (faux ami de « demander »). Pour une requête polie : \eng{I would like} ou \eng{Could I have}.
\item \eng{Can you give me some \textbf{advice} about hotels in Buea?} — \eng{Advice} est indénombrable : pas de « s ». Pour compter : \eng{a piece of advice}.
\item \eng{I \textbf{would like you to} repair my phone today.} — Calque de « je veux que tu ». L'anglais n'emploie pas \eng{want that} ; structure correcte : \eng{to want / would like + complément + to + infinitif}.
\item \eng{She \textbf{attended} the meeting of the workers' union.} — \eng{To assist} = aider ; « assister à » (être présent) se dit \eng{to attend}. Faux ami classique.
\item \eng{Please give me the \textbf{receipt} for my payment.} — \eng{Recipe} = recette de cuisine ; un reçu de paiement est \eng{a receipt}. Prononciation : « ri-site » (le « p » ne se prononce pas). On préfère la préposition \eng{for} : \eng{a receipt for a payment}.
\end{enumerate}

Conclusion : avant d'écrire un mot anglais qui ressemble au français, pose-toi la question : « est-ce le vrai sens, ou un faux ami ? » et vérifie la structure de la phrase (pas de \eng{that} après \eng{want}).
\end{exoresolu}

\courssub{Savoir-faire 4 : Exploiter les familles de mots, la prononciation et l'orthographe}

\begin{methode}[Comment enrichir son lexique par les familles de mots]
\begin{enumerate}
\item \textbf{Pars d'un mot connu et construis sa famille.} Exemple avec \eng{to serve} : \eng{service} (nom de l'action), \eng{servant} (nom de personne, « domestique »), \eng{server} (celui qui sert, au restaurant).
\item \textbf{Apprends les suffixes qui changent la nature du mot :}
\begin{itemize}
\item verbe $\rightarrow$ nom : \eng{to deliver} $\rightarrow$ \eng{delivery} ; \eng{to apply} $\rightarrow$ \eng{application} ; \eng{to assist} $\rightarrow$ \eng{assistance}
\item verbe $\rightarrow$ personne : \eng{to employ} $\rightarrow$ \eng{employee / employer} ; \eng{to assist} $\rightarrow$ \eng{assistant}
\item adjectif : \eng{help} $\rightarrow$ \eng{helpful / helpless}
\end{itemize}
\item \textbf{Attention aux changements orthographiques :} \eng{to apply} $\rightarrow$ \eng{application} (le « y » disparaît) ; \eng{to deliver} $\rightarrow$ \eng{delivery} (« y » après consonne).
\item \textbf{Note la prononciation en lettres françaises} : \eng{receipt} « ri-site », \eng{customer} « keu-sto-meure », \eng{enquiry} « ine-kouaï-ri », \eng{queue} « kiou ».
\item \textbf{Teste-toi} en rédigeant une phrase avec chaque membre de la famille : \eng{The assistant offered me assistance when I asked for help.} — puis crée tes propres phrases.
\end{enumerate}
\end{methode}

\begin{exoresolu}[Word formation — Familles de mots]
Énoncé — Complete the sentences with the correct form of the word in brackets.

\begin{enumerate}
\item \eng{The .......... of my parcel took three days.} (deliver)
\item \eng{You must write an .......... letter to get this job.} (apply)
\item \eng{The shop .......... helped me to choose a phone.} (assist)
\item \eng{Thank you for your .......... during my stay.} (hospitable)
\item \eng{The bank offers a good customer .......... .} (serve)
\end{enumerate}

\tcblower
Solution détaillée :

\begin{enumerate}
\item \textbf{delivery} — Après l'article \eng{the}, il faut un nom. Verbe \eng{to deliver} $\rightarrow$ nom \eng{delivery} (livraison). Orthographe : « y » final.
\item \textbf{application} — Devant le nom \eng{letter}, il faut un nom épithète : \eng{an application letter} (une lettre de candidature). Attention à l'orthographe : le « y » de \eng{apply} devient « i » devant \eng{-cation}.
\item \textbf{assistant} — Le sujet de \eng{helped} est une personne : \eng{the shop assistant} (le vendeur). Suffixe \eng{-ant} pour la personne ; \eng{assistance} est le nom de l'action.
\item \textbf{hospitality} — Après le possessif \eng{your}, il faut un nom : \eng{hospitality} (hospitalité). Adjectif \eng{hospitable} $\rightarrow$ nom \eng{hospitality}.
\item \textbf{service} — \eng{customer service} est une collocation figée : le service clientèle. Verbe \eng{to serve} $\rightarrow$ nom \eng{service}.
\end{enumerate}

Conclusion : pour trouver la bonne forme, repère d'abord la \cle{nature grammaticale} attendue (article $\rightarrow$ nom ; sujet qui agit $\rightarrow$ nom de personne ; devant un nom $\rightarrow$ adjectif ou nom épithète), puis applique le suffixe avec la bonne orthographe.
\end{exoresolu}

\courssub{Savoir-faire 5 : Comprendre un dialogue de service et rédiger une requête}

\begin{methode}[Comment réussir une question de compréhension ou une production écrite sur les services]
\begin{enumerate}
\item \textbf{Lis d'abord les questions, puis le texte.} Tu sauras quels mots-clés repérer : noms de services (\eng{booking, delivery, repair}), prix, horaires, lieux.
\item \textbf{Repère les formules de requête dans le dialogue} : elles signalent ce que le client veut (\eng{Could you...? I would like...}).
\item \textbf{Réponds par des phrases complètes} en réemployant le vocabulaire du texte, sans recopier de longs passages.
\item \textbf{Pour une production écrite} (lettre ou dialogue de demande de service), respecte ce plan : salutation $\rightarrow$ présentation du besoin $\rightarrow$ requête précise avec formule polie $\rightarrow$ justification $\rightarrow$ remerciement et formule finale.
\item \textbf{Relis-toi} en traquant : prépositions (\eng{ask for, apply for, pay for}), indénombrables (\eng{information, advice}), formes en \eng{-ing} après \eng{Would you mind}.
\end{enumerate}
\end{methode}

\begin{exoresolu}[Reading comprehension and guided writing]
Énoncé — Read the dialogue, then answer the questions in complete sentences.

\eng{Mrs Etame: Good morning. I would like to open a savings account, please.}

\eng{Bank clerk: Certainly, madam. Could you fill in this application form and show me your identity card?}

\eng{Mrs Etame: Here you are. Could you also tell me about your interest rates?}

\eng{Bank clerk: Of course. I will give you a brochure with all the information. Would you mind waiting a moment while I check your documents?}

\eng{Mrs Etame: Not at all. Thank you for your assistance.}

Questions :
\begin{enumerate}
\item What service does Mrs Etame request?
\item What two things must she do to get the service?
\item How does the clerk ask her to wait politely?
\end{enumerate}

\tcblower
Solution détaillée :

\begin{enumerate}
\item \eng{Mrs Etame requests to open a savings account.} — On repère la formule \eng{I would like to open a savings account} dans la première réplique. La réponse reprend la structure \eng{to request to + infinitif}.
\item \eng{To get the service, she must fill in an application form and show her identity card.} — Les deux actions sont signalées par \eng{fill in} et \eng{show}. On répond avec \eng{must} car le guichetier pose une condition.
\item \eng{The clerk asks her politely by saying: “Would you mind waiting a moment while I check your documents?”} — Formule polie \eng{Would you mind + -ing} ; la réponse \eng{Not at all} de Mrs Etame confirme qu'elle accepte.
\end{enumerate}

Production guidée (modèle de réponse) — \emph{Write four lines: you ask a mechanic in Bamenda to repair your motorbike.}

\eng{Good morning, sir. My motorbike has broken down near the market. Could you repair it today, please? I need it to go to school tomorrow morning. I would be grateful if you could also check the brakes. Thank you very much for your help.}

Conclusion : en compréhension, les réponses se trouvent dans les \cle{formules de requête} du dialogue ; en production, une requête réussie combine une formule polie, une justification et un remerciement.
\end{exoresolu}

\begin{attention}[Les 5 erreurs qui coûtent des points au Bac]
\begin{enumerate}
\item \textbf{Le calque « I want that you... »} — Après \eng{want} ou \eng{would like}, jamais de \eng{that} : écris \eng{I would like you to send me the form.} Structure : \eng{want / would like + quelqu'un + to + infinitif}.
\item \textbf{Les indénombrables avec un « s »} — \eng{informations}, \eng{advices}, \eng{luggages} sont toujours faux : \eng{some information}, \eng{a piece of advice}, \eng{my luggage}.
\item \textbf{Les faux amis de la requête} — \eng{to demand} = exiger (et non « demander » : \eng{to ask for / to request}) ; \eng{a receipt} = un reçu (et non une recette : \eng{a recipe}) ; \eng{to assist} = aider (et non « assister à » : \eng{to attend}).
\item \textbf{L'oubli des prépositions} — \eng{to ask \textbf{for} something}, \eng{to apply \textbf{for} a job}, \eng{to pay \textbf{for} a service}, \eng{to wait \textbf{for} a customer}, \eng{to listen \textbf{to} the clerk}. Une préposition oubliée ou fausse coûte un point à chaque occurrence.
\item \textbf{La forme du verbe après les formules de politesse} — \eng{Could you + infinitif} (\eng{Could you help me?}), mais \eng{Would you mind + -ing} (\eng{Would you mind helping me?}). Mélanger les deux structures (\eng{Would you mind to help}) est une faute de grammaire sanctionnée.
\end{enumerate}
\end{attention}

\newpage

\coursec{Exercices}

\courssub{Je vérifie mes connaissances}

\begin{exercice}[QCM : le bon mot]
Choisis la bonne réponse pour chaque phrase.
\begin{enumerate}
\item I would like to \textbf{……} a room for two nights, please.
\begin{itemize}
\item[a)] demand \item[b)] book \item[c)] do
\end{itemize}
\item Could you give me some \textbf{……} about your services?
\begin{itemize}
\item[a)] informations \item[b)] information \item[c)] an information
\end{itemize}
\item You must pay a registration \textbf{……} of 5,000 francs CFA.
\begin{itemize}
\item[a)] fee \item[b)] price \item[c)] fare
\end{itemize}
\item \textbf{……} you mind opening the window, please?
\begin{itemize}
\item[a)] Will \item[b)] Would \item[c)] Do
\end{itemize}
\item The bank \textbf{……} a new service to its customers last month.
\begin{itemize}
\item[a)] introduced \item[b)] presented \item[c)] assisted
\end{itemize}
\end{enumerate}
\end{exercice}

\begin{exercice}[Vrai ou faux ? Justifie]
Dis si chaque affirmation est vraie ou fausse, puis justifie ta réponse en une phrase.
\begin{enumerate}
\item \eng{“Information” is a countable noun in English, so we can say “an information”.}
\item \eng{“Could you help me, please?” is more polite than “Help me!”}
\item \eng{“To assist” in English means exactly “assister à” in French.}
\item \eng{“Would you mind” is followed by a verb in -ing.}
\item \eng{“To demand” is the most polite way to request a service.}
\end{enumerate}
\end{exercice}

\begin{exercice}[Vocabulaire : complète les phrases]
Complète les phrases avec le mot correct de la liste : \eng{request / demand / information / receipt / fee}.
\begin{enumerate}
\item Customers made a formal \textbf{……} for a refund after the poor service.
\item Please keep your \textbf{……} as proof of payment.
\item The workers \textbf{……} better working conditions from their employer.
\item For further \textbf{……}, please contact our office in Yaoundé.
\item The school charges a small \textbf{……} for the photocopying service.
\end{enumerate}
\end{exercice}

\begin{exercice}[Familles de mots]
Trouve le mot demandé de la même famille.
\begin{enumerate}
\item \eng{to serve} $\rightarrow$ (le nom de l'activité) : \textbf{……}
\item \eng{to apply} $\rightarrow$ (le nom de la personne) : \textbf{……}
\item \eng{help} $\rightarrow$ (l'adjectif positif) : \textbf{……} ; (l'adjectif négatif) : \textbf{……}
\item \eng{polite} $\rightarrow$ (le nom) : \textbf{……}
\item \eng{to employ} $\rightarrow$ (le nom de la personne qui donne le travail) : \textbf{……}
\end{enumerate}
\end{exercice}

\courssub{Je m'entraîne}

\begin{exercice}[Appariement : les collocations]
Relie chaque verbe de la colonne de gauche à la collocation correcte de la colonne de droite (les collocations sont mélangées).

\begin{center}
\begin{tabular}{|>{\eng}l|>{\eng}l|}
\hline
\rowcolor{popblueL} \textbf{Verbe} & \textbf{Collocation} \\
\hline
make & a room \\
\hline
book & a favour \\
\hline
fill in & a fee \\
\hline
pay & a form \\
\hline
do & a request \\
\hline
\end{tabular}
\end{center}

Puis écris une phrase complète en anglais avec chacune des cinq collocations correctes.
\end{exercice}

\begin{exercice}[Traduction]
Traduis les phrases suivantes en anglais.
\begin{enumerate}
\item Pourriez-vous m'aider à remplir ce formulaire, s'il vous plaît ?
\item Je voudrais réserver une table pour quatre personnes.
\item Combien coûte le service de livraison ?
\item Cela vous dérangerait-il de répéter, s'il vous plaît ?
\item Le client a demandé un reçu après avoir payé.
\end{enumerate}
\end{exercice}

\begin{exercice}[Corrige les erreurs de francophones]
Chaque phrase contient une erreur typique d'un élève francophone. Corrige-la et explique brièvement l'erreur.
\begin{enumerate}
\item \eng{I demand a service at the reception desk.}
\item \eng{Can you give me an information about the bus timetable?}
\item \eng{I assist to the meeting every Monday.}
\item \eng{Would you mind to help me with my luggage?}
\item \eng{I want that you repair my phone today.}
\end{enumerate}
\end{exercice}

\begin{exercice}[Requêtes polies]
Transforme ces demandes directes en requêtes polies, en utilisant \eng{Could you…}, \eng{Would you mind…} ou \eng{I would like…}.

Exemple : \eng{Give me a pen.} $\rightarrow$ \eng{Could you lend me a pen, please?}

\begin{enumerate}
\item \eng{Open the door.}
\item \eng{Tell me the price of this shirt.}
\item \eng{Bring me the menu.}
\item \eng{Show me your identity card.}
\item \eng{Wait for me at the entrance.}
\end{enumerate}
\end{exercice}

\courssub{Je me prépare au Bac}

\begin{exercice}[Compréhension et langue — type Bac]
Lis le texte suivant, puis réponds aux questions.

\begin{textebox}[At the microfinance office]
Last Tuesday, Mrs Etoundi went to a microfinance office in Bafoussam to request a small loan for her restaurant business. When she arrived, she politely asked the receptionist for information about the application procedure. The young woman smiled and handed her a form. “Could you fill in this form, please? And do not forget to attach a photocopy of your identity card,” she said.

Mrs Etoundi completed the form carefully and asked how much the registration fee was. “It is only two thousand francs,” the receptionist replied. “Would you mind waiting for a few minutes? The manager will attend to you shortly.” Mrs Etoundi did not mind at all. She sat down and observed the employees serving other customers with patience and courtesy.

Twenty minutes later, the manager received her warmly, examined her documents and promised to study her request within a week. Before leaving, Mrs Etoundi paid the fee and asked for a receipt. “Thank you very much for your assistance,” she said. “You are welcome, Madam. We are always at your service,” the manager answered with a smile.
\end{textebox}

\textbf{A. Comprehension questions} — Réponds en phrases complètes, en anglais.
\begin{enumerate}
\item Why did Mrs Etoundi go to the microfinance office?
\item What two things did the receptionist ask her to do?
\item How much was the registration fee?
\item How long did the manager promise to take to study her request?
\item Find in the text two polite expressions used to request a service.
\end{enumerate}

\textbf{B. Language work}
\begin{enumerate}
\item Find in the text the noun formed from the verb \eng{to assist}.
\item Give the opposite of \eng{politely}.
\item Rewrite in a more polite way : \eng{“Give me the form.”}
\item Complete : \eng{Would you mind …… (wait) a few minutes?}
\end{enumerate}
\end{exercice}

\begin{exercice}[Traduction — type Bac]
Traduis le passage suivant en anglais.

\begin{textebox}[Texte à traduire]
Hier matin, mon père est allé à la banque pour demander un service. Il a rempli un formulaire et a payé les frais. L'employé lui a donné un reçu et lui a promis que sa demande serait traitée rapidement. Mon père l'a remercié pour son aide et il est rentré à la maison très satisfait.
\end{textebox}
\end{exercice}

\begin{exercice}[Expression écrite — type Bac]
\textbf{Sujet :} You spent a night in a hotel in Douala and you were not satisfied with the service. Write a letter of complaint to the hotel manager. In your letter :
\begin{itemize}
\item say when you stayed at the hotel and what you had requested ;
\item describe the problems you faced (room, staff, fees) ;
\item request a refund or another solution politely.
\end{itemize}

Rédige ta lettre en \textbf{120 à 150 mots}, en anglais, avec la mise en forme d'une lettre formelle (adresse, date, formule d'appel, formule de politesse). Utilise au moins \textbf{six mots ou expressions de la leçon} (\eng{request, service, fee, receipt, assistance, I would like, Would you mind…}, etc.).
\end{exercice}



\newpage

\coursec{Corrigés des exercices}

\begin{corrige}[Exercice 1 -- QCM : le bon mot]
\begin{enumerate}
\item \textbf{b) book}. \eng{To book} = réserver (une chambre, un billet). \eng{I would like to book a room for two nights, please.} « Je voudrais réserver une chambre pour deux nuits, s'il vous plaît. » \eng{To demand} = exiger (trop fort) ; \eng{to do} ne convient pas.
\item \textbf{b) information}. \eng{Information} est un nom \textbf{incomptable} en anglais : jamais de \eng{s}, jamais d'article \eng{a/an}. On dit \eng{some information} ou \eng{a piece of information}.
\item \textbf{a) fee}. \eng{Fee} = frais (d'inscription, de scolarité, honoraires). \eng{Price} = prix d'un bien ; \eng{fare} = tarif de transport (bus, taxi).
\item \textbf{b) Would}. Structure de politesse : \eng{Would you mind + V-ing} = « Cela vous dérangerait-il de… ? ».
\item \textbf{a) introduced}. \eng{To introduce a service} = lancer un nouveau service. \eng{To present} = présenter, exposer ; \eng{to assist} = aider (faux ami de « assister à »).
\end{enumerate}
\end{corrige}

\begin{corrige}[Exercice 2 -- Vrai ou faux ?]
\begin{enumerate}
\item \textbf{Faux}. \eng{Information} est \textbf{incomptable} (uncountable) : on dit \eng{some information} ou \eng{a piece of information}, jamais \eng{an information}.
\item \textbf{Vrai}. \eng{Could you help me, please?} est une requête courtoise ; \eng{Help me!} est un impératif perçu comme un ordre brutal.
\item \textbf{Faux}. C'est un \textbf{faux ami} majeur : \eng{to assist} = \textbf{aider} ; « assister à » (être présent) = \eng{to attend}. Ex. : \eng{I attended the meeting.} « J'ai assisté à la réunion. »
\item \textbf{Vrai}. La structure est \eng{Would you mind + gérondif (V-ing)}. Ex. : \eng{Would you mind opening the door?}
\item \textbf{Faux}. \eng{To demand} = « exiger avec autorité » (ex. \eng{The union demands higher wages.} « Le syndicat exige des salaires plus élevés. »). Pour une requête polie : \eng{to request}, \eng{to ask for}, ou les modaux \eng{could} / \eng{would}.
\end{enumerate}
\end{corrige}

\begin{corrige}[Exercice 3 -- Vocabulaire : complète les phrases]
\begin{enumerate}
\item \textbf{request} : \eng{Customers made a formal request for a refund.} « Les clients ont fait une demande formelle de remboursement. » \eng{Demand} serait trop agressif pour une réclamation standard.
\item \textbf{receipt} : \eng{Please keep your receipt as proof of payment.} « Veuillez conserver votre reçu comme preuve de paiement. »
\item \textbf{demand} : \eng{The workers demand better working conditions.} « Les travailleurs exigent de meilleures conditions de travail. » Ici, la revendication justifie le verbe fort \eng{demand}.
\item \textbf{information} : \eng{For further information, please contact our office in Yaoundé.} « Pour de plus amples informations, veuillez contacter notre bureau de Yaoundé. » Incomptable : pas de \eng{s}.
\item \textbf{fee} : \eng{The school charges a small fee for the photocopying service.} « L'école facture de petits frais pour le service de photocopie. »
\end{enumerate}
\end{corrige}

\begin{corrige}[Exercice 4 -- Familles de mots]
\begin{enumerate}
\item \textbf{service} (le service rendu, l'activité). Autre dérivé : \eng{servant} (domestique), \eng{server} (serveur).
\item \textbf{applicant} (le candidat, le demandeur). \eng{Application} = la candidature, la demande écrite ; \eng{application form} = formulaire de demande.
\item \textbf{helpful} (serviable, utile) ; \textbf{helpless} (démuni, sans aide). Attention : \eng{helpful} ne prend qu'un seul \eng{l} à l'écrit.
\item \textbf{politeness} (la politesse). Adverbe : \eng{politely} ; contraire : \eng{impolite} / \eng{rude}.
\item \textbf{employer} (l'employeur, celui qui donne le travail). Ne pas confondre avec \eng{employee} (l'employé, celui qui travaille). Prononciation : \eng{employer} « èm-plo-yeur », \eng{employee} « èm-plo-yii ».
\end{enumerate}
\end{corrige}

\begin{corrige}[Exercice 5 -- Appariement : les collocations]
\textbf{Correspondances correctes :}
\begin{itemize}
\item \eng{make a request} (faire une demande)
\item \eng{book a room} (réserver une chambre)
\item \eng{fill in a form} (remplir un formulaire)
\item \eng{pay a fee} (payer des frais)
\item \eng{do a favour} (rendre un service)
\end{itemize}

\textbf{Phrases complètes (exemples de production) :}
\begin{enumerate}
\item \eng{The customer \textbf{made a request} for a refund after the poor service.} « Le client a fait une demande de remboursement après le mauvais service. »
\item \eng{I would like to \textbf{book a room} with a sea view, please.} « Je voudrais réserver une chambre avec vue sur la mer, s'il vous plaît. »
\item \eng{Please \textbf{fill in a form} at the reception desk.} « Veuillez remplir un formulaire à la réception. »
\item \eng{You must \textbf{pay a fee} of 10,000 francs for the visa.} « Vous devez payer des frais de 10 000 francs pour le visa. »
\item \eng{Could you \textbf{do me a favour} and carry this bag?} « Pourriez-vous me rendre un service et porter ce sac ? »
\end{enumerate}
\end{corrige}

\begin{corrige}[Exercice 6 -- Traduction]
\begin{enumerate}
\item \eng{Could you help me fill in this form, please?} (ou \eng{help me to fill in}). « Pourriez-vous m'aider à remplir ce formulaire, s'il vous plaît ? »
\item \eng{I would like to book a table for four people.} « Je voudrais réserver une table pour quatre personnes. »
\item \eng{How much does the delivery service cost?} (ou \eng{What is the delivery fee?}) « Combien coûte le service de livraison ? »
\item \eng{Would you mind repeating, please?} « Cela vous dérangerait-il de répéter, s'il vous plaît ? »
\item \eng{The customer asked for a receipt after paying.} (ou \eng{after he had paid}) « Le client a demandé un reçu après avoir payé. »
\end{enumerate}
\textbf{Points de vigilance :} « s'il vous plaît » = \eng{please} (jamais \eng{if it pleases you}) ; « réserver » = \eng{to book}, pas \eng{to reserve} dans le registre courant britannique ; après \eng{mind}, toujours le gérondif.
\end{corrige}

\begin{corrige}[Exercice 7 -- Corrige les erreurs de francophones]
\begin{enumerate}
\item \textbf{Correction :} \eng{I would like to request a service at the reception desk.} \\
\textbf{Explication :} \eng{To demand} = « exiger » (ton autoritaire). Pour une demande polie : \eng{to request} ou \eng{to ask for}.
\item \textbf{Correction :} \eng{Can you give me some information about the bus timetable?} \\
\textbf{Explication :} \eng{Information} est \textbf{incomptable} : ni article \eng{an}, ni \eng{s}. On dit \eng{some information} ou \eng{a piece of information}.
\item \textbf{Correction :} \eng{I attend the meeting every Monday.} \\
\textbf{Explication :} Faux ami. \eng{To assist} = \textbf{aider} ; « assister à » = \eng{to attend} (sans préposition \eng{to}).
\item \textbf{Correction :} \eng{Would you mind helping me with my luggage?} \\
\textbf{Explication :} Après \eng{mind}, on utilise le \textbf{gérondif (V-ing)}, jamais l'infinitif.
\item \textbf{Correction :} \eng{I want you to repair my phone today.} (plus poli : \eng{I would like you to repair my phone today.}) \\
\textbf{Explication :} Structure anglaise : \eng{want + objet + to + infinitif}. Le français « je veux que tu répares » (subjonctif) ne se traduit jamais par \eng{want that you repair}.
\end{enumerate}
\end{corrige}

\begin{corrige}[Exercice 8 -- Requêtes polies]
\begin{enumerate}
\item \eng{Could you open the door, please?} / \eng{Would you mind opening the door, please?}
\item \eng{Could you tell me the price of this shirt, please?} / \eng{I would like to know the price of this shirt, please.}
\item \eng{Could you bring me the menu, please?} / \eng{Would you mind bringing me the menu, please?}
\item \eng{Could you show me your identity card, please?} / \eng{Would you mind showing me your identity card, please?}
\item \eng{Would you mind waiting for me at the entrance, please?} / \eng{I would like you to wait for me at the entrance, please.}
\end{enumerate}
\textbf{Rappel méthodologique :} \eng{Could you + base verbale} ; \eng{Would you mind + V-ing} ; \eng{I would like + to + infinitif} (ou \eng{I would like you to + infinitif}). Ajouter \eng{please} renforce la courtoisie.
\end{corrige}

\begin{corrige}[Exercice 9 -- Compréhension et langue (type Bac)]
\textbf{A. Comprehension questions}
\begin{enumerate}
\item \eng{Mrs Etoundi went to the microfinance office to request a small loan for her restaurant business.} « Mme Etoundi s'est rendue au bureau de microfinance pour demander un petit prêt pour son restaurant. »
\item \eng{The receptionist asked her to fill in a form and to attach a photocopy of her identity card.} « La réceptionniste lui a demandé de remplir un formulaire et de joindre une photocopie de sa carte d'identité. »
\item \eng{The registration fee was two thousand francs.} « Les frais d'inscription étaient de deux mille francs. »
\item \eng{The manager promised to study her request within a week.} « Le directeur a promis d'étudier sa demande dans un délai d'une semaine. »
\item \eng{“Could you fill in this form, please?”} et \eng{“Would you mind waiting for a few minutes?”} (on accepte aussi \eng{she politely asked the receptionist for information}).
\end{enumerate}

\textbf{B. Language work}
\begin{enumerate}
\item \textbf{assistance} (\eng{Thank you very much for your assistance.}).
\item \textbf{impolitely} ou \textbf{rudely} (grossièrement).
\item \eng{Could you give me the form, please?} / \eng{Would you mind giving me the form, please?} / \eng{I would like to have the form, please.}
\item \textbf{waiting} : \eng{Would you mind waiting a few minutes?} (gérondif obligatoire après \eng{mind}).
\end{enumerate}

\textbf{Barème indicatif (sur 20) :} A. Compréhension : 10 pts (2 pts par réponse correcte en phrase complète ; 1 pt si l'idée est juste mais la langue fautive). B. Langue : 10 pts (2,5 pts par item).\\
\textbf{Points de vigilance :} répondre par des \textbf{phrases complètes} en reprenant les termes de la question ; ne pas recopier tout le paragraphe ; vérifier la concordance des temps (le texte est au prétérit).
\end{corrige}

\begin{corrige}[Exercice 10 -- Traduction (type Bac)]
\textbf{Proposition de traduction :}

\eng{Yesterday morning, my father went to the bank to request a service. He filled in a form and paid the fee. The clerk gave him a receipt and promised that his request would be processed quickly. My father thanked him for his assistance and went back home very satisfied.}

\textbf{Barème indicatif (sur 10) :} 2 pts par phrase correcte (lexique + grammaire).\\
\textbf{Points de vigilance :}
\begin{itemize}
\item « est allé » : prétérit \eng{went} (verbe irrégulier \eng{go -- went -- gone}).
\item « a rempli » : \eng{filled in} (phrasal verb ; \eng{fill} seul est accepté mais moins idiomatique).
\item « serait traitée » : conditionnel passif \eng{would be processed} (discours rapporté après \eng{promised that}).
\item « son aide » : \eng{his assistance} ou \eng{his help} (jamais \eng{his assist}).
\item « est rentré à la maison » : \eng{went back home} (pas de \eng{to} devant \eng{home}).
\end{itemize}
\end{corrige}

\begin{corrige}[Exercice 11 -- Expression écrite (type Bac) : lettre de réclamation]
\textbf{Proposition de rédaction modèle (environ 140 mots) :}

\begin{textebox}[Model letter of complaint]
12, Rue de la Paix\\
Bonanjo, Douala\\
15th June 2024\\[0.3em]
The Manager\\
Palm Beach Hotel\\
Akwa, Douala\\[0.3em]
Dear Sir or Madam,\\[0.3em]
I am writing to complain about the service I received during my stay at your hotel on 10th June. I had booked a double room and requested a quiet room with air conditioning.\\[0.3em]
Unfortunately, I was very disappointed. The air conditioning did not work, the room was not clean, and the receptionist was rude when I asked for assistance. In addition, I was charged an extra fee for breakfast, although it was included in my booking. I kept the receipt as proof.\\[0.3em]
I would like to request a partial refund of the amount I paid. Would you mind examining my complaint and replying within two weeks?\\[0.3em]
I look forward to hearing from you.\\[0.3em]
Yours faithfully,\\[0.3em]
Alain Mbarga
\end{textebox}

\textbf{Lexique de la leçon réemployé :} \eng{service, booked, requested, assistance, fee, receipt, I would like to request, Would you mind…}.

\textbf{Barème indicatif (sur 20) :}
\begin{itemize}
\item Respect de la tâche et du format de la lettre formelle (adresses, date, \eng{Dear Sir or Madam}, \eng{Yours faithfully}) : 4 pts.
\item Contenu : les trois points du sujet traités : 6 pts.
\item Réemploi d'au moins six mots ou expressions de la leçon : 4 pts.
\item Correction grammaticale et orthographique (prétérit, requêtes polies) : 4 pts.
\item Longueur respectée (120--150 mots) et clarté : 2 pts.
\end{itemize}

\textbf{Points de vigilance :}
\begin{itemize}
\item \eng{Dear Sir or Madam} appelle la formule finale \eng{Yours faithfully} (et non \eng{Yours sincerely}, réservé au cas où le nom du destinataire est connu).
\item Rester ferme mais \textbf{courtois} : une lettre de réclamation utilise \eng{I would like}, \eng{I am writing to complain}, jamais \eng{I demand}.
\item Employer le prétérit pour raconter les faits (\eng{I had booked, the air conditioning did not work, I was charged}).
\item Ne pas traduire « remboursement » par \eng{reimbursement} si l'on n'est pas sûr : \eng{refund} est le mot attendu.
\end{itemize}
\end{corrige}

\newpage

\coursec{Fiche bilan}

\begin{popfigure}
\begin{tikzpicture}[mindmap, grow cyclic, every node/.style={concept, text=white, font=\small\bfseries},
  root concept/.append style={concept color=popdark, font=\large\bfseries},
  level 1/.append style={level distance=3.4cm, sibling angle=60},
  level 2/.append style={level distance=2.1cm, font=\footnotesize}]
\node[root concept]{Requesting services}
  child[concept color=popblue]{ node {Asking for a service}
    child { node {book, order, hire} }
    child { node {appointment, delivery} }
  }
  child[concept color=poporange]{ node {Polite formulas}
    child { node {Could you\ldots?} }
    child { node {I would like\ldots} }
  }
  child[concept color=popgreen]{ node {Collocations}
    child { node {make a request} }
    child { node {provide a service} }
  }
  child[concept color=poppurple]{ node {False friends}
    child { node {demander $\neq$ demand} }
    child { node {assister $\neq$ assist} }
  }
  child[concept color=poppink]{ node {Word families}
    child { node {serve, service, servant} }
    child { node {deliver, delivery} }
  }
  child[concept color=popgold]{ node {In context}
    child { node {dialogues, letters} }
  };
\end{tikzpicture}
\legende{Carte mentale de la leçon : le vocabulaire de la demande de services.}
\end{popfigure}

\begin{bilan}
\textbf{1. Lexique clé : demander un service}
\begin{itemize}
\item \eng{to request a service} : demander un service ; \eng{to book / to reserve} : réserver ; \eng{to order} : commander.
\item \eng{to hire} : louer, engager ; \eng{to repair / to fix} : réparer ; \eng{to deliver} : livrer ; \eng{delivery} : livraison.
\item \eng{an appointment} : un rendez-vous ; \eng{a quotation / a quote} : un devis ; \eng{a bill} : une facture ; \eng{a receipt} : un reçu.
\item \eng{a customer} : un client ; \eng{a service provider} : un prestataire ; \eng{a plumber / an electrician / a tailor} : un plombier / un électricien / un tailleur.
\item \eng{free of charge} : gratuit ; \eng{available} : disponible ; \eng{as soon as possible} : dès que possible.
\end{itemize}

\textbf{2. Formules fonctionnelles de la requête}
\begin{center}
\begin{tabularx}{\linewidth}{|X|X|}\hline
\rowcolor{popblueL} \textbf{English} & \textbf{Français} \\ \hline
\eng{Could you help me, please?} & Pourriez-vous m'aider, s'il vous plaît ? \\ \hline
\eng{I would like to book a room.} & Je voudrais réserver une chambre. \\ \hline
\eng{Would it be possible to\ldots?} & Serait-il possible de\ldots ? \\ \hline
\eng{Could you send me a quotation?} & Pourriez-vous m'envoyer un devis ? \\ \hline
\eng{How much do you charge?} & Combien demandez-vous ? \\ \hline
\eng{Thank you for your help.} & Merci pour votre aide. \\ \hline
\end{tabularx}
\end{center}

\textbf{3. Collocations à connaître par cœur}
\begin{itemize}
\item \eng{to make a request} (faire une demande), \eng{to provide a service} (fournir un service), \eng{to render a service} (rendre un service).
\item \eng{to make an appointment} (prendre rendez-vous), \eng{to pay a bill} (payer une facture), \eng{to place an order} (passer une commande).
\end{itemize}

\textbf{4. Faux amis et pièges}
\begin{itemize}
\item \eng{demander} se dit \eng{to ask for} ; \eng{to demand} = exiger.
\item \eng{assister à} se dit \eng{to attend} ; \eng{to assist} = aider.
\item \eng{une demande d'emploi} = \eng{a job application}, pas « a demand ».
\item Pas de « s » à \eng{information} et \eng{advice} (indénombrables).
\end{itemize}

\textbf{5. Familles de mots et prononciation}
\begin{itemize}
\item \eng{serve} (v.) $\rightarrow$ \eng{service} (n.) $\rightarrow$ \eng{servant} (n.) ; \eng{deliver} (v.) $\rightarrow$ \eng{delivery} (n.) ; \eng{employ} (v.) $\rightarrow$ \eng{employee / employer} (n.).
\item Prononciation : \eng{service} « seur-viss », \eng{request} « ri-kouèste », \eng{receipt} « ri-ssite » (le \eng{p} est muet), \eng{charge} « tchardje ».
\end{itemize}

\textbf{6. Méthode express}
\begin{enumerate}
\item Identifier le service demandé et le prestataire.
\item Choisir une formule polie adaptée au degré de formalité.
\item Préciser le détail : date, prix, lieu, délai.
\item Remercier et confirmer.
\end{enumerate}
\end{bilan}

\begin{aretenir}[Checklist avant le Bac]
\begin{itemize}
\item[$\square$] Je connais au moins dix mots du champ lexical de la demande de services.
\item[$\square$] Je sais dire « demander » avec \eng{to ask for} et jamais avec \eng{to demand}.
\item[$\square$] Je maîtrise les formules polies : \eng{Could you\ldots?}, \eng{I would like\ldots}, \eng{Would it be possible to\ldots?}
\item[$\square$] Je sais employer les collocations \eng{make a request}, \eng{provide a service}, \eng{make an appointment}.
\item[$\square$] Je distingue \eng{to hire}, \eng{to rent} et \eng{to book}.
\item[$\square$] Je sais que \eng{information} et \eng{advice} sont indénombrables.
\item[$\square$] Je prononce correctement \eng{receipt} avec le \eng{p} muet.
\item[$\square$] Je connais les familles de mots : \eng{serve/service}, \eng{deliver/delivery}, \eng{employ/employee}.
\item[$\square$] Je sais rédiger un court dialogue ou une lettre de demande de service.
\item[$\square$] Je sais remercier et clore poliment une requête à l'oral comme à l'écrit.
\end{itemize}
\end{aretenir}

\findecours

\end{document}
